% concatenated from main.tex, packages.tex, cbc-macros.tex, introduction.tex, preliminaries.tex, Two-tiered.tex, algebras-variants.tex, bounding-lemma.tex, completeness.tex, translation-lemma.tex, conclusions.tex, appendix-prelims-algebras.tex, appendix-computational-strength.tex, appendix-characterisation-results-for-function-algebras.tex, appendix-completeness.tex
\pdfoutput=1 
%%
%% This is file `sample-manuscript.tex',
%% generated with the docstrip utility.
%%
%% The original source files were:
%%
%% samples.dtx  (with options: `manuscript')
%% 
%% IMPORTANT NOTICE:
%% 
%% For the copyright see the source file.
%% 
%% Any modified versions of this file must be renamed
%% with new filenames distinct from sample-manuscript.tex.
%% 
%% For distribution of the original source see the terms
%% for copying and modification in the file samples.dtx.
%% 
%% This generated file may be distributed as long as the
%% original source files, as listed above, are part of the
%% same distribution. (The sources need not necessarily be
%% in the same archive or directory.)
%%
%% Commands for TeXCount
%TC:macro \cite [option:text,text]
%TC:macro \citep [option:text,text]
%TC:macro \citet [option:text,text]
%TC:envir table 0 1
%TC:envir table* 0 1
%TC:envir tabular [ignore] word
%TC:envir displaymath 0 word
%TC:envir math 0 word
%TC:envir comment 0 0
%%
%%
%% The first command in your LaTeX source must be the \documentclass command.
%%%% Small single column format, used for CIE, CSUR, DTRAP, JACM, JDIQ, JEA, JERIC, JETC, PACMCGIT, TAAS, TACCESS, TACO, TALG, TALLIP (formerly TALIP), TCPS, TDSCI, TEAC, TECS, TELO, THRI, TIIS, TIOT, TISSEC, TIST, TKDD, TMIS, TOCE, TOCHI, TOCL, TOCS, TOCT, TODAES, TODS, TOIS, TOIT, TOMACS, TOMM (formerly TOMCCAP), TOMPECS, TOMS, TOPC, TOPLAS, TOPS, TOS, TOSEM, TOSN, TQC, TRETS, TSAS, TSC, TSLP, TWEB.
% \documentclass[acmsmall]{acmart}

%%%% Large single column format, used for IMWUT, JOCCH, PACMPL, POMACS, TAP, PACMHCI
% \documentclass[acmlarge,screen]{acmart}

%%%% Large double column format, used for TOG
% \documentclass[acmtog, authorversion]{acmart}

%%%% Generic manuscript mode, required for submission
%%%% and peer review
\documentclass[manuscript,screen]{acmart}
%% Fonts used in the template cannot be substituted; margin 
%% adjustments are not allowed.
%%
%% \BibTeX command to typeset BibTeX logo in the docs
\AtBeginDocument{%
  \providecommand\BibTeX{{%
    \normalfont B\kern-0.5em{\scshape i\kern-0.25em b}\kern-0.8em\TeX}}}

%% Rights management information.  This information is sent to you
%% when you complete the rights form.  These commands have SAMPLE
%% values in them; it is your responsibility as an author to replace
%% the commands and values with those provided to you when you
%% complete the rights form.
\setcopyright{acmcopyright}
\copyrightyear{2018}
\acmYear{2018}
\acmDOI{XXXXXXX.XXXXXXX}

%% These commands are for a PROCEEDINGS abstract or paper.
\acmConference[Conference acronym 'XX]{Make sure to enter the correct
  conference title from your rights confirmation emai}{June 03--05,
  2018}{Woodstock, NY}
%
%  Uncomment \acmBooktitle if th title of the proceedings is different
%  from ``Proceedings of ...''!
%
\acmBooktitle{Woodstock '18: ACM Symposium on Neural Gaze Detection,
 June 03--05, 2018, Woodstock, NY} 
\acmPrice{15.00}
\acmISBN{978-1-4503-XXXX-X/18/06}


% \usepackage[foot]{amsaddr}
%  \usepackage{amssymb}


\usepackage{hyperref}

\usepackage[utf8]{inputenc}


\usepackage{tikz}
% \usetikzlibrary{automata}
\usetikzlibrary{positioning,calc,arrows,shapes,  tikzmark, decorations.pathreplacing}


\usepackage{virginialake}
\usepackage{tikz-cd}%for gran tour diagram
\usepackage{cleveref}
 \usepackage[textsize=tiny, disable]{todonotes}
%   \setlength{\marginparwidth}{1.5cm} %can be commented for smaller todonotes
%  \setuptodonotes{inline} % to make todonotes inline

\usepackage{caption}
\usepackage{subfig}

 \usepackage{enumerate}%formatting

\newcommand{\dfn}{:=}
\renewcommand{\emptyset}{\varnothing}
\renewcommand{\setminus}{-}


\newcommand{\red}[1]{{\color{red}#1}}
\newcommand{\blue}[1]{{\color{blue}#1}}
\newcommand{\purple}[1]{{\color{purple}#1}}
\definecolor{mygreen}{rgb}{0, 0.5, 0}
\newcommand{\green}[1]{{\color{mygreen}#1}}
\newcommand{\orange}[1]{{\color{orange}#1}}

\newcommand{\anupam}[1]{\todo{A: #1}}
\newcommand{\gianluca}[1]{\todo{G: #1}}

\newcommand{\toappendix}[1]{\purple{#1}}
% \newcommand{\toappendix}[1]{} %to delete everything inside \toappendix command

%theorems (same as cyclic-bc file)

\newtheorem{thm}{Theorem}
\newtheorem{lem}[thm]{Lemma}
\newtheorem{prop}[thm]{Proposition}
\newtheorem{cor}[thm]{Corollary}
\newtheorem{conj}[thm]{Conjecture}
\newtheorem{fact}[thm]{Fact}
\newtheorem{cond}[thm]{Condition}

\theoremstyle{definition}
\newtheorem{defn}[thm]{Definition}
\newtheorem{exmp}[thm]{Example}
\newtheorem{conv}[thm]{Convention}

\theoremstyle{remark}
\newtheorem{rem}[thm]{Remark}
\newtheorem*{note}{Note}
\newtheorem*{nota}{Notation}


%sets

\newcommand{\Nat}{\mathbb{N}}
\newcommand{\N}{\Nat}
\newcommand{\Bool}{\{0,1 \}}

%relations

\newcommand{\wfpo}{\vartriangleleft}
\newcommand{\wfpoeq}{\trianglelefteq}
\newcommand{\pref}{\prec}
\newcommand{\prefeq}{\preceq}

\newcommand{\permpref}{\subset}
\newcommand{\permprefeq}{\subseteq}

\newcommand{\pp}{\permpref;\permprefeq}

%types
\newcommand{\n}{N}
\newcommand{\bool}{B}
\newcommand{\sq}{\Box}
\newcommand{\sqn}{{\sq\n}}
\newcommand{\sqb}{{\sq\bool}}
\newcommand{\sn}{\sqn}
\newcommand{\lists}[3]{{#2},\overset{#1}{\ldots},{#3} }

%complexity classes

\newcommand{\ptime}{\mathbf{PTIME}}
\newcommand{\pspace}{\mathbf{PSPACE}}
\newcommand{\elementary}{\mathbf{ELEMENTARY}}
\newcommand{\E}{\mathcal{E}}
\newcommand{\Ebin}{\mathcal{E}_{0,1}}

\newcommand{\f}{\mathbf{F}}
\newcommand{\fptime}{\f\ptime}
\newcommand{\fpspace}{\f\pspace}
\newcommand{\felementary}{\f\elementary}
\newcommand{\alogtime}{\mathbf{ALOGTIME}}
\newcommand{\nc}{\mathbf{NC}}
\newcommand{\flinspace}{\mathbf{FLINSPACE}}

%functions

\newcommand{\num}[1]{\underline{#1}}
\renewcommand{\succ}[1]{\mathsf{s}_{#1}}
\newcommand{\pred}{\mathsf{p}}
\newcommand{\floor}[1]{\lfloor {#1} \rfloor}
\newcommand{\hlf}[1]{\floor{#1 /2}}
\newcommand{\cnd}{\mathsf{cond}}
\newcommand{\proj}[2]{\pi^{#1}_{#2}}
\newcommand{\s}[1]{\vert {#1}\vert}
\newcommand{\depth}[1]{\mathsf{depth}{#1}}
\newcommand{\model}[1]{\denot{#1}}
\newcommand{\denot}[1]{f_{#1}}
\newcommand{\ex}{\mathsf{ex}}
\newcommand{\comp}{\mathsf{comp}}
\newcommand{\poly}{\mathrm{poly}}

\newcommand{\sumlen}[1]{|\hspace{-.15em}|#1|\hspace{-.15em}|}

\newcommand{\sexp}{\mathcal{E}}
\newcommand{\cconc}{\mathcal{C}}
\newcommand{\psearch}{\mathcal{B}}
\newcommand{\incrementation}{\mathcal{I}}
\newcommand{\successor}{\mathcal{S}}
\newcommand{\predecessor}{\mathcal{P}}
\newcommand{\completeness}{\mathcal{F}}
\newcommand{\primrec}{\mathcal{R}}
\newcommand{\lengths}{\mathcal{L}}
\newcommand{\unbounded}{\mathcal{U}}
\newcommand{\binarytounary}{\mathcal{N}}
\newcommand{\gfunction}{\mathcal{G}}
\newcommand{\hfunction}{\mathcal{H}}
\newcommand{\saferecursion}{\mathcal{SR}}

\newcommand{\numeral}[1]{\underline{#1}}
\newcommand{\permi}[1]{\vec{#1}}

%schemes

\newcommand{\saferec}{\mathsf{srec}}
\newcommand{\srec}{\saferec}
\newcommand{\saferecwfpo}[1]{\saferec_{#1}}
\newcommand{\srecwfpo}[1]{\saferecwfpo{#1}}

\newcommand{\srecpp}{\srecwfpo{\permpref}}

\newcommand{\safenesrec}{\mathsf{snrec}}
\newcommand{\snrec}{\safenesrec}
\newcommand{\safenesrecwfpo}[1]{\safenesrec_{#1}}
\newcommand{\snrecwfpo}[1]{\safenesrecwfpo{#1}}

\newcommand{\snrecpp}{\snrecwfpo{\permpref}}

\newcommand{\ssrecwfpo}[1]{\mathsf{ssrec}_{#1}}
\newcommand{\ssnrecwfpo}[1]{\mathsf{ssnrec}_{#1}}

\newcommand{\ssrecpp}{\ssrecwfpo{\permpref}}
\newcommand{\ssnrecpp}{\ssnrecwfpo{\permpref}}

%algebras

\newcommand{\bc}{\mathsf{B}}
\newcommand{\bcwfpo}[1]{\bc^{#1}}
\newcommand{\bcnorec}{\bc^{-}}

\newcommand{\nested}{\mathsf{N}}
\newcommand{\nbc}{\nested\bc}
\newcommand{\nbcwfpo}[1]{\nested\bcwfpo{#1}}

\newcommand{\shallow}{\mathsf{S}}
\newcommand{\sbc}{\shallow\bc}
\newcommand{\sbcwfpo}[1]{\sbc^{#1}}

\newcommand{\bcpp}{\bcwfpo{\permpref}}
\newcommand{\nbcpp}{\nbcwfpo{\permpref}}
\newcommand{\sbcpp}{\sbcwfpo\permpref}

%proofs

\newcommand{\seqar}{\Rightarrow}
\newcommand{\der}{\mathcal{D}}
\newcommand{\pder}[1]{\der_{#1}}
\newcommand{\rd}{{R}_\der}
\newcommand{\derrd}{\langle \der, \rd \rangle}
\newcommand{\bud}{\mathit{Bud}}
\newcommand{\nodordereq}{\sqsubseteq}
\newcommand{\nodorder}{\sqsubset}
\newcommand{\close}[1]{{C}_{#1}}
\newcommand{\open}[1]{{O}_{#1}}
 \newcommand\omicron{o}
 \newcommand{\hnu}{\nu_0}

\newcommand{\id}{\mathsf{id}}
\newcommand{\wk}{\mathsf{w}}
\newcommand{\exch}{\mathsf{e}}
\newcommand{\cntr}{\mathsf{c}}
\newcommand{\cut}{\mathsf{cut}}
\newcommand{\pcut}{\mathsf{pcut}}
\newcommand{\boxlef}{\sq_l}
\newcommand{\boxrig}{\sq_r}
\newcommand{\sql}{\boxlef}
\newcommand{\sqr}{\boxrig}
\newcommand{\zero}{0}
\newcommand{\com}{\mathsf{dis}}
\newcommand{\rules}{\mathsf{r}}





%systems

\newcommand{\cyclic}{\mathsf{C}}
\newcommand{\cbc}{\cyclic\bc}
\newcommand{\ncbc}{\cyclic\nested\bc}
\newcommand{\scbc}{\shallow\cbc}
\newcommand{\slr}{\mathsf{SLR}}
\newcommand{\PR}{\mathsf{PR}}
\newcommand{\tgodel}{\mathsf{T}}
\newcommand{\ntgodel}[1]{\tgodel_{#1}}
\newcommand{\ct}{\mathsf{CT}}

\begin{document}

%%
%% The "title" command has an optional parameter,
%% allowing the author to define a "short title" to be used in page headers.
\title{Cyclic Implicit Complexity}

%%
%% The "author" command and its associated commands are used to define
%% the authors and their affiliations.
%% Of note is the shared affiliation of the first two authors, and the
%% "authornote" and "authornotemark" commands
%% used to denote shared contribution to the research.
\author{Gianluca Curzi}
\affiliation{%
  \institution{University of Gothenburg}
  % \streetaddress{1 Th{\o}rv{\"a}ld Circle}
  % \city{Sweden}
  \country{Sweden}}
  \affiliation{\institution{University of Birmingham} \country{UK}}
\email{gianluca.curzi@gu.se}
\orcid{57216760760}

\author{Anupam Das}
\affiliation{%
  \institution{University of Birmingham}
  % \streetaddress{1 Th{\o}rv{\"a}ld Circle}
  % \city{}
  \country{UK}}
\email{A.Das@bham.ac.uk}
\orcid{0000-0002-0142-3676}



%%
%% By default, the full list of authors will be used in the page
%% headers. Often, this list is too long, and will overlap
%% other information printed in the page headers. This command allows
%% the author to define a more concise list
%% of authors' names for this purpose.
\renewcommand{\shortauthors}{G.~Curzi and A.~Das}

%%
%% The abstract is a short summary of the work to be presented in the
%% article.
\begin{abstract}
  \emph{Circular} (or \emph{cyclic}) proofs have received increasing attention in recent years, and have been proposed as an alternative setting for studying (co)inductive reasoning. 
    In particular, now several type systems based on circular reasoning have been proposed.
    % 
    However, little is known about the complexity theoretic aspects of circular proofs, which exhibit sophisticated loop structures atypical of more common `recursion schemes'. 
    
    This paper attempts to bridge the gap between circular proofs and \emph{implicit computational complexity} (ICC). 
    Namely we introduce a circular proof system based on Bellantoni and Cook's famous safe-normal function algebra, and we identify proof theoretical constraints, inspired by ICC, to characterise the polynomial-time and elementary computable functions.
    Along the way we introduce new recursion theoretic implicit characterisations of these classes that may be of interest in their own right.
\end{abstract}

%%
%% The code below is generated by the tool at http://dl.acm.org/ccs.cfm.
%% Please copy and paste the code instead of the example below.
%%
\begin{CCSXML}
<ccs2012>
   <concept>
       <concept_id>10003752.10003777.10003787</concept_id>
       <concept_desc>Theory of computation~Complexity theory and logic</concept_desc>
       <concept_significance>500</concept_significance>
       </concept>
   <concept>
       <concept_id>10003752.10003790.10003792</concept_id>
       <concept_desc>Theory of computation~Proof theory</concept_desc>
       <concept_significance>500</concept_significance>
       </concept>
 </ccs2012>
\end{CCSXML}

\ccsdesc[500]{Theory of computation~Complexity theory and logic}
\ccsdesc[500]{Theory of computation~Proof theory}

%%
%% Keywords. The author(s) should pick words that accurately describe
%% the work being presented. Separate the keywords with commas.
\keywords{Cyclic proofs, implicit complexity, function algebras,  safe recursion, higher-order complexity}

%% A "teaser" image appears between the author and affiliation
%% information and the body of the document, and typically spans the
%% page.
% \begin{teaserfigure}
%   \includegraphics[width=\textwidth]{sampleteaser}
%   \caption{Seattle Mariners at Spring Training, 2010.}
%   \Description{Enjoying the baseball game from the third-base
%   seats. Ichiro Suzuki preparing to bat.}
%   \label{fig:teaser}
% \end{teaserfigure}

\received{20 February 2007}
\received[revised]{12 March 2009}
\received[accepted]{5 June 2009}

%%
%% This command processes the author and affiliation and title
%% information and builds the first part of the formatted document.
\maketitle


\section{Introduction}
Formal proofs are traditionally seen as finite objects modelling  logical or mathematical reasoning. 
\emph{Non-wellfounded} proofs are a generalisation of this notion to an infinitary (but finitely branching) setting, in which consistency is maintained by a standard global condition: the `progressing' criterion. 
% 
% In non-wellfounded proof theory,  
Special attention is devoted to \emph{regular} (or \emph{circular} or \emph{cyclic}) proofs, i.e.~those non-wellfounded proofs having only finitely many distinct sub-proofs, and which may thus be represented by finite (possibly cyclic) directed graphs. 
For such proofs the progressing criterion may be effectively decided by reduction to the universality problem for B\"uchi automata.

% Regular proofs can be turned into finite structures  called  \emph{cycle normal forms},  usually given as finite trees  with additional ‘backpointers’. \anupam{unnecessary sentence} However, regular proofs   admit  fallacious reasoning.   To recover consistency, a standard solution is based on the  introduction of non-local correctness criteria, e.g.~\emph{progressiveness}, typically checked by B\"{u}chi automata on infinite words. 

% Regular proofs, and their corresponding cycle normal forms, 

Non-wellfounded proofs have been employed to  reason about the modal $\mu$-calculus and fixed-point logics~\cite{niwinski1996games,dax2006proof}, first-order inductive definitions~\cite{brotherston2011sequent},  Kleene algebra~\cite{das2017cut,DP18},  linear logic~\cite{baelde2016infinitary,DeSaurin-infinets}, arithmetic~\cite{Simpson17-cyclic-arithmetic,BerardiTatsuta,das2018logical},  system $\tgodel$~\cite{Das2021-preprint,Kuperberg-Pous21, Das2021}, and continuous cut-elimination~\cite{mints1978finite, fortier2013cuts,bouncing-threads}. 
In particular,  \cite{Kuperberg-Pous21} and \cite{Das2021,Das2021-preprint} investigate the computational expressivity of circular proofs, with respect to the proofs-as-programs paradigm, in the setting of higher-order primitive recursion.

% Due to their coinductive nature, non-wellfounded proofs are able to  define any number-theoretic (partial) function. Definability can be then restricted by combining  regularity and progressiveness:  the former rules out uncomputable functions, while the latter   recovers termination.
% \anupam{compress all of above and get to point}

However little is known about the complexity-theoretic aspects of circular proofs.
Usual termination arguments for circularly typed programs are nonconstructive, proceeding by contradiction and using a non-recursive `totality' oracle (cf.~\cite{Kuperberg-Pous21,Das2021,Das2021-preprint}). 
As a result, these arguments are not appropriate for delivering feasible complexity bounds (cf.~\cite{das2018logical}).
% \toappendix{Kuperberg, Pinault and Pous~\cite{kuperberg2019cyclic} introduced a  circular proof system based on Kleene algebras and seen as a computational machinery for recognising languages. The system and its affine version capture, respectively, the regular languages and the languages accepted in logarithmic time by random-access Turing machines. Their work does not consider the cut rule, which  allows for a proof-theoretical description of function composition and is the key ingredient to express recursion in a cyclic-theoretic setting. }
%

The present paper aims to bridge this gap by proposing a circular foundation for Implicit Computational Complexity (ICC), a branch of computational complexity studying machine-free characterisations of complexity classes. 
Our starting point is  Bellantoni and Cook's famous function algebra $\bc$ characterising the polynomial time computable  functions ($\fptime$) using \emph{safe recursion} \cite{BellantoniCook}. 
The prevailing idea behind safe recursion (and its predecessor, `ramified' recursion \cite{Leivant91}) is to organise data into strata in a way that prevents recursive calls being substituted into recursive parameters of previously defined functions.
This approach has been successfully employed to give resource-bound-free characterisations of polynomial-time \cite{BellantoniCook}, levels of the polynomial-time hierarchy \cite{Bellantoni-fph}, and levels of the Grzegorczyk hierarchy \cite{Wirz99characterizingthe}, and has been extended to higher-order settings too \cite{Hofmann97,Leivant99}.

% \toappendix{Functions of $\bc$ have shape $f(x_1,...,x_n;y_{1},...,y_{m})$, where the semicolon separates the \emph{normal} arguments $x_1,...,x_n$ from the \emph{safe} arguments $y_{1},...,y_{m}$. The idea behind safe recursion is that only normal arguments can be used as recursive  parameters, while recursive calls can only appear in safe position. This prevents recursive calls becoming    recursive parameters of other previously defined functions. In the spirit of the    Curry-Howard paradigm, which identifies proofs and programs,  $\bc$  can be alternatively designed as a sequent calculus proof-system where safe recursion is introduced by a specific inference rule $\saferec$. In this system, a two-sorted function $f(x_1,...,x_n;y_{1},...,y_{m})$ is represented by a derivation of the sequent   $\lists{n}{\sn}{\sn}, \lists{m}{\n}{\n}\seqar \n$ where $\n$ is the ground type for natural numbers and $\sn$  is its modal version.
% }

\subsection*{Circular systems for implicit complexity}
Construing $\bc$ as a type system, we consider non-wellfounded proofs, or \emph{coderivations},   generated by its recursion-free subsystem $\bcnorec$. 
The circular  proof system  $\ncbc$ is then obtained by considering the regular and progressing coderivations of $\bcnorec$ which satisfy a further criterion, \emph{safety}, motivated by the eponymous notion from Bellantoni and Cook's work (cf.~also `ramification' in Leivant's work \cite{Leivant91}). 
On the one hand, regularity and progressiveness ensure that coderivations of $\ncbc$ define total computable functions; on the other hand, the latter criterion ensures that the corresponding equational programs are `safe': the recursive call of a function is never substituted into the recursive parameter of a step function.
    
Despite $\ncbc$ having only ground types, it is able to define equational programs that nest   recursive calls, a property that  typically  arises only in higher-order recursion (cf., e.g., \cite{Hofmann97,Leivant99}).  
In fact, we show that this system defines precisely the elementary computable functions ($\felementary$).
Let us point out that the capacity of circular proofs to simulate some higher-order behaviour reflects an emerging pattern in the literature. 
For instance in~\cite{Das2021,Das2021-preprint} it is  shown that the number-theoretic functions definable by type level $n$  proofs of a circular version of system  $\tgodel$ are exactly those definable by  type level $n+1$ proofs of $\tgodel$. 

In the setting of ICC, Hofmann~\cite{Hofmann97} and Leivant~\cite{Leivant99} already observed that higher-order safe recursion mechanisms can be used to characterise $\felementary$. 
In particular, in~\cite{Hofmann97}, Hofmann presents the type system $\slr$ (Safe Linear Recursion) as a higher-order version of $\bc$  imposing a `linearity' restriction on the higher-order safe recursion operator. 
He shows that this system defines just the polynomial-time computable functions on natural numbers ($\fptime$). 
% , which prevents  duplication of the  recursive calls, and hence their nesting. 

Inspired by~\cite{Hofmann97}, we too introduce a linearity requirement for $\ncbc$ that is able to control the interplay between cycles and the cut rule, called the \emph{left-leaning criterion}. 
The resulting circular  proof system is called $\cbc$, which we show defines precisely $\fptime$. 
% 
% Intuitively, $\cbc$ can be seen as a circular version of $\bc$, while $\ncbc$ can be seen as the circular version of a new function algebra, called  $\nbc$, which generalises $\bc$  by permitting nested versions of safe recursion.
% Let us point out that $\nbc$ is comparable to the type 2 fragment of Leivant's extension of ramified recursion to finite types~\cite{Leivant99}, 

% \red{More specifically, $\nbc$ shares some features with the second-order fragment of  Leivant's extension of ramified recursion to finite types~\cite{Leivant99}, and indeed they both characterise $\felementary$.  However, in order to define nesting of recursive calls  without working with type 2 functionals, $\nbc$  crucially requires the introduction of ‘auxiliary functions', i.e.~oracles.}
% \gianluca{Here stress that the introduction of oracles is crucial to define type 2 recursion (and hence nesting) without actually working with functionals as in Leiveant.}

\subsection*{Function algebras for safe nested recursion and safe recursion along well-founded relations}
As well as introducing the circular systems $\cbc$ and $\ncbc$ just mentioned, we also develop novel function algebras for $\fptime$ and $\felementary$ that allow us to prove the aforementioned complexity characterisations via a `sandwich' technique, cf.~\Cref{fig:picture-main-results}.
This constitutes a novel (and more direct) approach to reducing circularity to recursion, 
that crucially takes advantage of safety.

We give a relativised formulation of $\bc$, i.e.\ with \emph{oracles}, that allows us to define a form of safe \emph{nested} recursion.
The resulting function algebra is called $\nbc$ and is comparable to the type 2 fragment of Leivant's extension of ramified recursion to finite types~\cite{Leivant99}.
The algebras $\bc$ and $\nbc$ will serve as lower bounds for $\cbc$ and $\ncbc$ respectively.

The relativised formulation of function algebras also admits a robust notion of (safe) recursion along a well-founded relation.
We identify a particular well-founded preorder $\permpref$ (`permutation of prefixes') whose corresponding safe recursor induces algebras $\bcpp$ and $\nbcpp$ that will serve as upper bounds for $\cbc$ and $\ncbc$ respectively.  



% \anupam{Our new function algebras should also be highlighted as new contributions. (1) the significance of NB and nested recursion; (2) (and related), algebras based on arbitrary wf relations thanks to our relativised approach.}
% The main results of our paper are the following:
% \begin{itemize}
%     \item a function is definable in $\ncbc$ iff it is $\nbc$ iff it is in $\felementary$;
%     \item a function is definable in $\cbc$ iff it is in $\fptime$.
% \end{itemize}
% First, we define the function algebras $\nbcpp$ and $\bcpp$, respectively obtained by extending  $\nbc$ and $\bc$ with a safe recursion mechanism based on permutation of prefixes. Then we show that $\nbc$ and $\nbcpp$ capture precisely the  elementary time computable functions, while $\bcpp$ captures  the polytime ones. Completeness for $\ncbc$ can be  achieved by showing how to represent the functions of $\nbc$. In a similar way, we represent the functions of $\bc$ in  $\cbc$. Soundness is subtler, as it relies on a translation of coderivations of $\ncbc$ and $\cbc$, which are essentially coinductive objects, into the inductively defined functions of $\nbcpp$ and $\bcpp$, respectively.




\subsection*{Outline}
This paper is structured as follows. Section~\ref{sec:preliminaries} presents $\bc$ as a proof system.
In Section~\ref{sec:two-tiered-circular-systems-on-notation} we define non-wellfounded proofs and their semantics
% We then analyse some  proof-theoretical conditions able to restrict the set of definable functions of $\bcnorec$. 
% This leads us to 
and present the circular proof systems $\ncbc$ and $\cbc$.
In Section~\ref{sec:some-variants} we present the function algebras  $\nbc$, $\bcpp$ and $\nbcpp$. 
Section~\ref{sec:characterization-results-for-function-algebras} shows that $\bcpp$ captures $\fptime$ (\Cref{cor:bcpp-bc-fptime}) and that both $\nbc$ and $\nbcpp$ capture $\felementary$ (Corollary~\ref{cor:nb-elementary-characterization}). 
These results require a delicate Bounding Lemma (Lemma~\ref{lem:boundinglemma}) and an encoding of the elementary functions into $\nbc$ (Theorem~\ref{thm:elementary-in-nbc}). In Section~\ref{sec:completeness} we show that any function definable in $\bc$ is also definable in $\cbc$ (Theorem~\ref{thm:bc-in-cbc}), and that any function definable in $\nbc$ is also definable in $\ncbc$ (Theorem~\ref{thm:nbc-in-ncbc}). 
Finally, in Section~\ref{sec:translation}, we present a translation of $\ncbc$ into $\nbcpp$ that maps $\cbc$ coderivations into $\bcpp$ functions (Lemma~\ref{lem:short-translation-lemma}), by reducing circularity to a form of simultaneous recursion on $\permpref$. 

The main results of this paper are summarised in Figure~\ref{fig:picture-main-results}. 
% Full proofs of our results, as well as further discussion and examples,  can be found in an extended version of this paper available \cite{CurziDas}.

\begin{figure}[t]
    \centering
    \small
\begin{tikzcd}[column sep=normal]
&&&  \felementary  \arrow[rrrd,bend left=10, leftarrow, pos=0.3, "\text{Theorem}~\ref{thm:fp-soundness}"]&&& \\ 
\nbc \arrow[rrr, "\text{Theorem}~\ref{thm:nbc-in-ncbc}"]\arrow[rrru,bend left=10, leftarrow, pos=0.7, "\text{Theorem}~\ref{thm:elementary-in-nbc}"] && &\ncbc 
  \arrow[rrr,  "\text{Lemma}~\ref{lem:short-translation-lemma}"] &  && \nbcpp   \\
  \bc \arrow[rrr, "\text{Theorem}~\ref{thm:bc-in-cbc}"]\arrow[rrrd,bend right=10, leftrightarrow, dashed, pos=0.7, swap,  "\text{\cite{BellantoniCook}}"] &&& \cbc 
      \arrow[rrr, "\text{Lemma}~\ref{lem:short-translation-lemma}"]& && \bcpp  \\ 
      &&& \fptime \arrow[rrru,bend right=10, leftarrow, pos=0.3, swap, "\text{Theorem}~\ref{thm:fp-soundness}"]&& &
\end{tikzcd}
    \caption{Summary of the main results of the paper, where $\rightarrow$ indicates an inclusion ($\subseteq$) of function classes. 
    }
    \label{fig:picture-main-results}
\end{figure}



\subsection*{Comparison with conference version}
This article is a full version of the conference paper \cite{CurziDas:lics22}, which appeared in the proceedings of \emph{Logic in Computer Science 2022}. 
The current version expands upon \cite{CurziDas:lics22} by including full proofs, as well as many further examples and remarks to aid the narrative.\section{Preliminaries}\label{sec:preliminaries}

Bellantoni and Cook introduced in \cite{BellantoniCook} an algebra of functions based on a simple two-sorted structure. This idea was itself inspired by Leivant's characterisations, one of the founding works in Implicit Computational Complexity (ICC) \cite{Leivant91}.
The resulting `tiering' of the underlying sorts has been a recurring theme in the ICC literature since, and so it is this structure that shall form the basis of the systems we consider in this work.

We consider functions on the natural numbers with inputs distinguished into two sorts: `safe' and `normal'.
We shall write functions explicitly indicating inputs, namely writing $f(x_1, \dots, x_{m}; y_1, \dots, y_{n})$ when $f$ takes $m$ normal inputs $\vec x$ and $n$ safe inputs $\vec y$.
Both sorts vary over the natural numbers, but their roles will be distinguished by the closure operations of the algebras and rules of the systems we consider. We will sometimes call these functions \emph{safe-normal}.

Throughout this work, we write $|x|$ for the length of $x$ (in binary notation), and if $\vec x = x_1, \dots, x_m$ we write $|\vec x|$ for the list $|x_1|,\dots, |x_m|$.


\subsection{Bellantoni-Cook characterisation of $\fptime$}
We first recall Bellantoni-Cook in its original guise.


\begin{defn}
[Bellantoni-Cook algebra]\label{defn:bc-function-algebra}
$\bc$ is defined as the smallest class of (two-sorted) functions containing,
\begin{itemize}
    \item $0(;) := 0 \in \Nat$.
    \item $\proj {m;n} {j;} (x_1, \dots, x_{m}; y_1, \dots, y_{n}) := x_j$, for $1 \leq j\leq m$.
    \item $\proj {m;n} {;j} (x_1, \dots, x_{m}; y_1, \dots, y_{n}) := y_j$, for $1 \leq j\leq n$.
    \item $\succ i (;x):= 2x+i$, for $i \in \{0,1\}$.
    \item $\pred (;x) := \hlf x$.
    \item $\cnd(;w,x,y,z) := \begin{cases}
    x & w = 0 \\
    y & w = 0 \mod 2, w\neq 0 \\
    z & w = 1 \mod 2
    \end{cases}$
\end{itemize}
and closed under the following:
\begin{itemize}
    \item (Safe composition) 
    \begin{itemize}
        \item If $f(\vec x, x; \vec y)\, ,\, g(\vec x;) \in \bc$ then $f(\vec x, g(\vec x;);\vec y) \in \bc$.
        \item If $f(\vec x; \vec y, y) \, , \, g(\vec x; \vec y) \in \bc$ then $f(\vec x; \vec y, g(\vec x; \vec y)) \in \bc$.
    \end{itemize}
     \item (Safe recursion on notation)
    %  \footnote{This was originally called `predicative recursion on notation' by Bellantoni and Cook in \cite{BellantoniCook}, in reference to Parsons' predicative higher-order characterisation of the primitive recursive functions (\`a la Kleene).}
    If $g(\vec x;\vec y) \, , \, h_i (x, \vec x; \vec y, y) \in \bc$ for $i=0,1$ then so is $f(x, \vec x;\vec y) $ given by:
    \[
    \arraycolsep=2pt
    \begin{array}{rcll}
        f(0,\vec x; \vec y) & := & g(\vec x;\vec y) \\
        f(\succ 0 x, \vec x; \vec y) & := & h_0 (x, \vec x; \vec y, f(x, \vec x; \vec y)) & \text{if $x \neq 0$}\\
        f(\succ 1 x, \vec x; \vec y) & := & h_1 (x, \vec x; \vec y, f(x, \vec x; \vec y))
    \end{array}
    \]
\end{itemize}
\end{defn}
Intuitively, in a function $f(\vec{x}; \vec y)\in \bc$ only the  normal arguments $\vec x$  can be used as recursive parameters. The idea  behind safe recursion  is that recursive calls can only appear in safe position, and hence they can never be used as recursive parameters of other previously defined functions. Safe composition preserves the distinction between normal and safe arguments by requiring that, when composing  along a normal parameter,  the pre-composing function has no safe parameter at all. As a result, we can effectively substitute a normal parameter into a safe position but not vice-versa.





Writing $\fptime$ for the class of functions computable in polynomial-time, the main result of Bellantoni and Cook is:
\begin{thm}
[\cite{BellantoniCook}] \label{thm:bellantoni}
$f(\vec x;) \in \bc $ if and only if $ f(\vec x) \in \fptime$.
\end{thm}


\subsection{Proof theoretic presentation of Bellantoni-Cook}
We shall work with a formulation of Bellantoni and Cook's algebra as a type system with modalities to distinguish the two sorts (similarly to \cite{Hofmann97}). 
% 
In order to facilitate the definition of the circular  system that we present later, we here work with sequent-style typing derivations. 

We only consider \emph{types} (or \emph{formulas}) $\n$ (`safe') and $\sq \n$ (`normal') which intuitively vary over the natural numbers.
% Morally, $\n$ corresponds to the `safe' sort and $\sn$ corresponds to the `normal' sort in $\bc$.
We write $A,B,$ etc.\ to vary over types.

A \emph{sequent} is an expression $\Gamma \seqar A$, where $\Gamma$ is a list of types (called the \emph{context} or \emph{antecedent}) and $A$ is a type (called the \emph{succedent}). For a list of types $\Gamma = \lists{k}{\n}{\n}$, we write $\sq\Gamma$ for $\lists{k}{\sn}{\sn}$. 



In what follows, we shall essentially identify $\bc$ with the $S4$-style type system in \Cref{fig:bc-type-system}.
The colouring of type occurrences may be ignored for now, they will become relevant in the next section.
Derivations in this system are simply called \emph{$\bc$-derivations}, and will be denoted $\der, \sexp, \ldots$. We write $\der: \Gamma \seqar A$ if the derivation $\der$ has end-sequent  $ \Gamma \seqar A$. We may   write $\der= \rules (\der_1, \ldots, \der_n)$  
(for $ n\leq 3$)
if $\rules$ is the last inference step of $\der$ whose immediate subderivations are, respectively,  $\der_1, \ldots, \der_n$. 
% If $n=0$, we  write $\rules$ in place of $\rules ()$.  
Unless otherwise indicated, we assume that the $\rules$-instance is as typeset in \Cref{fig:bc-type-system}.
 




\begin{figure*}
\[
\hspace{1.3cm}
   \def\arraystretch{3}
   \begin{array}{c}
\vlinf{\id}{}{\n \seqar \n}{}
\quad 
\vliinf{\cut_\n}{}{\orange \Gamma \seqar B}{\orange \Gamma\seqar \n}{\orange \Gamma, \blue \n \seqar B}
\quad
\vliinf{\cut_\sq}{}{\orange \Gamma \seqar B}{\orange \Gamma\seqar \sn}{\blue{\sn},\orange \Gamma \seqar B}
\quad
\vlinf{\wk_\n}{}{\orange\Gamma, \blue \n \seqar B}{\orange \Gamma \seqar B}
\quad
\vlinf{\wk_{\sq}}{}{\blue{\sn},\orange\Gamma \seqar B}{\orange \Gamma \seqar B}
\quad
\vlinf\exch{}{\orange \Gamma, \blue B,\red A, \purple{\Gamma'} \seqar C}{\orange \Gamma, \red A, \blue B, \purple{\Gamma'} \seqar C}
\\
\vlinf{\sql}{}{\blue \sn, \orange \Gamma \seqar A}{\orange \Gamma, \blue \n \seqar A}
\quad
\vlinf{\sqr}{}{\orange {\sq \Gamma} \seqar \sn}{ \orange {\sq \Gamma} \seqar \n}
\quad
\vlinf{\zero}{}{\seqar \n}{}
\quad
\vlinf{\succ 0 }{}{\orange \Gamma \seqar A}{\orange \Gamma \seqar A}
\quad
\vlinf{\succ 1 }{}{\orange \Gamma \seqar A}{\orange \Gamma \seqar A}
\quad 
\vliiinf{\saferec}{}{\blue\sn, \orange \Gamma \seqar \n}{\orange \Gamma \seqar \n}{\blue\sn, \orange \Gamma  , \red \n \seqar \n}{\blue\sn, \orange \Gamma, \red \n \seqar \n}
\\ 
\vliiinf{\cnd_\n}{}{\orange \Gamma, \blue \n \seqar \n}{\orange \Gamma \seqar \n}{\orange \Gamma, \blue \n \seqar \n}{\orange \Gamma , \blue \n \seqar \n}
\qquad
\vliiinf{\cnd_\sq}{}{\blue{\sn},\orange \Gamma\seqar \n}{\orange \Gamma \seqar \n}{\blue{\sn},\orange \Gamma \seqar \n}{\blue{\sn}, \orange \Gamma  \seqar \n}
\end{array}
\]
    \caption{System $\bc$, as a sequent-style type system.}
    \label{fig:bc-type-system}
\end{figure*}

\begin{conv}
[Left normal, right safe]
In what follows, we  assume that sequents have shape ${\lists{}{\sn}{\sn}, \lists{}{\n}{\n} \seqar A}$, i.e.\ in the LHS all $\sn$ occurrences are placed before all $\n$ occurrences.
Note that this invariant is maintained by the typing rules of \Cref{fig:bc-type-system}, as long as we insist that $A=B$ in the exchange rule $\exch$.
This effectively means that exchange steps have one of the following two forms:
\[
\vlinf{\exch_\n}{}{\orange \Gamma, \blue \n, \red \n, \purple{\vec \n'} \seqar A }{\orange \Gamma, \red \n, \blue \n, \purple{\vec \n'}\seqar A}
\qquad
\vlinf{\exch_\sq}{}{\orange{\sq \vec \n} , \blue \sn , \red \sn, \purple{\Gamma'} \seqar A}{\orange{\sq \vec \n}, \red \sn, \blue \sn, \purple{\Gamma'} \seqar A}
\]
Let us point out that this convention does not change the class of definable functions with only normal inputs, under the semantics we are about to give.

{Henceforth, we shall often refrain from indicating explicit instances of the exchange rule when constructing derivations.}
\end{conv}

We construe the system of $\bc$-derivations as a class of safe-normal functions by identifying each rule instance as an operation on safe-normal functions.
% and identifying the system with the class of thus definable functions. 
Formally:

\begin{defn}
[Semantics]
\label{defn:semantics-bc}
Given a $\bc$-derivation $\der$ with conclusion  $\lists{m}{\sn}{\sn}, \lists{n}{\n}{\n} \seqar A$ 
we define a two-sorted function $\denot \der (x_1, \dots, x_m;y_1, \dots, y_n)$ by induction on the structure of $\der$ as follows (all rules as typeset in \Cref{fig:bc-type-system}):
\begin{itemize}
    % \item $\denot{\id(;y)} \dfn y $. \anupam{I think rather $\denot \id (y) \dfn y$?}
    \item If $\der = \id$ then $\denot{\der} (;y) \dfn y $.
    % \item   $\denot {\wk_\n(\der_0)} (\vec x; \vec y, y) \dfn \denot{\der_0} (\vec x;\vec y)$.
    \item If $\der =\wk_\n(\der_0)$ then  $\denot {\der} (\vec x; \vec y, y) \dfn \denot{\der_0} (\vec x;\vec y)$.
    % \item $\denot {\wk_\sq(\der_0)} (x, \vec x; \vec y) \dfn \denot{\der_0} (\vec x;\vec y)$.
    \item If $\der=\wk_\sq(\der_0)$ then $\denot {\der} (x, \vec x; \vec y) \dfn \denot{\der_0} (\vec x;\vec y)$.
    % \item $\denot {\exch_\n(\der_0)} (\vec x; \vec y, y,y', \vec y') \dfn \denot{\der_0} (\vec x; \vec y, y',y, \vec y')$.
    \item If $\der=\exch_\n(\der_0)$ then $\denot {\der} (\vec x; \vec y, y,y', \vec y') \dfn \denot{\der_0} (\vec x; \vec y, y',y, \vec y')$.
    % \item  $\denot {\exch_\sq(\der_0)} (\vec x, x, x', \vec x'; \vec y) \dfn \denot{\der_0} (\vec x, x',x , \vec x'; \vec y)$.
    \item If $\der=\exch_\sq(\der_0)$ then $\denot {\der} (\vec x, x, x', \vec x'; \vec y) \dfn \denot{\der_0} (\vec x, x',x , \vec x'; \vec y)$.
    % \item  $\denot { \sql(\der_0)} (x, \vec x; \vec y) \dfn \denot{\der_0}(\vec x; \vec y, x)$.
    \item If $\der=\sql(\der_0)$ then $\denot { \der} (x, \vec x; \vec y) \dfn \denot{\der_0}(\vec x; \vec y, x)$.
    % \item $\denot {\sqr(\der_0)} (\vec x;) \dfn \denot{\der_0} (\vec x;)$.
    \item If $\der=\sqr(\der_0)$ then $\denot {\der} (\vec x;) \dfn \denot{\der_0} (\vec x;)$.
    % \smallskip
    % \item  $\denot {0}(;) \dfn 0$.
    \item If $\der=0$ then $\denot {\der}(;) \dfn 0$.
    % \item  $\denot {\succ 0(\der_0)} (\vec x;\vec y) \dfn \succ 0 (; \denot{\der_0}(\vec x;\vec y))$.
    \item If $\der=\succ 0(\der_0)$ then $\denot {\der} (\vec x;\vec y) \dfn \succ 0 (; \denot{\der_0}(\vec x;\vec y))$.
    % \item  $\denot {\succ 1(\der_0)} (\vec x;\vec y) \dfn \succ 1 (; \denot{\der_0}(\vec x;\vec y))$.
    \item  If $\der=\succ 1(\der_0)$ then $\denot {\der} (\vec x;\vec y) \dfn \succ 1 (; \denot{\der_0}(\vec x;\vec y))$.
    % \smallskip
    % \item  $\denot {\cut_\n(\der_0, \der_1)} (\vec x;\vec y) \dfn \denot{\der_1}(\vec x; \vec y, \denot{\der_0} (\vec x; \vec y))$.
    \item If $\der = \cut_\n(\der_0, \der_1)$ then $\denot {\der} (\vec x;\vec y) \dfn \denot{\der_1}(\vec x; \vec y, \denot{\der_0} (\vec x; \vec y))$.
    % \item $\denot {\cut_\sq(\der_0, \der_1)} (\vec x;\vec y) \dfn \denot{\der_1}(\denot{\der_0} (\vec x; \vec y), \vec x; \vec y)$.
    \item If $\der = \cut_\sq(\der_0, \der_1)$ then $\denot {\der} (\vec x;\vec y) \dfn \denot{\der_1}(\denot{\der_0} (\vec x; \vec y), \vec x; \vec y)$.
    \item If $\der=\cnd_\n(\der_0, \der_1, \der_2)$ then:
    \[
        \arraycolsep=2pt
    \begin{array}{rcll}
         \denot {\der} (\vec x;\vec y,0) & \dfn &  \denot{\der_0} (\vec x;\vec y) \\
         \denot {\der} (\vec x; \vec y, \succ 0 y) & \dfn & \denot{\der_1} (\vec x; \vec y, y)  & \text{if $y\neq 0$} \\
         \denot {\der} (\vec x; \vec y, \succ 1 y) & \dfn & \denot{\der_2} (\vec x; \vec y, y) 
    \end{array}
    \]
    \item  If $\der=\cnd_\sq(\der_0, \der_1, \der_2)$  then:
    \[
            \arraycolsep=2pt
    \begin{array}{rcll}
         \denot \der (0, \vec x;\vec y) & \dfn &  \denot{\der_0} (\vec x;\vec y) \\
         \denot \der (\succ 0 x, \vec x; \vec y) & \dfn & \denot{\der_1} (x, \vec x; \vec y) & \text{if $x \neq 0$} \\
         \denot \der (\succ 1 x , \vec x; \vec y) & \dfn & \denot{\der_2} (x, \vec x; \vec y) 
    \end{array}
    \]
    \item If $\der =\srec(\der_0, \der_1, \der_2)$ then:
    \[
     \arraycolsep=2pt
    \begin{array}{rcll}
         \denot \der (0, \vec x;\vec y) & \dfn &  \denot{\der_0} (\vec x;\vec y) \\
         \denot \der (\succ 0 x, \vec x; \vec y) & \dfn & \denot{\der_1} (x, \vec x; \vec y, \denot\der(x, \vec x; \vec y)) & \text{if $x \neq 0$} \\
         \denot \der (\succ 1 x , \vec x; \vec y) & \dfn & \denot{\der_2} (x, \vec x; \vec y,  \denot\der(x, \vec x; \vec y)) 
    \end{array}
    \]
\end{itemize}
\end{defn}




This formal semantics exposes how $\bc$-derivations and $\bc$ functions relate. The rule $\saferec$ in~\Cref{fig:bc-type-system} corresponds to safe recursion, and safe composition along safe parameters is expressed by means of the rules $\cut_\n$. 
Note, however, that the function $f_\der$ is not quite defined according to function algebra $\bc$, due to the  interpretation of the $\cut_\sq$ rule apparently not satisfying the required constraint on safe composition along a normal parameter.
However, this admission turns out to
be harmless, as exposited in the following proposition:




\begin{prop}\label{prop:modal-succedent-derivations}
% \marginpar{\footnotesize Proof in \cref{prf:modal-succedent-derivations}.}
 Given a $\bc$-derivation $\der: \sq \Gamma , \vec \n \seqar \sn$, there is a $\bc$-derivation   $\der^*:\sq\Gamma \seqar \sn$ no larger than $\der$ such that, for any $\vec x$, $\vec y$:
\begin{equation*}
 \denot\der(\vec{x}; \vec{y})= \denot{\der^*}(\vec{x};).
\end{equation*}
\end{prop} 
\begin{proof}
The proof is by induction on the size of the derivation. The case where the last rule of $\der$ is an instance of $\id, \zero, \sqr, \cnd_{\n}, \cnd_\sq$, or $\srec$ holds vacuously. If  the last rule of $\der$ is an instance of $\exch_{\n},\exch_{\sq},  \sql,\succ i$, and $\wk_\sq$ then we apply the induction hypothesis. Let us now suppose that $\der$ has been obtained from a derivation $\der_0$ by applying an instance of  $\wk_\n$. By induction hypothesis, there exists a derivation  $\pder0^*:\lists{n}{\sn}{\sn} \seqar \sn$ such that $\model{\pder0}(\vec x; \vec y)= \model{\pder0^*}(\vec x; )$. Since $\model{\der}(\vec x; \vec y, y)= \model{\pder0}(\vec x; \vec y)=\model{\pder0^*}(\vec x; )$ we just set $\der^*= \pder{0}^*$.  Suppose now that $\der$ is obtained from two derivations $\der_0$ and $\der_1$ by applying an instance of  $\cut_\n$. By induction hypothesis, there exists $\pder{1}^*$ such that $\model{\pder{1}}(\vec x;\vec y,y)=\model{\pder{1}^*}(\vec x;)$. Since 
 $\model{\pder{1}}(\vec x; \vec y,  \model{\pder{0}}(\vec x; \vec y))= \model{\pder{1}^*}(\vec x;)$, we set $\der^*=\pder{1}^*$. As for the case where the last rule is $\cut_\sq$, by induction hypothesis, there exist derivations $\pder{0}^*$ and $\pder{1}^*$ such that $\model{\pder{0}}(\vec x; \vec y)= \model{\pder{0}^*}(\vec x;)$ and $\model{\pder{1}}(\vec x,x; \vec y)= \model{\pder{1}^*}(\vec x, x;)$, so that we define $\der^*$ as the derivation obtained from $\pder{0}^*$ and $\pder{1}^*$ by applying the rule $\cut_{\sq}$.
\end{proof}

% \begin{prop}\label{prop:modal-succedent-derivations}
%  Given a $\bc$-derivation $\der: \sq \Gamma , \vec \n \seqar \sn$, there is a $\bc$-derivation   $\der^*:\sq\Gamma \seqar \sn$   (of smaller size) such that:
% \begin{equation*}
%  \denot\der(\vec{x}; \vec{y})= \denot{\der^*}(\vec{x};).
% \end{equation*}
% \end{prop} 

% \begin{proof}[Proof sketch]  The proof is by induction on $\der$, using the fact that   $\cnd_{\n}$, $\cnd_{\sn}$ and $\saferec_{\sn}$  have only non-modal succedents.
% \end{proof}


Our overloading of the notation $\bc$, for both a function algebra and for a type system, is
now justified by:

\begin{prop}
% [Folklore]
\label{prop:bc-type-system-characterisation}
 $f(\vec x;\vec y)\in \bc$ iff there is a $\bc$-derivation $\der$ for which $\denot \der (\vec x; \vec y) = f(\vec x;\vec y)$.
\end{prop}
\begin{proof}
Let us first prove the left-right implication by induction on $f \in \bc$. The cases where $f$ is $\zero$ or $\succ i$ are straightforward. If $f$ is a projection then we construct $\der$ using the rules $\id$, $\wk_{\n}$, $\wk_\sq$, $\exch_\n$, and $\exch_{\sq}$. If $f=\pred(; x)$ then $\der$ is as follows: 
\[
\small
\vliiinf{\cnd_\n}{}{\n \seqar \n}{\vlinf{\zero}{}{\seqar \n}{\vlhy{\ }}}{\vlinf{\id}{}{\n \seqar \n}{\vlhy{\ }}}{\vlinf{\id}{}{\n \seqar \n}{\vlhy{\ }}}
\]
If $f$ is $\cnd(;x,y,z,w)$ then $\der$ is constructed using $\id$, $\wk_\n$, $\exch_\n$ and $\cnd_\n$.
% \[
% \vlderivation{
% \vliiin{\cnd_\n}{}{\underline{\n}, \n,\n, \n \seqar \n}
% {\vlin{\wk_\n}{}{\n,\n, \underline{\n}\seqar \n}{\vlin{\wk_\n}{}{\n, \underline{\n}\seqar \n}{\vlhy{\ }}}}
% {\vlin{\wk_\n}{}{\underline{\n}, \n, \n, \n \seqar \n}{\vlin{\wk_\n}{}{\underline{\n}, \n, \n \seqar \n}{\vlin{\wk_\n}{}{\n, \underline{\n}  \seqar \n}{\vlin{\id}{}{\n \seqar \n}{\vlhy{\ }}}}}}
% {\vlin{\wk_\n}{}{\underline{\n}, \n, \n, \n \seqar \n}{\vlin{\wk_\n}{}{\underline{\n}, \n, \n \seqar \n}{\vlin{\wk_\n}{}{\underline{\n}, \n  \seqar \n}{\vlin{\id}{}{\n \seqar \n}{\vlhy{\ }}}}}}
% }
% \]
Suppose now that $f(\vec x; \vec y)= h(\vec x, g(\vec x;); \vec y)$. Then $\der$ is as follows:
\[
\small
\vlderivation{
\vliin{\cut_\sq}{}
{\sq \vec \n, \vec \n \seqar \n}
{\vlin{\sqr}{}{\sq \vec \n\seqar \red{\sn}}{\vltr{\der_0}{\sq \vec \n\seqar \n}{\vlhy{\ }}{\vlhy{\ }}{\vlhy{\ }}}}
{\vltr{\der_1}{ \sq \vec \n, \red{\sn}, \vec \n \seqar \n}{\vlhy{\ }}{\vlhy{\ }}{\vlhy{\ }}}
}
\]
where $\der_0$ and $\der_1$ are such that $\denot {\der_0}=g$ and $\denot{\der_1}=h$. If $f(\vec x; \vec y)= h(\vec x; \vec y, g(\vec x; \vec y))$, then $\der$ is as follows:
\[
\small
\vlderivation{
\vliin{\cut_\n}{}
{\sq \vec \n, \vec \n \seqar \n}
{\vltr{\der_0}{\sq \vec \n, \vec \n \seqar \red{\n}}{\vlhy{\ }}{\vlhy{\ }}{\vlhy{\ }}}
{\vltr{\der_1}{ \sq \vec \n, \vec \n, \red{\n} \seqar \n}{\vlhy{\ }}{\vlhy{\ }}{\vlhy{\ }}}
}
\]
where $\der_0$ and $\der_1$ are such that $\denot {\der_0}=g$ and $\denot{\der_1}=h$. Last, suppose that $f(x, \vec x; \vec y)$ has been obtained by safe recursion from $g(\vec x; \vec y)$ and $h_i(x, \vec x; \vec y, y)$ with $i=0,1$. Then,  $\der$ is as follows:
\[
\small
\vlderivation{
\vliiin{\srec}{}
{\sn, \sq \vec \n, \vec \n \seqar \n}
{\vltr{\der_{0}}{\sq \vec \n, \vec \n \seqar \n}{\vlhy{\ }}{\vlhy{\ }}{\vlhy{\ }}}
 {\vltr{\der_{1}}{\sn, \sq \vec \n, \vec \n, \n   \seqar \n}{\vlhy{\ }}{\vlhy{\ }}{\vlhy{\ }}}
 {\vltr{\der_{2}}{\sn, \sq \vec \n, \vec \n, \n   \seqar \n}{\vlhy{\ }}{\vlhy{\ }}{\vlhy{\ }}}
}
\]
where $\pder0$, $\pder1$, and $\pder2$ are such that $\denot{\der_0}=g$, $\denot{\der_1}=h_0$ and $\denot{\der_2}=h_1$. 

For the right-left implication, we prove by  induction on the size of derivations that $\der: \lists{n}{\sn}{\sn},\lists{m}{\n}{\n} \seqar C$ implies $\denot \der\in \bc$. The cases where the last rule of $\der$ is $\id, \zero, \succ i, \wk_{\n}, \wk_\sq, \exch_\n,\exch_\sq, \sql, \sqr$ are all straightforward using constants, successors and projections. If $\der$ has been obtained from two derivations $\pder0$ and $\pder1$ by applying an instance of $\cut_\n$ then, by applying the induction hypothesis,  $\denot \der(\vec x; \vec y)= \denot{\pder1}(\vec x; \vec y, \denot{\pder0}(\vec x; \vec y))\in \bc$. As for the case where the last rule is $\cut_\sq$, by  Proposition~\ref{prop:modal-succedent-derivations}  there exists a derivation $\pder0^*$ with size no larger than $\der_0$ such that  $\denot{\pder0}(\vec x;\vec y)=\denot{\pder0^*}(\vec x;)$. By applying the induction hypothesis, we have $\denot\der(\vec x; \vec y)= \denot{\pder1}(\vec x, \denot{\pder0^*}(\vec x;); \vec y)\in \bc$. If $\der$ has been obtained from derivations $\pder0$, $\pder1$, $\pder2$ by applying an instance of  $\cnd_\n$ then, by using the induction hypothesis,  $\denot\der(\vec x; \vec y, y)= \cnd(; y,\denot{\pder0}(\vec x; \vec y), \denot{\pder1}(\vec x; \vec y,\pred(;y)), \denot{\pder2}(\vec x; \vec y,\pred(;y)) )\in \bc$. If last rule is $\cnd_{\sq}$, by  induction hypothesis, 
$\denot\der(x,\vec x; \vec y)= \cnd(;  x, \denot{\pder0}(\vec x; \vec y), \denot{\pder1}(\pred(x;),\vec x; \vec y), \denot{\pder2}(\pred(x;),\vec x; \vec y))$, where    $\pred(x;)= \pred(;\proj{1;0}{1;}(x; )) $, and so $\denot\der(x,\vec x; \vec y)\in \bc$. Last, if $\der$ has been obtained from derivations $\pder0$, $\pder1$, $\pder2$ by applying an instance of $\srec$, then  $\denot\der\in \bc$  using the induction hypothesis and the safe recursion scheme.
\end{proof}
% \begin{proof}[Proof sketch]
% The left-right implication follows by a routine induction on the definition of $f(\vec x; \vec y)$.
% For the right-left implication we proceed by induction on the structure of $\der$. The only non-trivial case is when the last rule is $\cut_\sq$, for which we appeal to~\Cref{prop:modal-succedent-derivations}
% to recover a correct instance of safe composition along a normal parameter.
% % 
% % For the right-left implication we proceed by induction on the structure of $\der$. 
% % % The difficulty is in handling normal-valued derivations, i.e.\ those with modal succedent.
% % In particular, notice that the interpretation $\denot\der$ of a derivation $\der$ given in \Cref{defn:semantics-bc} almost induces a bona fide interpretation of typing derivations into $\bc$, but for the case where the last rule of $\der$ is $\cut_\sq$. To handle this we appeal to Proposition~\ref{prop:modal-succedent-derivations}, which allows us to infer that, if    $\pder0: \sq\Gamma, \vec \n \seqar \sn$ is the sub-derivation of the leftmost premise of $\cut_\sq$, it can be replaced by the (smaller) subderivation $\pder0^*: \sq \Gamma \seqar \sn$ constructed in \Cref{prop:modal-succedent-derivations}. From here, the interpretation from \Cref{defn:semantics-bc} induces a correct instance of safe composition along a normal parameter.
% \end{proof}

% The above proposition justifies the overuse of the notation $\bc$, which names both a function algebra and a type system.
% \anupam{repetition of earlier paragraph}

\begin{conv}
Given \Cref{prop:bc-type-system-characterisation} above, 
we shall henceforth freely write $f(\vec x ; \vec y) \in \bc$ if there is a derivation $\der:\sq\Gamma, \vec \n \seqar \n$ with $\denot \der (\vec x; \vec y) = f(\vec x;\vec y) $.
\end{conv}


\section{Two-sorted circular systems on notation}\label{sec:two-tiered-circular-systems-on-notation}

In this section we  introduce a `coinductive' version of $\bc$,  and we study global  criteria that tame its computational strength. 
This proof-theoretic investigation will lead us to two relevant circular systems: $\ncbc$, which morally permits  `nested' versions of safe recursion, and $\cbc$, which will turn out to be closer to usual safe recursion.

Throughout this section we shall work with the set of typing rules $\bcnorec \dfn \bc \setminus \{\saferec\}$.

\subsection{Non-wellfounded typing derivations}
To begin with, we define the notion of `coderivation', which is the fundamental formal object of this section.

\begin{defn}
[Coderivations]
A ($\bcnorec$-)\emph{coderivation} $\der$ is a possibly infinite {rooted} tree (of height $\leq \omega$) generated by the rules of $\bcnorec$.
Formally, we identify $\der$ with a prefix-closed subset of $ \{0,1,2\}^*$ (i.e.\ a ternary tree) where each node  is labelled by an inference step from $\bcnorec$ such that, whenever $\nu\in \der$ is labelled by a step $\vliiinf{}{}{S}{S_1}{\cdots}{S_n}$, for $n\leq 3$, $\nu$ has $n$ children in $\der$ labelled by steps with conclusions $S_1, \dots, S_n$ respectively. Sub-coderivations of a coderivation $\der$ rooted at position $\nu \in \{0,1,2 \}^*$ are typically denoted $\pder\nu$, so that $\pder\epsilon=\der$. 
We write  $\nu \sqsubseteq \mu$ (or  ~$\nu \sqsubset \mu$) if $\nu$ is a prefix (respectively, a strict prefix) of $\mu$, and in this case we say that $\mu$ is \emph{above} (respectively, \emph{strictly above}) $\nu$ or that $\nu$ is \emph{below} (respectively, \emph{strictly below}) $\mu$. 
% \red{We extend this order to node labels, i.e., sequents.} 
\anupam{this is unclear to me: check if this is ever used, if not delete. }


We say that a coderivation is \emph{regular} (or \emph{circular}) if it has only finitely many distinct sub-coderivations.
\end{defn}



Note that, while usual derivations may be naturally written as finite trees or dags, regular coderivations may be naturally written as finite directed (possibly cyclic) graphs.
Some examples of regular coderivations can be found in~\Cref{fig:examples-regular-coderivations}, employing the following writing conventions:

\begin{conv}[Representing coderivations]\label{conv:convention}
Henceforth, we may mark steps by $\bullet$ (or similar) in a regular coderivation to indicate roots of identical sub-coderivations. Moreover, to avoid ambiguities and  to ease parsing of (co)derivations, we shall often underline principal formulas of a  rule instance in a given coderivation and omit instances of $\wk_{\sq}$ and $\wk_\n$ as well as certain structural steps, e.g.\ during a cut step. 
% Double inference lines mean application of several rules, and are typically labelled by a list of those rules that have been applied from top.

Finally, when the sub-coderivations $\der_0$ and $\der_1$ above the second and the third premise of the conditional rule (from left) are similar (or identical), we may compress them into a single parametrised sub-coderivation $\der_i$ (for $i=0,1$).
\end{conv}



\begin{figure*}
    \centering
\[
\small
\vlderivation{
\vliin{\cut_{\sq}}{\bullet}{\blue{\sn} \seqar \n}{
    \vlin{ \sqr}{}{\blue{\sn} \seqar \sn}{\vlin{\sql}{}{\blue{\sn} \seqar \n}{\vlin{ \succ 1}{}{\n \seqar \n}{\vlin{\id}{}{\n \seqar \n}{\vlhy{}}}}}
}{
   \vlin{\cut_{\sq}}{\bullet}{\sn \seqar \n}{\vlhy{\vdots}}
}
}
% \qquad 
% \vlderivation{
% \vliiin{\cnd_{\sq}}{\bullet}{\blue{\sn} \seqar \n}
% {
%   \vliin{\cut_{\n}}{}{\seqar \n}{
%           \vlin{\zero}{}{\seqar \n}{\vlhy{\ }}
%           }{
%           \vlin{\id, \succ 1}{}{
%           \n \seqar \n
%           }{\vlhy{\ }
%           }
%           }
% }
% {
%     \vlin{\sql}{}{\blue{\sn} \seqar \n}{\vliq{\id, \succ 1}{}{\n \seqar \n }{\vlhy{\ }}}
% }{
%     \vliin{\cut_{\n}}{}{\blue{\sn} \seqar \n}{
%       \vlin{\cnd_{\sq}}{\bullet}{\blue{\underline{\sn}} \seqar \n}{\vlhy{\vdots}}
%           }{
%           \vliq{\id, \succ 0}{}{\n \seqar \n}{
%             \vlhy{\ }
%           }
%           }
% }
% }
\qquad \qquad 
\vlderivation{
\vliin{\cnd_\sq}{\bullet}{\blue{\underline{\sn}}, \sq \vec \n \seqar \n}
{
\vltr{\gfunction}{\sq \vec \n \seqar \n}{\vlhy{\ }}{\vlhy{\ }}{\vlhy{\ }}
}
{
  \vliin{\cut_\sq}{{\scriptstyle{i=0,1}}}{\blue{\sn}, \sq \vec \n \seqar \n}
  {
  \vlin{\sqr}{}{\blue{\sn}, \sq \vec \n  \seqar \red{\sn}}{
    \vlin{\cnd_{\sq}}{\bullet}{\blue{\sn}, \sq \vec \n  \seqar {\n}}{\vlhy{\vdots}}
  }
  }{
  \vltr{\hfunction_i}{\blue{\sn}, \sq \vec \n, \red{\sn}  \seqar \n}{\vlhy{\ }}{\vlhy{\ }}{\vlhy{\ }}
  }
}
{
}
}
\]
\[
\small
\vlderivation{
\vliin{\cnd_\sq}{\bullet}{\blue{\underline{\sn}}, \red{\sn} , \n \seqar \n }
 {
   \vliin{\cnd_\sq}{\circ}{\red{\underline{\sn}}, \n \seqar \n }
   {
    \vlin{\id}{}{\n \seqar \n}{\vlhy{}}
   }
   {
   \vlin{\succ i}{{\scriptstyle{i=0,1}}}{\red{\sn},  \n \seqar \n}{\vlin{\cnd_\sq}{\circ}{ \red{\sn}, \n \seqar \n}{\vlhy{\vdots}}}
   }
 }
 {
  \vlin{\succ i}{{\scriptstyle{i=0,1}}}{\blue{\sn}, \red{\sn}, \n \seqar \n}{\vlin{\cnd_\sq}{\bullet}{\blue{\sn}, \red{\sn}, \n \seqar \n}{\vlhy{\vdots}}}
 }
}
\]
    \caption{Examples of regular coderivations $\incrementation$, $\primrec$ and $ \cconc$, from left (assuming $\gfunction$, $\hfunction_0$ and $\hfunction_1$ regular).}
    \label{fig:examples-regular-coderivations}
\end{figure*}



% \begin{exmp}[Regular coderivations]\label{exmp:regular-coderivations} The coderivations $\incrementation$, $\successor$ and $\cconc$ below are regular.
% \end{exmp}
As discussed in~\cite{Das2021,Das2021-preprint,Kuperberg-Pous21}, coderivations can be identified with Kleene-Herbrand-G\"odel style equational programs, in general computing partial recursive functionals  (see, e.g., \cite[\S 63]{Kleene71:intro-to-metamath} for further details).
We shall specialise this idea to our two-sorted setting.


\begin{defn}
[Semantics of coderivations]
\label{defn:semantics-coderivations}
To each $\bcnorec$-coderivation $\der$ we associate a two-sorted Kleene-Herbrand-G\"odel partial function $\denot \der$ obtained by construing the semantics of \Cref{defn:semantics-bc} as a (possibly infinite) equational program.
% 
Given a two-sorted function $f(\vec x; \vec y)$, we say that $f$ is \emph{defined} by a $\bcnorec$-coderivation $\der$ if $\denot\der(\vec x; \vec y)=f(\vec x; \vec y)$. 
\end{defn}



\begin{rem}\label{rem:regular-finite-equation}
Note, in particular, that from a regular coderivation $\der$ we obtain a \emph{finite} equational program determining $\denot \der$.
Of course, our overloading of the notation $\denot \der$ is suggestive since it is consistent with that of \Cref{defn:semantics-bc}.
\end{rem}

\begin{exmp}[Regular coderivations, revisited] Let us consider the semantics of coderivations $\incrementation$, $\primrec$ and $\cconc$ from~\Cref{fig:examples-regular-coderivations}.  
\begin{itemize}
    \item The  partial functions $\denot{\incrementation \nu}$  are given by the following equational program:
\[
 \arraycolsep=2pt
\begin{array}{rcll}
    \denot{\incrementation_{\epsilon}}  (x;) & = & \denot{\incrementation_{1}}(\denot{\incrementation_{0}}(x;);) \\
   \denot{\incrementation_{0}}(x;) & = & \denot{\incrementation_{00}}(x;)\\
   \denot{\incrementation_{00}}(x;) & = & \denot{\incrementation_{000}}(;x)\\
   \denot{\incrementation_{000}}(;x) & = & \succ 1(;x)\\
   \denot{\incrementation_{1}}(x;) & = & \denot{\incrementation_{\epsilon}}(x;)
\end{array}
\]
By purely equational reasoning, we can simplify this program to obtain $\denot{\incrementation_{\epsilon}}(x;)  =  \denot{\incrementation_{\epsilon}}(\succ 1x;)$. 
% \[
% \begin{array}{rcll}
%       \denot{\incrementation_{\epsilon}}(x;) & = & \denot{\incrementation_{\epsilon}}(\succ 1x;)
% \end{array}
% \]
 Since the above equational program keeps increasing the input, the function $\denot\incrementation=\denot{\incrementation_{\epsilon}}$ is always undefined. 
%     \item  The equational program associated with the coderivation  $\successor$ is as follows:
% \[
% \begin{array}{rcll}
%     \denot{\successor_{\epsilon}}  (0;) & = & \denot{\successor_{0}}(;) \\
%      \denot{\successor_{\epsilon}}  (\succ 0x;) & = & \denot{\successor_{1}}(x;) \\ 
%      \denot{\successor_{\epsilon}}  (\succ 1x;) & = & \denot{\successor_{2}}(x;) \\
%     \denot{\successor_{0}}(;)& = & \succ 10\\
%     \denot{\successor_{1}}(x;)& = & \succ 1x\\
%       \denot{\successor_{2}}(x;)& = & \denot{\successor_{21}}(\denot{\successor_{20}}(x;);)\\
%       \denot{\successor_{20}}(x;)& = &\denot{\successor_{\epsilon}}(x;) \\
%          \denot{\successor_{21}}(x;)& = & \denot{\successor_{210}}(;x) \\ 
%           \denot{\successor_{210}}(;x)& = & 
%          \succ 0(;x)
% \end{array}
% \]
% By equational reasoning again we  obtain:
% \[
% \begin{array}{rcll}
%     \denot{\successor_{\epsilon}}  (0;) & = & \succ 10\\
%      \denot{\successor_{\epsilon}}  (\succ 0x;) & = & \succ 1x & x \neq 0 \\ 
%      \denot{\successor_{\epsilon}}  (\succ 1x;) & = & \succ 0\denot{\successor_{\epsilon}}(x;) 
% \end{array}
% \]
%  Note, therefore, that the function   $\denot\successor=\denot{\successor_{\epsilon}}$ computes the unary successor function $x \mapsto x+1$. 
% \item 
\item Let $\denot{\gfunction}(\vec x; \vec y)$ and $\denot{\hfunction_i}( x, \vec x, z; \vec y)$ ($i=0,1$) be the functions defined by the  regular  $\bcnorec$-coderivations $\gfunction$ and $\hfunction_i$, respectively. Then the equational program for $\primrec$  can be rewritten as follows:
\begin{equation}
    \label{eq:norm-prim-rec-scheme}
     \arraycolsep=2pt
    \begin{array}{rcll}
 \denot{\primrec}(0, \vec x;) & = & \denot{\gfunction}(\vec x; )\\
  \denot{\primrec}(\succ ix,\vec x; ) & = &\denot{\hfunction_i}( x, \vec x, \denot{\primrec}(x, \vec x; ); )  
\end{array}
\end{equation}
which is an instance of a \emph{non-safe}  recursion scheme (on notation).
    \item The equational program of  $\cconc$ can be written as:
\[
 \arraycolsep=2pt
\begin{array}{rcll}
    \denot{\cconc_{\epsilon}}  (0, 0;z) & = & z\\
     \denot{\cconc_{\epsilon}}  ( 0, \succ i y;z) & = &   \succ i \denot{\cconc_{\epsilon}} ( x, y;z) & \neq 0 \\ 
     \denot{\cconc_{\epsilon}}  (\succ i x, y;z) & = &   \succ i\denot{\cconc_{\epsilon}}  ( x, y;z) & \neq 0 
\end{array}
\]
which computes concatenation of the binary representation of three natural numbers.
\end{itemize}

 

 
\end{exmp}
The above examples illustrate two undesirable features of regular $\bcnorec$-coderivations, from the point of view of implicit complexity:
\begin{enumerate}[I.]
    \item \label{enum:problem1} on the one hand, despite being finitely presentable, they can define \emph{partial functions};
    \item \label{enum:problem2} on the other hand, despite the presence of modalities implementing the normal/safe distinction of function arguments, they  can define \emph{non-safe recursion schemes}.
\end{enumerate}

\subsection{The progressing criterion}

To address Problem~\ref{enum:problem1} we shall adapt to our setting a well-known `termination criterion' from non-wellfounded proof theory.
% able to identify certain coderivations (even the non-regular ones) that define total functions. 
First, let us recall some standard  proof theoretic concepts about (co)derivations, similar to those in \cite{Das2021,Das2021-preprint,Kuperberg-Pous21}.
\begin{defn}
[Ancestry]
\label{defn:ancestry}
Fix a coderivation $\der$. We say that a type occurrence $A$ is an \emph{immediate ancestor} of a type occurrence $B$ in $\der$ if they are types in a premiss and conclusion (respectively) of an inference step and, as typeset in \Cref{fig:bc-type-system}, have the same colour.
If $A$ and $B$ are in some $\orange \Gamma$ or $\purple{ \Gamma'}$, then furthermore they must be in the same position in the list.
\end{defn}




Being a binary relation, immediate ancestry forms a directed graph upon which our correctness criterion is built. 

\gianluca{I guess that the figure with the examples must be expanded for readability.}
\begin{defn}
[Progressing coderivations]
\label{defn:progressing}
A \emph{thread} is a maximal path in the graph of immediate ancestry.
We say that a (infinite) thread is \emph{progressing} if it is eventually constant $\sn$ and infinitely often principal for a $\cnd_{\sq}$ rule.

A coderivation is \emph{progressing} if each of its infinite branches has a progressing thread.
\end{defn}

\begin{exmp}[Regular coderivations, re-revisited] In~\Cref{fig:examples-regular-coderivations},  $\incrementation$ has precisely one infinite branch (that loops on $\bullet$) which contains no instances of  $\cnd_{\sq}$ at all. Therefore, $\incrementation$ is not progressing. 
% By contrast, $\successor$ is progressing, since the  infinite branch has a progressing thread on the blue formulas $\blue{\sn}$ indicated there. Last, 
On the other hand, $\cconc$ has two simple loops, one on $\bullet$ and the other one on $\circ$. For any infinite branch $B$ we have two cases:
\begin{itemize}
    \item if $B$ crosses the bottommost conditional infinitely many times, it contains a progressing \blue{blue} thread;
    \item otherwise, $B$ crosses the topmost conditional infinitely many times, so that it  contains  a progressing \red{red} thread.
\end{itemize}
Therefore, $\cconc$ is progressing. By similar reasoning, we can conclude that $\primrec$ is progressing whenever $\gfunction$ and $\hfunction_i$ are.
\end{exmp}




Like in \cite{Kuperberg-Pous21,Das2021-preprint,Das2021}, the progressing criterion is sufficient to guarantee that the partial function computed is, in fact, a well-defined total function:

\begin{prop}
% [Totality]
\label{prop:prog-total}
If $\der$ is progressing then $\denot \der$ is total.
\end{prop}
\begin{proof}
[Proof sketch]
We proceed by contradiction.
If $\denot \der$ is non-total then, since each rule preserves totality top-down, we must have that $\denot{\der'}$ is non-total for one of $\der$'s immediate sub-coderivations $\der'$.
Continuing this reasoning we can build an infinite leftmost `non-total' branch $B=(\der^{i})_{i<\omega}$.
Let $(\sn^i)_{i\geq k}$ be a progressing thread along $B$, and assign to each $\sn^i$ the least natural number $n_i \in \Nat$ such that $\denot{\der^{i}}$ is non-total when $n_i$ is assigned to the type occurrence $\sn^i$.

Now, notice that:
\begin{itemize}
    \item $(n_i)_{i\geq k}$ is monotone non-increasing, by inspection of the rules and their interpretations from \Cref{defn:semantics-bc}.
    \item $(n_i)_{i\geq k}$ does not converge, since $(\sn^i)_{i\geq k}$ is progressing and so is infinitely often principal for $\cnd_\sq$, where the value of $n_i$ must strictly decrease (cf., again, \Cref{defn:semantics-bc}).
\end{itemize}
This contradicts the well-ordering property of the natural numbers.
\end{proof}

One of the most appealing features of the progressing criterion is that, while being rather expressive and admitting many natural programs, e.g.~as we will see in the next subsections, it remains effective (for regular coderivations) thanks to well known arguments in automaton theory:

\begin{fact}[Folklore]
 It is decidable whether a regular coderivation is progressing.
\end{fact}
This well-known result (see, e.g.,~\cite{DaxHL06} for an exposition for a similar circular system) follows from the fact that the progressing criterion is equivalent to the universality of a  B\"{u}chi automaton of size determined by the (finite) representation of the input coderivation. This problem is decidable in polynomial space, though the correctness of this algorithm requires nontrivial infinitary combinatorics, as formally demonstrated in~\cite{KolodziejczykMP19}. 


Let us finally observe that the progressing condition turns out to be sufficient to restate~\Cref{prop:modal-succedent-derivations} in the setting of  non-wellfounded coderivations:


\begin{prop} \label{prop:cutbox} Given a progressing  $\bcnorec$-coderivation  $\der: \sq \Gamma , \vec \n \seqar \sn$, there is a progressing  $\bcnorec$-coderivation $\der^*:\sq\Gamma \seqar \sn$   such that, for any $\vec x$, $\vec y$:
\begin{equation*}
    \model{\der}(\vec{x}; \vec{y})= \model{\der^*}(\vec{x};).
\end{equation*}
\end{prop}
\begin{proof} By progressiveness, any infinite branch contains a $\cnd_\sq$-step, which has non-modal succedent. 
Thus there is a set of $(\cnd_{\sq}, 0, \id)$-occurrences  that forms a bar across $\der$.
% Since the conclusion has a modal succedent, by inspection on the rules of $\bcnorec$ we conclude that each infinite branch contains an instance of $\sqr$. Hence, the set of the conclusions of such $\sqr$-steps forms  a bar across $\der$. 
By K\"{o}nig Lemma, the set of all nodes of $\der$ below this bar, say $X_\der$, is finite. The proof now follows by induction on the size of $X_\der$, and proceeds analogously to~\Cref{prop:modal-succedent-derivations}. If the size is $0$ then the last rule of $\der$ must be an instance of $\cnd_{\sq}$, $0$, or $\id$, in which case the statement holds vacuously. Otherwise, we do case analysis. If the last rule of $\der$ is an instance of $\sqr$ or $\cnd_{\n}$ then the statement holds vacuously. If  the last rule of $\der$ is an instance of $\exch_{\n},\exch_{\sq},  \sql,\succ i$, and $\wk_\sq$ then we apply the induction hypothesis. Let us now suppose that $\der$ has been obtained from a coderivation $\der_0$ by applying an instance of  $\wk_\n$. By induction hypothesis, there exists a coderivation  $\pder0^*:\lists{n}{\sn}{\sn} \seqar \sn$ such that $\model{\pder0}(\vec x; \vec y)= \model{\pder0^*}(\vec x; )$. Since $\model{\der}(\vec x; \vec y, y)= \model{\pder0}(\vec x; \vec y)=\model{\pder0^*}(\vec x; )$ we just set $\der^*= \pder{0}^*$.  Suppose now that $\der$ is obtained from two coderivations $\der_0$ and $\der_1$ by applying an instance of  $\cut_\n$. By induction hypothesis, there exists $\pder{1}^*$ such that $\model{\pder{1}}(\vec x;\vec y,y)=\model{\pder{1}^*}(\vec x;)$. Since 
 $\model{\pder{1}}(\vec x; \vec y,  \model{\pder{0}}(\vec x; \vec y))= \model{\pder{1}^*}(\vec x;)$, we set $\der^*=\pder{1}^*$. As for the case where the last rule is $\cut_\sq$, by induction hypothesis, there exist coderivations $\pder{0}^*$ and $\pder{1}^*$ such that $\model{\pder{0}}(\vec x; \vec y)= \model{\pder{0}^*}(\vec x;)$ and $\model{\pder{1}}(\vec x,x; \vec y)= \model{\pder{1}^*}(\vec x, x;)$, so that we define $\der^*$ as the coderivation obtained from $\pder{0}^*$ and $\pder{1}^*$ by applying the rule $\cut_{\sq}$.
\end{proof}



Note that the above proof only depends on progressiveness, and so it holds for non-regular progressing $\bcnorec$-coderivations as well.





\subsection{Computational expressivity of coderivations} \label{subsec:definability in typing derivations.}



% The sequent calculus system $\bc$ (Figure~\ref{fig:bc-type-system}) is a $S4$-style modal calculus transplanting  Bellantoni and Cook's  function algebra into a proof-theoretical setting (as in~\cite{Hofmann97}). In particular, $\bc$-derivations  locally introduce safe recursion  by means of the rule $\saferec$, and Proposition~\ref{prop:modal-succedent-derivations} ensures that the composition schemes defined by the cut rules are safe. As hinted by~\Cref{enum:problem2}, a different situation arises when we move  $\bcnorec$-coderivations, where the lack of further constraints on non-wellfounded coderivations makes the distinction between modal and non-modal inputs becomes redundant.

% The following example illustrates the problem.
% \begin{exmp}
% [Non-safety in regular progressing coderivations]
% \label{exmp:non-safe-recursion}
% Let $g(\vec x; \vec y)$ and $h_i( x, \vec x, z; \vec y)$ with $i=0,1$ be functions definable by some (regular and progressing) $\bcnorec$-coderivations $\gfunction$ and $\hfunction_i$, respectively. Then the following \emph{non-safe} recursion scheme,
% \begin{equation}
%     \label{eq:norm-prim-rec-scheme}
%     \begin{array}{rcll}
%  f(0, \vec x;) & = & g(\vec x; )\\
%   f(\succ ix,\vec x; ) & = &h_i( x, \vec x, f(x, \vec x; ); )  
% \end{array}
% \end{equation}
% can be defined by the following (regular and progressing) $\bcnorec$-coderivation:
% \[
% \vlderivation{
% \vliin{\cnd_\sq}{\bullet}{\underline{\blue{\sn}}, \sq \vec \n \seqar \n}
% {
% \vltr{\gfunction}{\sq \vec \n \seqar \n}{\vlhy{\ }}{\vlhy{\ }}{\vlhy{\ }}
% }
% {
%   \vliin{\cut_\sq}{{\scriptstyle{i=0,1}}}{\blue{\sn}, \sq \vec \n \seqar \n}
%   {
%   \vlin{\sqr}{}{\blue{\sn}, \sq \vec \n  \seqar \red{\sn}}{
%     \vlin{\cnd_{\sq}}{\bullet}{\blue{\sn}, \sq \vec \n  \seqar {\n}}{\vlhy{\vdots}}
%   }
%   }{
%   \vltr{\hfunction_i}{\blue{\sn}, \sq \vec \n, \red{\sn}  \seqar \n}{\vlhy{\ }}{\vlhy{\ }}{\vlhy{\ }}
%   }
% }
% {
% }
% }
% \]
% Note in particular that we can use the scheme in \Cref{eq:norm-prim-rec-scheme} to define all primitive recursive functions.
% \end{exmp}
% We can go further and show that, without  additional  restrictions on (regular, progressing) coderivations, the distinction between modal and non-modal inputs becomes redundant:
% \begin{prop}\label{prop:modal-nonmodal-distinction}
% For any coderivation $\der: \lists{n}{\sn}{\sn},  \lists{m}{\n}{\n}\seqar \n$ of $\bcnorec$ there is a coderivation $\der^\sq: \lists{n+m}{\sn}{\sn}\seqar \n$ such that:
% \begin{itemize}
%     \item $\denot\der(\vec x; \vec y)= \denot{\der^\sq}(\vec x ,\vec y;)$;
%     \item $\der^\sq$ does not contain instances of $\wk_\n$, $\exch_\n$, $\cut_\n$, $\cnd_\n$.
% \end{itemize}
% Moreover, $\der^\sq$ is regular (resp.~progressing) if $\der$ is.
% \end{prop}
% \begin{proof}
% We construct $\der^\sq$ coinductively. The only  interesting cases are  when $\der$ is $\id$ or it is  obtained from $\pder0:\Gamma \seqar \n$ and $\pder1:\Gamma, \n\seqar A$ by applying a $\cut_\n$-step. Then, $\der^\sq$ is constructed,  respectively, as follows:
% \[
% \vlderivation{
% \vlin{\sql}{}{\sn \seqar \n}{\vlin{\id}{}{\n \seqar \n}{\vlhy{\ }}}}\qquad  \qquad 
%   \vlderivation{
%   \vliin{\cut_\sn}{}{\sq \Gamma \seqar \n}{
%         \vlin{\sqr}{}{\sq \Gamma \seqar \sn}{ \vltr{\der_0^\sq}{\sq \Gamma \seqar \n}{\vlhy{\ }}{\vlhy{\ }}{\vlhy{\ }}}
%     }{
%         \vlin{\sql}{}{\sq \Gamma, \sn\seqar \n}{\vltr{\der_1^\sq}{\sq \Gamma, \n\seqar \n}{\vlhy{\ }}{\vlhy{\ }}{\vlhy{\ }}}
%     }
%     }
%     \qedhere
% \]
% \end{proof}

% As we shall see in the next subsections, the above issues can be avoided by imposing a further global criterion  to recover safe recursion in a non-wellfounded setting. On the other hand,  Proposition~\ref{prop:modal-nonmodal-distinction} turns out to be  particularly useful to illustrate how regularity and progressiveness affect the class of $\bcnorec$-definable functions. Indeed, thanks to the above proposition,

Problem~\ref{enum:problem2} indicates that the modal/non-modal distinction for (progressing) $\bcnorec$-coderivations is somewhat redundant and does not suffice, by itself, to control complexity.  
Indeed, by inspection of~\Cref{fig:bc-type-system} it is not hard to see that we can stepwise replace each $\n$ by $\sn$ in the coderivation while preserving progressiveness, so that all definable functions with normal and safe inputs can be defined by using only normal inputs. 

% and so $\bcnorec$-coderivation $\der$ can be also defined by a  $\bcnorec$-coderivation  that  never makes (essential) use of non-modal formulas and that is progressing (resp.~regular) whenever $\der$ is. 
\begin{prop}\label{prop:modal-nonmodal-distinction}
Let $\der: \lists{n}{\sn}{\sn},  \lists{m}{\n}{\n}\seqar \n$ be a  $\bcnorec$-coderivation. Then,    there exists   a  $\bcnorec$-coderivation $\der^\sq: \lists{n+m}{\sn}{\sn}\seqar \n$    s.t.:
\begin{itemize}
    \item $\denot\der(\vec x; \vec y)= \denot{\der^\sq}(\vec x ,\vec y;)$;
    \item $\der^\sq$ does not contain instances of $\wk_\n$, $\exch_\n$, $\cut_\n$, $\cnd_\n$.
\end{itemize}
Moreover, $\der^\sq$ is regular (resp.~progressing) if $\der$ is.
\end{prop}
% \begin{proof}
% [Proof sketch]
% % [Proof of \Cref{prop:modal-nonmodal-distinction}]
% We construct $\der^\sq$ coinductively. The only  interesting cases are  when $\der$ is $\id$ or it is  obtained from $\pder0:\Gamma \seqar \n$ and $\pder1:\Gamma, \n\seqar A$ by applying a $\cut_\n$-step. Then, $\der^\sq$ is constructed,  respectively, as follows:
% \[
% \small
% \vlderivation{
% \vlin{\sql}{}{\sn \seqar \n}{\vlin{\id}{}{\n \seqar \n}{\vlhy{\ }}}}\qquad  \quad 
%   \vlderivation{
%    \vliin{\cut_\sn}{}{\sq \Gamma \seqar \n}{
%         \vlin{\sqr}{}{\sq \Gamma \seqar \sn}{ \vltr{\der_0^\sq}{\sq \Gamma \seqar \n}{\vlhy{\ }}{\vlhy{\ }}{\vlhy{\ }}}
%     }{
%         \vlin{\sql}{}{\sq \Gamma, \sn\seqar \n}{\vltr{\der_1^\sq}{\sq \Gamma, \n\seqar \n}{\vlhy{\ }}{\vlhy{\ }}{\vlhy{\ }}}
%     }
%     } \qedhere
% \]
% \end{proof}
\begin{proof}
We construct $\der^\sq$ coinductively. We only consider two  interesting cases,  when $\der$ is $\id$ or it is  obtained from $\pder0:\Gamma \seqar \n$ and $\pder1:\Gamma, \n\seqar A$ by applying a $\cut_\n$-step. Then, $\der^\sq$ is constructed,  respectively, as follows:
\[
\small
\vlderivation{
\vlin{\sql}{}{\sn \seqar \n}{\vlin{\id}{}{\n \seqar \n}{\vlhy{\ }}}}\qquad  \quad 
  \vlderivation{
   \vliin{\cut_\sq}{}{\sq \Gamma \seqar \n}{
        \vlin{\sqr}{}{\sq \Gamma \seqar \sn}{ \vltr{\der_0^\sq}{\sq \Gamma \seqar \n}{\vlhy{\ }}{\vlhy{\ }}{\vlhy{\ }}}
    }{
        \vlin{\sql}{}{\sq \Gamma, \sn\seqar \n}{\vltr{\der_1^\sq}{\sq \Gamma, \n\seqar \n}{\vlhy{\ }}{\vlhy{\ }}{\vlhy{\ }}}
    }
    }
\]
the remaining interesting cases are treated similarly.
\end{proof}


Consequently, we may view (regular, progressing) $\bcnorec$-coderivations as the type 0 (regular, progressing) fragment of the system $\ct$ from \cite{Das2021-preprint,Das2021,Kuperberg-Pous21}. 


As a result we inherit the following characterisations:


\begin{prop}
\label{prop:computational-props-coderivations}
We have the following:
\begin{enumerate}
    % \item\label{item:prim-rec-regprogcoder} If $f(\vec x)$ is primitive recursive then $f(\vec x;)$ is defined by a regular and progressing $\bcnorec$-coderivation.
    \item\label{item:type1-complete} Any function $f(\vec x;)$ is defined by a progressing $\bcnorec$-coderivation.
    \item\label{item:turing-complete} The class of regular $\bcnorec$-coderivations is Turing-complete, i.e.\ they define every partial recursive function.
    \item\label{item:type1-prim-rec} $f(\vec x;)$ is defined by a regular progressing $\bcnorec$-coderivation
if and only if 
$f(\vec x)$ is type-1-primitive-recursive, i.e.\ it is in the level 1 fragment $\ntgodel 1$ of G\"odel's $\tgodel$.
\end{enumerate}
\end{prop}

The  proof of the above proposition can be found in~\Cref{sec:computational-strength}.



Given the computationally equivalent system $\ct_0$ with contraction from \cite{Das2021}, we can view the above result as a sort of `contraction admissibility' for regular progressing $\bcnorec$-coderivations. 
Call $\bcnorec+ \{\cntr_\n, \cntr_\sn \}$ the extension of $\bcnorec$ with the rules $\cntr_\n$ and $\cntr_\sn$ below:
\begin{equation} \label{eqn:explicit-contraction}
  \vlinf{\cntr_\n}{}{\orange \Gamma, \blue \n\seqar B}{\orange \Gamma, \blue \n , \blue \n \seqar B}
\qquad 
\vlinf{\cntr_\sq}{}{\blue \sn , \orange \Gamma\seqar B}{\blue \sn,\blue \sn , \orange \Gamma \seqar B}
\end{equation}
 where  the semantics for the new system extends the one for $\bcnorec$ in the obvious way, and the notion of (progressing) thread is induced by the given colouring.\footnote{Note that the totality argument of \Cref{prop:prog-total} still applies in the presence of these rules, cf.~also \cite{Das2021}.} 
% The next proposition shows that, in the setting of regular progressing $\bcnorec$-coderivations,  the rules $\cntr_\n$ and $\cntr_\sn$  do not increase the class of definable functions.
We have:
\begin{cor} $f(\vec x; )$ is definable by a regular progressing  $\bcnorec+ \{\cntr_\n, \cntr_\sn \}$-coderivation  if and only if it is definable by a regular progressing $\bcnorec$-coderivation.
\end{cor}
% \begin{proof}
% [Proof idea]
% The former system is equivalent to $\ct_0$ from \cite{Das2021}, whose type 1 functions are just those of $\ntgodel 1$ by \cite[Corollary 80]{Das2021}, which are all defined by regular progressing $\bcnorec$-coderivations by \Cref{prop:computational-props-coderivations} \Cref{item:type1-prim-rec}.
% \end{proof}

























\subsection{Proof-level conditions motivated by implicit complexity}\label{sec:Proof-level conditions motivated by implicit complexity}
% The sequent calculus system $\bc$ (Figure~\ref{fig:bc-type-system}) is a $S4$-style modal calculus transplanting  Bellantoni and Cook's  function algebra into a proof-theoretical setting (as in~\cite{Hofmann97}). In particular,
$\bc$-derivations  locally introduce safe recursion  by means of the rule $\saferec$, and Proposition~\ref{prop:modal-succedent-derivations} ensures that the composition schemes defined by the cut rules are safe. As suggested by Problem~\ref{enum:problem2}, a different situation arises when we move to   $\bcnorec$-coderivations, where the lack of further constraints 
means that we can define `non-safe' equational programs.
We may recover safety by a natural proof-level condition:
% on non-wellfounded coderivations makes the distinction between modal and non-modal inputs becomes redundant.

% To weaken the computational strength of progressing coderivations, we need to introduce a further criterion on $\bcnorec$-coderivations  to rule out recursion schemes that are `not safe'. \anupam{turn this sentence around: we want to consider only safe equational programs} To this end we propose the following simple criterion:


\begin{defn}
[Safety]
\label{defn:safety}
A $\bcnorec$-coderivation is \emph{safe} if each infinite branch crosses only finitely many $\cut_\sq$-steps.
\end{defn}

The corresponding equational programs of safe coderivations indeed only have safe inputs in hereditarily safe positions, as we shall soon see.
Let us illustrate this by means of examples.

\begin{exmp}
The coderivation $\incrementation$ in~\Cref{fig:examples-regular-coderivations} is not safe, as there is an instance of $\cut_{\sq}$ in the loop on $\bullet$, which means that there is an infinite branch crossing infinitely many $\cut_{\sq}$-steps. By contrast,  the coderivation  $\cconc$ is safe because it does not contain instances of the rule $\cut_\sq$.  Finally, by inspecting the coderivation $\primrec$  of~\Cref{fig:examples-regular-coderivations}, we notice that the infinite branch that loops on $\bullet$  contains infinitely many $\cut_{\sq}$ steps, 
so it  is not safe.
\end{exmp}


Perhaps surprisingly, however, the safety condition is not enough to restrict the set of $\bcnorec$-definable functions to $\fptime$, as the following example shows. 





\begin{exmp}
[Safe exponentiation]
\label{exmp:safe-exponentiation}
Consider the following coderivation $\sexp$,
\[
\vlderivation{
    \vliin{\cnd_{\sq}}{\bullet}{\blue{\sn}, \orange \n \seqar \n}{
        \vlin{\succ{0}}{}{\orange \n \seqar \n}{
        \vlin{\id}{}{\orange \n \seqar \n}{\vlhy{}}
        }
    }{
        \vliin{\cut_\n}{\scriptstyle{i=0,1}}{\blue{\sn}, \orange \n \seqar \n}{
            \vlin{\cnd_{\sq}}{\bullet}{\blue{\sn}, \orange \n \seqar \red \n}{\vlhy{\vdots}}
        }{
            \vlin{\cnd_{\sq}}{\bullet}{\blue{\sn} , \red \n \seqar \n}{\vlhy{\vdots}}
        }
    }
}
\]
where we identify the sub-coderivations above the second and the third premises of the conditional. The coderivation is clearly progressing. Moreover it is safe, as $\sexp$ has no instances of $\cut_{\sq}$. Its  associated equational program can be written as follows:
\begin{equation}
 \arraycolsep=2pt
    \label{eq:sexp-eq-prog}
    \begin{array}{rcll}
    \denot \sexp (0;y) & = & \succ 0 (;y) \\
    \denot \sexp (\succ 0 x;y) & = & \denot \sexp (x; \denot \sexp(x;y))  &\quad  x \neq 0\\
    \denot \sexp (\succ 1 x ; y) & = & \denot \sexp (x; \denot \sexp(x;y))
\end{array}
\end{equation}
The above equational program has already appeared in \cite{Hofmann97,Leivant99}. It is not hard to show, by induction on $x$, that  $\denot\sexp(x;y)=2^{2^{\s{x}}} \cdot  y$.
Thus $f_\sexp$ has exponential growth rate (as long as $y\neq 0$), despite being defined by a `safe' recursion scheme. 
\end{exmp}





The above coderivation exemplifies a  safe recursion scheme that is able to \emph{nest} one recursive call inside another in order to obtain exponential growth rate. This is in fact a peculiar feature of circular proofs, and it is worth discussing. 




\begin{rem}
[On nesting and higher-order recursion]
As we
have seen, namely in~\Cref{prop:computational-props-coderivations}.\ref{item:type1-prim-rec},
(progressing) $\bcnorec$-coderivations are able to simulate some sort of higher-order recursion, namely at type 1 (cf.~also \cite{Das2021}).
In this way it is arguably not so surprising that the sort of `nested recursion' in \Cref{eq:sexp-eq-prog} is definable since type 1 recursion, in particular, allows such nesting of the recursive calls. 
To make this point more apparent, consider the following higher-order `safe' recursion operator:
\begin{equation*}
  \mathsf{rec}_A: \sq{\n}\rightarrow (\sq{\n}\rightarrow A  \rightarrow     A)\rightarrow A \rightarrow A   
\end{equation*}
with $A= \n \rightarrow \n$, and      $f(x)=\mathsf{rec}_A(x,  h, g)$ is defined as $  f(0)=g$ and    $f(\succ i x)=h(x,f (x) )$  for $x>0$. By setting $g\dfn \lambda y:\n. \succ 0y$ and $h\dfn  \lambda x: \sn . \lambda u: \n \rightarrow \n . \lambda y: \n . u(u\, y)$ we can easily check that $\denot\sexp(x; y)= \mathsf{rec}_A(x,  h, g)(y)$, where $\sexp$ is as in Example~\ref{exmp:safe-exponentiation}. Hence, the function $\denot\sexp(x; y)$ can be defined by means of a higher-order version of safe recursion. 

As noticed by Hofmann~\cite{Hofmann97}, and formally proved by Leivant~\cite{Leivant99}, $\felementary$ can be characterised using higher-order safe recursion, thanks to this capacity to nest recursive calls.
Moreover, Hofmann showed in~\cite{Hofmann97} that by introducing a `linearity' restriction on the operator $ \mathsf{rec}_A$, which prevents duplication of recursive calls, it is possible to recover the class $\fptime$. The resulting type system, called $\slr$ (`Safe Linear Recursion'), can thus be regarded as a higher-order formulation of $\bc$.
\end{rem}


Following~\cite{Hofmann97}, we shall impose a linearity criterion to rule out those coderivations that nest recursive calls. This is achieved by observing that the duplication of the recursive calls of $\denot\sexp$ in Example~\ref{exmp:safe-exponentiation} is due to the presence in $\sexp$ of loops on $\bullet$ crossing both premises of a $\cut_\n$ step. Hence, our circular-proof-theoretic counterpart of Hofmann's linearity restriction  can be  obtained  by demanding that only the left premiss of each $\cut_\n$ step is crossed by such loops. 
Again, we rather give a more natural proof-level criterion which does not depend on our intuitive notion of loop.

 

\begin{defn}
[Left-leaning]
\label{defn:left-leaning}
A $\bcnorec$-coderivation is said to be \emph{left-leaning} if each infinite branch goes right at a $\cut_\n$-step only finitely often.
\end{defn}
\begin{exmp}
In~\Cref{fig:examples-regular-coderivations}, $\incrementation$ is trivially left-leaning, as it contains no instances of $\cut_\n$ at all. The  coderivations $\cconc$  and $\primrec$ are also left-leaning, since no  infinite branch can go right at the $\cut_\n$ steps. By contrast, the coderivation $\sexp$ in Example~\ref{exmp:safe-exponentiation} is not left-leaning, as there is an infinite branch looping at $\bullet$ and crossing infinitely many times the rightmost premise of the $\cut_\n$-step.
\end{exmp}


We are now ready to present our circular systems:
\begin{defn}
[Circular Implicit Systems]
\label{defn:cyclic-systems}
$\ncbc$ is the class of safe regular progressing $\bcnorec$-coderivations. 
$\cbc$ is the restriction of $\ncbc$ to only left-leaning coderivations. 
A two-sorted function $f(\vec x;\vec y)$ is \emph{$\ncbc$-definable} (or \emph{$\cbc$-definable}) if there is a coderivation $\der\in \ncbc$ (resp., $\der \in \cbc$) such that $\denot \der (\vec x;\vec y) = f(\vec x;\vec y)$.
\end{defn}


Let us point out that Proposition~\ref{prop:cutbox} can be strengthened to preserve safety and left-leaningness:

\begin{prop} \label{prop:cutbox-for-left-leaning} Let    $\der: \sq \Gamma , \vec \n \seqar \sn$ be a  coderivation in $\ncbc$ (or $\cbc$). There exists a $\ncbc$-coderivation (resp., $\cbc$-coderivation) $\der^*:\sq\Gamma \seqar \sn$  such that
\(
    \model{\der}(\vec{x}; \vec{y})= \model{\der^*}(\vec{x};)
\).
\end{prop}






\subsection{On the complexity of proof-checking}

Note that both the safety and the left-leaning conditions above are defined at the level of arbitrary coderivations, not just regular and/or progressing ones. Moreover, since these conditions are defined at the proof-level rather than the thread-level, they are easy to check on regular coderivations:
\begin{prop} \label{prop:safety-left-leaning-NL} The safety and the left-leaning condition are $\mathbf{NL}$-decidable for regular coderivations.
\end{prop}
\begin{proof} We can represent a regular coderivation $\der$ as a  finite directed (possibly cyclic) graph $G_\der$ labelled with inference rules. Then,  the problem of deciding whether $\der$ is not safe (resp.~left leaning) comes down to the problem of deciding whether  a cycle $\pi$ of $G_\der$  exists crossing a node labelled $\cut_\sq$ (resp.~crossing the rightmost child of a node labelled $\cut_\n$).  W.l.o.g.~we can assume that $\pi$ is a simple cycle. Given (an encoding of)  $G_\der$ as read-only input and (an encoding of) $\pi$ as a read-only  certificate, we can easily construct a deterministic Turing machine $M$ verifying that the cycle $\pi$ crosses $\cut_\sq$ (resp. the rightmost child of a node $\cut_\n$) in $G_\der$. More specifically:
\begin{itemize}
    \item the size of the certificate is smaller than the size of the description of $G_\der$;
    \item  $M$ reads once from left to right the addresses of the certificate which provide the information about where to move the pointers;
    \item $M$ works in logspace as it only stores addresses in memory. \qedhere
\end{itemize}
\end{proof}

The idea here is that, for regular coderivations, checking that no branch has infinitely many occurrences of a particular rule can be reduced to checking acyclicity of a certain subgraph, which is well-known to be in $\mathbf{coNL}=\mathbf{NL}$.
% \begin{proof}[Proof of~\Cref{prop:safety-left-leaning-NL}.] We can represent a regular coderivation $\der$ as a  finite directed (possibly cyclic) graph $G_\der$ labelled with inference rules. Then,  the problem of deciding whether $\der$ is not safe (resp.~left leaning) comes down to the problem of deciding whether  a cycle $\pi$ of $G_\der$  exists crossing a node labelled $\cut_\sq$ (resp.~crossing the rightmost child of a node labelled $\cut_\n$).  W.l.o.g.~we can assume that $\pi$ is a simple cycle. Given (an encoding of)  $G_\der$ as read-only input and (an encoding of) $\pi$ as a read-only  certificate, we can easily construct a deterministic Turing machine $M$ verifying that the cycle $\pi$ crosses $\cut_\sq$ (resp. the rightmost child of a node $\cut_\n$) in $G_\der$. More specifically:
% \begin{itemize}
%     \item the size of the certificate is smaller than the size of the description of $G_\der$;
%     \item  $M$ reads once from left to right the addresses of the certificate which provide the information about where to move the pointers;
%     \item $M$ works in logspace as it only stores addresses in memory. \qedhere
% \end{itemize}
% \end{proof}




Recall that progressiveness of regular coderivations is decidable by reduction to universality of B\"uchi automata, a $\pspace$-complete problem.
Indeed progressiveness itself is $\pspace$-complete in many settings, cf.~\cite{NolletST19}.
It is perhaps surprising, therefore, that the safety of a regular coderivation also allows us to decide progressiveness efficiently too, thanks to the following reduction:
\begin{prop}\label{prop:regular-safe-progressing-iff-infinite-branch-infinite-cond}
A  safe  $\bcnorec$-coderivation is progressing iff every infinite branch has infinitely many $\cnd_\sq$-steps.
\end{prop}
\begin{proof}
The left-right implication is trivial. For the right-left implication, let us consider an infinite branch $B$ of a  safe $\bcnorec$-coderivation $\der$. By safety, there exists  a node $\nu$ of $B$ such that any sequent  above $\nu$ in $B$ is not the conclusion of a $\cut_\sq$-step. Now, by inspecting the rules of $\bcnorec \setminus \{ \cut_\sq \}$ we  observe that:
\begin{itemize}
    \item every modal formula occurrence in $B$ has a unique thread along $B$;
    \item infinite threads along $B$ cannot start strictly above $\nu$. 
\end{itemize}
Hence, setting $k$ to be the number of $\sn$ occurrences in the antecedent of $\nu$, $B$ has (at most) $ k$ infinite threads. Moreover, since $B$ contains infinitely many $\cnd_\sq$-steps, by the Infinite Pigeonhole Principle we conclude that  one of these  threads is infinitely often principal for the $\cnd_\sq$ rule. 
\end{proof}
% \begin{proof}
% % [Proof of~\Cref{prop:regular-safe-progressing-iff-infinite-branch-infinite-cond}.]
% The left-right implication is trivial. For the right-left implication, let us consider an infinite branch $B$ of a  safe $\bcnorec$-coderivation $\der$. By safety, there exists  a node $\nu$ of $B$ such that any sequent  above $\nu$ in $B$ is not the conclusion of a $\cut_\sq$-step. Now, by inspecting the rules of $\bcnorec \setminus \{ \cut_\sq \}$ we  observe that:
% \begin{itemize}
%     \item every modal formula occurrence in $B$ has a unique thread along $B$;
%     \item infinite threads along $B$ cannot start strictly above $\nu$. 
% \end{itemize}
% Hence, setting $k$ to be the number of $\sn$ occurrences in the antecedent of $\nu$, $B$ has (at most) $ k$ infinite threads. Moreover, since $B$ contains infinitely many $\cnd_\sq$-steps, by the Infinite Pigeonhole Principle we conclude that  one of these  threads is infinitely often principal for the $\cnd_\sq$ rule. 
% \end{proof}




% This result follows from the observation that each branch in a safe coderivation has only finitely many distinct $\sn$-threads. 
Thus, using similar reasoning to that of~\Cref{prop:safety-left-leaning-NL}  we may conclude from~\Cref{prop:regular-safe-progressing-iff-infinite-branch-infinite-cond} the following:
\begin{cor} \label{rem:NL} Given a regular  $\bcnorec$-coderivation $\der$, the problem of deciding if $\der$ is in $\ncbc$ (resp.~$\cbc$)  is in  $\mathbf{NL}$.
\end{cor}

Let us point out that the reduction above is similar to (and indeed generalises) an analogous one for cut-free extensions of Kleene algebra, cf.~\cite[Proposition 8]{DP18}.



% \begin{rem}\label{rem:NL-checkability-with-ecplicit-contraction}
% Recall that in Proposition~\ref{prop:modal-nonmodal-distinction} we showed by means of an indirect argument that explicit contraction is ‘admissible' for regular progressing $\bcnorec$-coderivations, and in  Remark~\ref{rem:explicit-contraction} we pointed out that a more direct proof of this result cannot easily be found. 
% In fact, we could not easily find such a direct contraction-admissibility even for $\cbc$ or $\ncbc$.
% However, let us point out that~\Cref{prop:regular-safe-progressing-iff-infinite-branch-infinite-cond} still holds when explicit contraction is added to $\bcnorec$.
% On the one hand, this observation  witnesses the robustness of the safety criterion; on the other hand, it stresses that, whatever class of functions $\cbc$ and $\ncbc$ define in the presence explicit contraction, the resulting circular proof systems are still $\mathbf{NL}$-checkable (cf.~\Cref{rem:NL}). A detailed argument for this can be found in \Cref{sec:prf-nl-checkability-with-cntr}.
% \end{rem}

\section{Some variants of safe recursion}\label{sec:some-variants}
In this section we shall introduce various extensions of $\bc$ to ultimately classify the expressivity of the circular systems $\cbc$ and  $\ncbc$. 
First, starting from the analysis of   Example~\ref{exmp:safe-exponentiation} and  the subsequent system $\ncbc$, we shall define a version of $\bc$ with safe \emph{nested} recursion, called $\nbc$. 
Second, motivated by the more liberal way of defining functions in both $\cbc$ and $\ncbc$, we shall endow the function algebras $\bc$ and $\nbc$ with forms of safe recursion over a well-founded relation $\permpref$ on lists of normal parameters.  Figure~\ref{fig:function-algebras-table} summarises the function algebras considered and their relations.


\begin{figure}[t]
    \[
    \small
    % \hspace{4cm}
     % \arraycolsep=3pt
   \def\arraystretch{1.3}
    \begin{array}{|c|c|c|}
    \hline
     \textit{\textbf{safe recursion}} & \textit{ on notation} &\textit{ on $\permpref$ }\\       \hline
        \textit{ unnested}& \qquad \bc\qquad  & \qquad \bcpp \qquad \\
        % \textit{ unnested, with compositions}&\sbc  & \sbcpp \\
        \textit{ nested} &  \nbc  & \nbcpp \\ \hline 
    \end{array}
    \]
        \caption{The function algebras considered in \Cref{sec:some-variants}. Any algebra is included in one below it and to the right of it.}
    \label{fig:function-algebras-table}
\end{figure}


\subsection{Relativised algebras and nested recursion}
One of the key features of the Bellantoni-Cook algebra $\bc$ is that `nesting' of recursive calls is not permitted.
For instance, let us recall the equational program from \Cref{exmp:safe-exponentiation}:
\begin{equation}
 \arraycolsep=2pt
    \label{eq:safe-exp-by-snrec}
    \begin{array}{rcl}
    \ex(0;y) & = & \succ 0 y \\
    \ex(\succ i x ; y) & = & \ex(x;\ex(x;y))
\end{array}
\end{equation}
Recall that $\ex(x;y)=f_\sexp (x;y) = 2^{2^{\s{x}}}\cdot y$.
The `recursion step' on the second line is compatible with safe composition, in that safe inputs only occur in hereditarily safe positions, but one of the recursive calls takes another recursive call among its safe inputs.
% 
In Example~\ref{exmp:safe-exponentiation}  we showed how  the above function  $\ex(x;y)$ can be   $\ncbc$-defined. 
We thus seek a suitable extension of $\bc$ able to formalise such nested recursion to serve as a function algebraic counterpart to $\ncbc$. 

% \subsection{Algebras with oracles and safe nested recursion}
It will be convenient for us to work with generalisations of  $\bc$ including \emph{oracles}.  
% Formally speaking, these can be seen as variables ranging over two-sorted functions, though we shall often treat them as explicit functions too.


\begin{defn}\label{defn:oracles-functions}
For all sets of oracles $\vec a = a_1, \dots, a_k$, we define the algebra of $\bcnorec$ functions \emph{over} $\vec a$ to include all the initial functions of $\bcnorec$ and,
\begin{itemize}
    \item (oracles). $a_i(\vec x;\vec y)$ is a function over $\vec a$, for $1\leq i\leq k$, (where $\vec x$,$\vec y$ have appropriate length for $a_i$).
\end{itemize}
and closed under:
\begin{itemize}
% \item (Safe recursion on notation).
% From $g(\vec x;\vec y)$ over oracles $\vec a$ and $h(x,\vec x; \vec y,y)$ over oracles $\vec a$ define $f(x,\vec x;\vec y) $ over oracles $\vec a$ by:
% \[
% \begin{array}{rcll}
%     f(0,\vec x; \vec y) & = & g(\vec x; \vec y)  \\
%     f(\succ 0 x, \vec x; \vec y) & = & h_0 (x, \vec x; \vec y , f(x, \vec x ; \vec y)) & x \neq 0 \\
%     f(\succ 1 x, \vec x; \vec y) & = & h_1 (x, \vec x; \vec y , f(x, \vec x ; \vec y))
% \end{array}
% \]
    \item (Safe Composition). 
    \begin{enumerate}[(1)]
        \item \label{enum:safe-composition-condition-1} from $g(\vec x;\vec y), h(\vec x;\vec y,y)$ over $\vec a$ define $f(\vec x;\vec y)$ over $\vec a$  by $f(\vec x;\vec y) = h(\vec x; \vec y, g(\vec x;\vec y))$.
        \item \label{enum:safe-composition-condition} from $g(\vec x;)$ over $\emptyset$ and $h(\vec x, x;\vec y) $ over $\emptyset$ define $f(\vec x; \vec y)$ over $\vec a$ by $f(\vec x; \vec y) = h(\vec x, g(\vec x;); \vec y)$.
    \end{enumerate}
\end{itemize}
We write $\bcnorec(\vec a)$ for the class of functions over $\vec a$ generated in this way.
\end{defn}


Henceforth, we may write $f(\vec a)(\vec x; \vec y)$ to stress that the function $f(\vec x; \vec y)$ is over oracles $\vec a$.

\begin{rem}\label{rem:variables-vs-functions}
Formally, a function $f(\vec a)(\vec x; \vec y)$ can be understood as a type 2 functional that takes type 1 inputs $\vec a$ (safe-normal functions), and type 0 inputs $\vec x$ (normal) and $\vec y$ (safe). With little abuse of notation, we may represent this functional as $\lambda \vec a. f(\vec a)(\vec x; \vec y)$, where $\vec a$ are now \emph{variables} ranging over first-order functions (with appropriate sorts). Notice that this justifies replacing an oracle with another first-order function (with matching sorts). However, to avoid any reference to higher-order functions, in this paper we prefer to adopt the parametric notation $f(\vec a)(\vec x; \vec y)$ where oracles are arbitrary but fixed \emph{functions}. 
\end{rem}


Note that Safe Composition along normal parameters (\Cref{enum:safe-composition-condition} above) comes with the condition that $g(\vec x; )$ is oracle-free. This restriction prevents oracles appearing in normal position. Since our recursion schemes replace oracles with recursive calls, the latter will also be required to appear in safe position. 
The same condition on $h(\vec x, x; \vec y)$ being oracle-free is not strictly necessary for the complexity bounds we are after, as we shall see in the next section when we define more expressive algebras, but is convenient in order to facilitate the `grand tour' strategy of this paper (cf.~\Cref{fig:picture-main-results}).

We shall write, say, $\lambda \vec v .f(\vec x; \vec v)$ for the function taking only safe arguments $\vec v$ with $(\lambda \vec v .f(\vec x; \vec v))(;\vec y) = f(\vec x; \vec y)$ (here $\vec x$ may be seen as parameters).
Nested recursion can be formalised in the setting of algebras-with-oracles as follows:

\begin{defn}
[Safe Nested Recursion]
\label{defn:safe-nested-recursion}
We write $\snrec$ for the scheme:
\begin{itemize}
    \item from $g(\vec x; \vec y)$ over $\emptyset$ and $h_i(a)(x,\vec x;\vec y)$ over $a,\vec a$, define $f(x, \vec x;\vec y) $ over $\vec a$ by:
    \[
     \arraycolsep=2pt
    \begin{array}{rcll}
     f(0, \vec x; \vec y) &=&g(\vec x;\vec y)
    \\
        f(\succ 0 x, \vec x; \vec y) &=&h_0(\lambda \vec v . f(x,\vec x;\vec v))(x,\vec x;\vec y) &\quad  x \neq 0 \\
        f(\succ 1 x, \vec x; \vec y) &=&h_1(\lambda \vec v . f(x,\vec x;\vec v))(x,\vec x;\vec y)
    \end{array}
    \]
    % over oracles $\vec a$.
    % \anupam{Shouldn't $g$ take oracles $\vec a$ too?}
    % \gianluca{No, it's quite like with circular systems: safety prevents loops to cross the first premise of $\cnd_\sq$. It's crucial for the completeness theorem $\nbc \subset\ncbc$.}
\end{itemize}
We write $\nbc(\vec a)$ for the class of functions over $\vec a$ generated from $\bcnorec$ under $\snrec$ and Safe Composition (from \Cref{defn:oracles-functions}), and write simply $\nbc$ for $\nbc(\emptyset)$.
\end{defn}


% \red{
% \begin{rem}\label{rem:fast-growing}
% W.l.o.g.~we can assume that the base function  of $\snrec$ has an empty set of oracles. Indeed, given $f(x, \vec x; \vec y)$ defined from  $g(\vec x; \vec y)$ over oracles $\vec a$  and $h(a)(x,\vec x;\vec y)$ over oracles $a,\vec a$, we set $\hat{f}(x, \vec x; \vec y)$ defined by $\snrec$ as follows:
%  \[
%     \begin{aligned}
%      \hat{f}(0, \vec x; \vec y) &=0
%     \\
%         \hat{f}(\succ i x, \vec x; \vec y) &=\cnd(;x, g(\vec x; \vec y), h(\lambda \vec v . \hat{f}(x,\vec x;\vec v))(\pred(x),\vec x;\vec y) ) 
%     \end{aligned}
%     \]
% It is easy to check that $\hat{f}(x, \vec x; \vec y)=f(\pred(x), \vec x; \vec y)$ for $x>0$,  so that we set $f(x, \vec x; \vec y)\dfn \hat{f}(\succ 1 x, \vec x; \vec y)$. 
% \end{rem}
% }



% \begin{rem}
% Notice that w.l.o.g.~we can assume that the base function $g$ in the Safe nested recursion has an empty set of oracles. Indeed, given a function $f(x, \vec x; \vec y)$ over oracles $\vec a$ defined by safe nested recursion from the function $g(\vec x; \vec y)$ over oracles $\vec a$ and the step function $h(a)(x, \vec x; \vec y)$ over oracles $\vec a, a$, we can define a function $f^*(x, \vec x; \vec y)$ over oracles $\vec a$ such that $ f^*(0, \vec x; \vec y) =0$  and $   f^*(\succ i x, \vec x; \vec y) =h(\lambda \vec v . f^*(x,\vec x;\vec v))(x,\vec x;\vec y) $. It is easy to show that $f(x, \vec x; \vec y)= \cnd(; x, g(\vec x; \vec y),  f^*(x, \vec x; \vec y), f^*(x, \vec x; \vec y))$.  
% \end{rem}




% Note that Safe Nested Recursion also admits definitions that are not morally `nested' but rather use a form of `composition during recursion':
% \begin{itemize}
%     \item from $g(\vec x; \vec y), \left(\vec g_j(x, \vec x; \vec y)\right)_{1 \leq j \leq k}$ and  $h_i\left(x, \vec x; \vec y, (z_j)_{1 \leq j \leq k}\right)$  define:
%     \begin{equation}
%         \label{eq:safe-rec-w-comp-dur-rec-no-oracles}
%          \arraycolsep=2pt
%     \begin{array}{rcll}
%          f(0, \vec x; \vec y) &=&g(\vec x;\vec y)\\
%          f(\succ i x , \vec x; \vec y) &=& 
%     h_i \left(x, \vec x; \vec y, \left(f(x, \vec x; \vec g_j(x, \vec x; \vec y))\right)_{1 \leq j \leq k}\right)
%         \end{array}
%     \end{equation}
% \end{itemize}
% In this case, note that we have allowed the safe inputs of $f$ to take arbitrary values given by previously defined functions, but at the same time $f$ never calls itself in a safe position, as in~\eqref{eq:safe-exp-by-snrec}.
% This safe recursion scheme does not require the use of oracles in its statement, but is equivalent to the following one using a notion of `nesting' 
% for Safe Composition:
% \begin{defn}
% [Safe Recursion with Composition During Recursion]
% \label{defn:safe-rec-with-comp-dur-rec}
% An instance of Safe Composition along Safe Parameters  (\Cref{enum:safe-composition-condition-1} from \Cref{defn:oracles-functions}) is called \emph{unnested} if (at least) one of $g(\vec x;\vec y)$ and $h(\vec x; \vec y, y)$ is over $\emptyset$ (i.e.\ is oracle-free). 
% % 
% We define $\sbc(\vec a)$ to be the restriction of $\nbc(\vec a)$ using only unnested Safe Composition along Safe Parameters, 
% % \Cref{enum:safe-composition-condition-1} from \Cref{defn:oracles-functions}, 
% and write simply $\sbc$ for $\sbc(\emptyset)$.
% \end{defn}

% It is not hard to see that, indeed,   $\sbc$ defines the same class of functions of  the extension of $\bcnorec$ by the scheme given in \Cref{eq:safe-rec-w-comp-dur-rec-no-oracles}.
% The `$\mathsf S$' in $\sbc$ thus indicates that the recursion scheme is `shallow', as opposed to nested.
% % \anupam{remainder of paragraph can be omitted} 
% %  On the one hand, given $h_i$ and $\vec{g}_j$, the scheme in~\eqref{eq:safe-rec-w-comp-dur-rec-no-oracles} can be defined in $\nbc(\vec a)$ using $\snrec$  and  only unnested instances of safe composition. On the other hand, given $h(a)(x,\vec x;\vec y)$ over oracles $a,\vec a$, the scheme $\snrec$ allows us to define  $f(\succ i x, \vec x; \vec y)$ by replacing the oracle $a$ in $h$ with the recursive calls $\lambda \vec v . f(x,\vec x;\vec v)$. However, since $h(a)$  is defined using unnested safe composition, no recursive call of $f$ can be found along the safe parameters of another one.
% % 
% It should be said that we will not actually use $\sbc$ in this work, but rather an extension of it defined in the next subsection, but we have included it for completeness of the exposition.

\begin{rem}[Safe Composition During Safe Recursion]
\label{rem:errata}
    Note that Safe Nested Recursion also admits variants that are not morally `nested' but rather use a form of `composition during recursion':
\begin{itemize}
    \item from $g(\vec x; \vec y), \left(\vec g_j(x, \vec x; \vec y)\right)_{1 \leq j \leq k}$ and  $h_i\left(x, \vec x; \vec y, (z_j)_{1 \leq j \leq k}\right)$  define:
    \begin{equation}
        \label{eq:safe-rec-w-comp-dur-rec-no-oracles}
         \arraycolsep=2pt
    \begin{array}{rcll}
         f(0, \vec x; \vec y) &=&g(\vec x;\vec y)\\
         f(\succ i x , \vec x; \vec y) &=& 
    h_i \left(x, \vec x; \vec y, \left(f(x, \vec x; \vec g_j(x, \vec x; \vec y))\right)_{1 \leq j \leq k}\right)
        \end{array}
    \end{equation}
\end{itemize}
In this case, note that we have allowed the safe inputs of $f$ to take arbitrary values given by previously defined functions, but at the same time $f$ never calls itself in a safe position, as in~\eqref{eq:safe-exp-by-snrec}. 
In the conference version of this paper \cite[p.8]{CurziDas:lics22} we made an unsubstantiated claim that a function algebra extending $\bcnorec$ by \eqref{eq:safe-rec-w-comp-dur-rec-no-oracles} above is equivalent to a restriction $\sbc$ of $\nbc$ requiring that oracles are \emph{unnested} in Safe Composition: in \cref{enum:safe-composition-condition-1}, one of $g(\vec x;\vec y)$ and $h(\vec x;\vec y,y)$ must be oracle-free, i.e.\ over $\emptyset$. 
Such an equivalence is in fact not immediate, as our Safe Composition scheme allows substitution only along a single parameter, so there does not seem to be a direct way of expressing~\Cref{eq:safe-rec-w-comp-dur-rec-no-oracles} in $\nbc$ without nesting oracles, in the sense just described. 
While we did not rely on this equivalence for any of the main results of \cite{CurziDas:lics22}, in this work we simply omit reference to $\sbc$ and its variants to avoid any confusion.
% One could modify our Safe Composition schema to perform  substitution along multiple parameters, an alternative presentation of composition that is widely used in the literature (and, indeed, it  appears in the original presentation of $\bc$ (see~\cite{BellantoniCook})). However, to avoid any confusion arisen by our claim, we decided to omit references to $\sbc$ and related variants in this paper, as they play no essential role in establishing our results.

% introduced two more function algebras based on~\Cref{eq:safe-rec-w-comp-dur-rec-no-oracles}, $\sbc$


% This safe recursion scheme does not require the use of oracles in its statement, but is equivalent to the following one using a notion of `nesting' 
% for Safe Composition.

%  In the conference version of this paper (see~\cite{CurziDas}) we introduced two more function algebras based on~\Cref{eq:safe-rec-w-comp-dur-rec-no-oracles}, $\sbc$ and its extension to safe recursion on permutation of prefixed  $\sbcpp$, and we claimed that they can be equivalently defined from $\nbc$ and $\nbcpp$ by requiring that oracles never nest. This claim is imprecise, as to avoid nesting of oracles 






% \begin{defn}
% [Safe Recursion with Composition During Recursion]
% \label{defn:safe-rec-with-comp-dur-rec}
% An instance of Safe Composition along Safe Parameters  (\Cref{enum:safe-composition-condition-1} from \Cref{defn:oracles-functions}) is called \emph{unnested} if (at least) one of $g(\vec x;\vec y)$ and $h(\vec x; \vec y, y)$ is over $\emptyset$ (i.e.\ is oracle-free). 
% % 
% We define $\sbc(\vec a)$ to be the restriction of $\nbc(\vec a)$ using only unnested Safe Composition along Safe Parameters, 
% % \Cref{enum:safe-composition-condition-1} from \Cref{defn:oracles-functions}, 
% and write simply $\sbc$ for $\sbc(\emptyset)$.
% \end{defn}

% It is not hard to see that, indeed,   $\sbc$ defines the same class of functions of  the extension of $\bcnorec$ by the scheme given in \Cref{eq:safe-rec-w-comp-dur-rec-no-oracles}.
% The `$\mathsf S$' in $\sbc$ thus indicates that the recursion scheme is `shallow', as opposed to nested.
% % \anupam{remainder of paragraph can be omitted} 
% %  On the one hand, given $h_i$ and $\vec{g}_j$, the scheme in~\eqref{eq:safe-rec-w-comp-dur-rec-no-oracles} can be defined in $\nbc(\vec a)$ using $\snrec$  and  only unnested instances of safe composition. On the other hand, given $h(a)(x,\vec x;\vec y)$ over oracles $a,\vec a$, the scheme $\snrec$ allows us to define  $f(\succ i x, \vec x; \vec y)$ by replacing the oracle $a$ in $h$ with the recursive calls $\lambda \vec v . f(x,\vec x;\vec v)$. However, since $h(a)$  is defined using unnested safe composition, no recursive call of $f$ can be found along the safe parameters of another one.
% % 
% It should be said that we will not actually use $\sbc$ in this work, but rather an extension of it defined in the next subsection, but we have included it for completeness of the exposition.
\end{rem}


\subsection{Safe recursion on well-founded relations}
Relativised function algebras may be readily extended by recursion on arbitrary well-founded relations.
For instance, given a well-founded preorder $\wfpoeq$, and writing $\wfpo$ for its strict variant,\footnote{To be precise, for a preorder $\wfpoeq$ we write $x \wfpo y  $ if $x\wfpoeq y$ and $ y \not\wfpoeq x$. As abuse of terminology, we say that $\wfpoeq$ is well-founded just when $\wfpo$ is.}
`safe recursion on $\wfpo$' is given by the scheme:
\begin{itemize}
    \item from $h(a)(x,\vec x;\vec y)$ over $a,\vec a$, define $f(x, \vec x;\vec y)$ over $\vec a$ by: 
    \[ f(x,\vec x; \vec y) = h(\lambda v \wfpo x . f(v,\vec x;\vec y ))(x,\vec x; \vec y )\] 
\end{itemize}
Note here that we employ the notation $\lambda v \wfpo x$ for a `guarded abstraction'. Formally:
\[
(\lambda v \wfpo x . f(v,\vec x; \vec y))(z)
\dfn
\begin{cases}
f(z, \vec x; \vec y) & z \wfpo x \\
0 & \text{otherwise}
\end{cases}
\]
It is now not hard to see that total functions (with oracles) are closed under the recursion scheme above, by reduction to induction on the well-founded relation $\wfpo$. 


Note that such schemes can be naturally extended to preorders on \emph{tuples} of numbers too, by abstracting several inputs.
% 
We shall specialise this idea to a particular well-founded preorder that will be helpful later to bound the complexity of definable functions in our systems $\cbc$ and $\ncbc$.

Recall that we say that $x$ is a \emph{prefix} of $y$ if $y$ has the form $xz$ in binary notation, i.e.\ $y $ can be written $x2^n + z$ for some $n\geq 0$ and some $z<2^n$.
We say that $x$ is a \emph{strict prefix} of $y$ if $x$ is a prefix of $y$ but $x \neq y$.
% \begin{itemize}
%     \item $x=y$; or,
%     \item $x$ is a prefix of $\pred y$. 
% \end{itemize}


\anupam{Some example earlier should highlight why the permuation of prefixes order is important for our considerations}
\gianluca{Recall example about "max" function, which merge two standard recursions into 1 on two parameters. Perhaps not the best example but it is nstructive.}

\begin{defn}
[Permutations of prefixes]
\label{defn:perm-pref}
Let $[n]$ denote $\{0,\dots, n-1\}$.
We write $(x_0, \dots, x_{n-1}) \permprefeq (y_0, \dots, y_{n-1})$ if, for some permutation $\pi:[n] \to [n] $, we have that $x_i$ is a prefix of $y_{\pi i}$, for all $i<n$.
We write $\vec x \permpref \vec y$ if $\vec x \permprefeq \vec{y} $ but $ \vec y \not\permprefeq \vec{x}$, i.e.\ there is a permutation $\pi : [n]\to [n]$ with $x_i$ a prefix of $y_i$ for each $i<n$ and, for some $i<n$, $x_i$ is a strict prefix of $y_i$.
\end{defn}

It is not hard to see that $\permprefeq$ is a well-founded preorder, by reduction to the fact that the prefix relation is a well-founded partial order.
As a result, we may duly devise a version of safe (nested) recursion on $\permpref$:

\begin{defn}
[Safe (nested) recursion on permutations of prefixes]
\label{defn:saferec-pp}
We write $\nbcpp(\vec a)$ for the class of functions over $\vec a$ generated from $\bcnorec$ under Safe Composition, the scheme $\snrecpp$,
\begin{itemize}
    \item from $h(a)(\vec x; \vec y)$ over $a, \vec a$ define $f(\vec x;\vec y)$ over $\vec a$ by:
    \[
    f(\vec x; \vec y) = 
    h(\lambda \vec u \permpref \vec x, \lambda \vec v. f(\vec u; \vec v))(\vec x; \vec y)
    \]
\end{itemize}
and
the following generalisation of Safe Composition along a Normal Parameter:
\begin{itemize}
 \item[(2)$'$] \label{item:safe-composition-normal-param-with-nonempty-oracles} from $g(\vec x;)$ over $\emptyset$ and $h(\vec x, x;\vec y) $ over $\vec a$ define $f(\vec x; \vec y)$ over $\vec a$ by $f(\vec x; \vec y) = h(\vec x, g(\vec x;); \vec y)$.
\end{itemize}
% Recalling \Cref{defn:safe-rec-with-comp-dur-rec},
% we write $\sbcpp(\vec a)$ to be the restriction of   $\nbcpp(\vec a)$  using only unnested Safe Composition along Safe Parameters.
We define $\bcpp(\vec a)$ to be the restriction of $\nbcpp(\vec a)$
where every instance of $\snrecpp$ has the form:
\begin{itemize}
    \item from $h(a)(\vec x; \vec y)$ over $a, \vec a$, define $f(\vec x;\vec y) $ over $\vec a$:
    \[
    f(\vec x; \vec y) = 
    h(\lambda \vec u \permpref \vec x , \lambda \vec v \permprefeq \vec y . f(\vec u; \vec v))(\vec x ; \vec y)
    \]
\end{itemize}
We call this latter recursion scheme $\srecpp$, e.g.\ if we need to distinguish it from $\snrecpp$.
\end{defn}
Note that the version of safe composition along a normal parameter above differs from the previous one, \Cref{enum:safe-composition-condition} from \Cref{defn:oracles-functions}, since the function $h$ is allowed to use oracles.
Again, this difference is inessential in terms of computational complexity, as we shall see. However, as we have mentioned, the greater expressivity of $\bcpp$ and $\nbcpp$ will facilitate our overall strategy for characterising $\cbc$ and $\ncbc$, cf.~\Cref{fig:picture-main-results}.

% Again, let us note that $\bc(\vec a)$ contains only total functions (with oracles).
% 
% Combining safe nested recursion and safe recursion on the well-founded relation $\subset$, we obtain:

% \begin{defn}
% [Safe nested recursion on permutations of prefixes]
% We define $\nbcpp (\vec a)$ just like $\bcpp(\vec a)$ from \Cref{defn:saferec-pp}, but instead of $\srecpp$ using the scheme $\snrecpp$:
% \begin{itemize}
%     \item from $h(a)(\vec x; \vec y)$ over $a, \vec a$ define $f(\vec x;\vec y)$ over $\vec a$ by:
%     \[
%     f(\vec x; \vec y) = 
%     h(\lambda \vec u \permpref \vec x, \lambda \vec v. f(\vec u; \vec v))(\vec x; \vec y)
%     \]
% \end{itemize}
% Recalling \Cref{defn:safe-rec-with-comp-dur-rec},
% we write $\sbcpp(\vec a)$ to be the restriction of   $\nbcpp(\vec a)$  using only unnested Safe Composition along Safe Parameters.
% \end{defn}
Let us take a moment to point out that $\nbcpp(\vec a)\supseteq  \bcpp(\vec a)$ indeed contain only well-defined total functions over the oracles $\vec a$, by reduction to induction on $\permpref$.



% \anupam{
% Is this a further constraint on the definition of $\nbc$ or is it a consequence of the definition? The latter would be nice, or perhaps a slight weakening of the statement that nonetheless makes the translation from $\nbc$ to $\ncbc$ go through. E.g., can we make the translation work by assuming that all oracles only use safe inputs? I do not immediately see how this would avoid the issue of safety, but I will reflect some more.
% }
% \anupam{.....alright, I think I've got it: I think we can show that every function $f(\vec x;\vec y)$ in $\nbc$ can be written as $f'(\vec x, f_1(\vec x;), \dots, f_k(\vec x;); \vec y)$ where $f'(\vec x,\vec z;\vec y)$ satisfies the condition above. I think it should work in the case of recursion because, crucially, oracles take only safe inputs. I will check this at some point.}

% \anupam{I have commented an attempted proof below. The recursion case needs to be thought about.}
% % \begin{prop}
% % Let  $f(a)(\vec x; \vec y)\in \nbc(a, \vec a)$.  There exist $n_f \geq 0$,   $f^*( \vec x; \vec y, z_1, \ldots, z_{n_f}) \in \nbc( \vec a)$ and a family $(\vec{f}_{i}( \vec x; \vec y))_{1 \leq i\leq n_f}\in \nbc(a, \vec a) $ for some $n$ such that:
% % \[
% % f(\lambda \vec v. a( \vec x; \vec v))( \vec x; \vec y)=f^*( \vec x; \vec y, a( \vec x; \vec f_1(\lambda \vec v. a( \vec x; \vec v))( \vec x; \vec y)), \ldots, a( \vec x;  \vec f_{n_f}(\lambda \vec v. a( \vec x; \vec v))( \vec x; \vec y)))
% % \]
% % \end{prop}
% % \begin{proof}
% % By induction on the construction of $f$. If $f$ is an oracle among $a, \vec a$ then we define $f^*$ as a projection function. The case where $f$ is an initial function is straightforward. If $f(\vec x; \vec y)=h(\vec x; \vec y, g(\vec x; \vec y))$ then we set $n_f= n_h+n_g$,  $f^*(\vec x; \vec y, z_1, \ldots, z_{n_f})=h^*(\vec x; \vec y,g^*(\vec x; \vec y, z_1, \ldots, z_{n_g}),  z_1, \ldots, z_{n_h})$ and $\vec f_i(\vec x; \vec y)= \vec g_i(\vec x; \vec y)$ for all $1 \leq i \leq n_g$ and $\vec f_i(\vec x; \vec y)= \vec h_i(\vec x; \vec y, g(\vec x; \vec y))$. Let us suppose that $f(\vec x; \vec y)= h(\vec x, g(\vec x;); \vec y)$, where $g$ has no oracles. Then we have two cases: if $h$ has no oracles, then we are done. Otherwise, since the normal parameters 
% % \end{proof}

% \red{
% Let us write $\nbc'$ for the algebra obtained by using the version of Safe Composition on Normal Parameters above. (Note that, clearly, $\nbc' \subseteq \nbc$). 
% Notably, we have:
% \begin{lem}
% Suppose $f(\vec x;\vec y) \in \nbc(\vec a)$ where each $a_i$ takes only safe inputs. 
% Then there is some $f'(\vec x,\vec z;\vec y) \in \nbc' (\vec a)$ and $f_1(\vec x;), \dots, f_k(\vec x;) \in \nbc'(\emptyset)$ such that:
% \[
% f(\vec x;\vec y) = f'(\vec x, \vec f(\vec x;);\vec y)
% \]
% In particular, if $f(\vec x; \vec y) \in \nbc(\emptyset)$ then $f(\vec x; \vec y)\in \nbc'(\emptyset)$ (by safe composition).
% \end{lem}
% \begin{proof}
% The case of initial functions and oracles are immediate.

% Suppose $f(\vec x;\vec y) = h(\vec x; \vec y, g(\vec x; \vec y))$ and $g'(\vec x,\vec w;\vec y),h'(\vec x, \vec z;\vec y, y),\vec g(\vec x;),\vec h(\vec x;)$ are given by the inductive hypothesis. 
% Then by safe composition in $\nbc'$ we can set,
% \[
% f'(\vec x,\vec w, \vec z;\vec y) = h'(\vec x,\vec h(\vec x;) ; \vec y, g'(\vec x,\vec w; \vec y))
% \]
% and also $\vec f = (\vec g,\vec h)$.

% Suppose $f(\vec x; \vec y) = h(\vec x, g(\vec x;);\vec y)$.
% Since $g$ has no oracles, by the inductive hypothesis and Safe Composition we have that $g(\vec x;)\in \nbc'(\emptyset).$
% Let $h'(\vec x, x , \vec z;\vec y), \vec h(\vec x,x;\vec y)$ be given by the inductive hypothesis for $h$. We simply set,
% \[
% f'(\vec x, z,\vec z;\vec y) = h' (\vec x,z,\vec z;\vec y)
% \]
% and $\vec f = (g,\vec g)$.

% Finally, suppose $f(x, \vec x; \vec y)$ is given by safe nested recursion,
% \[
% \begin{array}{rcl}
%     f(0,\vec x;\vec y) & = & g(\vec x;\vec y)\\
%     f(\succ 0 x, \vec x; \vec y) &= & h_0(\lambda \vec v. f(x,\vec x;\vec v))(x, \vec x; \vec y)\\
%     f(\succ 1 x, \vec x; \vec y) &= & h_1(\lambda \vec v. f(x,\vec x;\vec v))(x, \vec x; \vec y)
% \end{array}
% \]
% PROBLEM HERE: the $\vec g,\vec h_i$ will take $\succ i x$ as input, not $x$.
% \end{proof}
% }
% However,  a similar restriction can not be restated for the function algebras $\bcpp$, $\nbcpp$, and $\sbcpp$, due to an essential difference  between the recursion schemes $\snrec$ and $\snrecpp$. Indeed, the former has a  single  recursion parameter, which is explicitly decreased by one. By contrast, in the latter the whole set of  normal arguments is regarded as a recursive parameter, and the way such a parameter   decreases  (w.r.t.~$\subset$) depends on how $h(a)$ is defined. To see this, consider a function $f(\vec x; \vec y)$  defined by $\snrecpp$  from a given function $h(a)(\vec x; \vec y)$, i.e.
% \[ f(\vec x; \vec y) = 
%     h(\lambda \vec u \permpref \vec x, \lambda \vec v. f(\vec u; \vec v))(\vec x; \vec y)
%     \]
% by definition of guarded abstraction, in order to apply the recursive call of $f$ to a set of normal parameters $\vec z$ we require $\vec z \subset \vec x$, which can only happen if a predecessor $\pred$ is applied to some normal argument, i.e.~if $h(a)$ is defined using an instance of safe composition in Definition~\ref{defn:oracles-functions} (point~\ref{enum:safe-composition-condition}) with $\vec a\neq \emptyset$.

\subsection{Simultaneous recursion schemes}
For our main results, we will ultimately need a typical property of function algebras, that $\bcpp$ and $\nbcpp$ are closed under simultaneous versions of their recursion schemes.


\begin{defn}[Simultaneous schemes]
\label{defn:simultaneous-schemes}
We define schemes $\ssrecpp$ and $\ssnrecpp$, respectively, as follows, for arbitrary $\vec a = a_1, \dots, a_{k}$:
\begin{itemize}
    \item from $h_i(\vec a)(\vec x; \vec y)$ over $\vec a, \vec b$, for $1 \leq i\leq k$, define $f_i(\vec x;\vec y)$ over $\vec b$, for $1 \leq i \leq k$, by:
    \[
    f_i(\vec x; \vec y) = 
    h_i((\lambda \vec u \permpref \vec x, \lambda \vec v\subseteq \vec y. f_j(\vec u; \vec v))_{1 \leq j\leq  k})(\vec x; \vec y)
    \]
    \item from $h_i(\vec a)(\vec x; \vec y)$ over $\vec a, \vec b$, for $1\leq i \leq k$, define $f_i(\vec x;\vec y)$ over $\vec b$, for $1 \leq i \leq  k$, by:
    \[
    f_i(\vec x; \vec y) = 
    h_i((\lambda \vec u \permpref \vec x, \lambda \vec v. f_j(\vec u; \vec v))_{1 \leq j \leq k})(\vec x; \vec y)
    \]
\end{itemize}
\end{defn}
\begin{prop} \label{prop:simultaneous-recursion-admissible}
We have the following:
\begin{enumerate}
    \item\label{item:sim-rec-admiss-bcpp} If $\vec f(\vec x;\vec y)$ over $\vec b$ are obtained by applying  $\ssrecpp$ to $\vec h(\vec a)(\vec x;\vec y) \in \bcpp(\vec a,\vec b)$, then also $\vec f(\vec x;\vec y) \in \bcpp (\vec b)$.
\item\label{item:sim-rec-admiss-nbcpp-sbcpp}
If $\vec f(\vec x;\vec y)$ over $\vec b$ are obtained by applying  $\ssnrecpp$ to $\vec h(\vec a)(\vec x;\vec y) \in \nbcpp(\vec a,\vec b)$, then also $\vec f(\vec x;\vec y) \in \nbcpp (\vec b)$.
\end{enumerate}
\end{prop}


\begin{proof}
% [Proof of~\Cref{prop:simultaneous-recursion-admissible}.]
We only prove \Cref{item:sim-rec-admiss-bcpp}, i.e. that $\bcpp$ is closed under $\ssrecpp$, but the same argument works for \Cref{item:sim-rec-admiss-nbcpp-sbcpp}: just ignore all guards on safe inputs in recursive calls.


Let $f_i(\vec x;\vec y)$ and $h_i(a_1, \dots, a_k)(\vec x; \vec y)$ be as given in \Cref{defn:simultaneous-schemes}, and temporarily write $f_j^{\vec x;\vec y}$ for $\lambda \vec u \permpref \vec x, \lambda \vec v \permprefeq \vec y . f_j (\vec u;\vec v)$, so we have:
\[
f_i(\vec x;\vec y) = h_i (f_1^{\vec x; \vec y}, \dots, f_k^{\vec x;\vec y})(\vec x; \vec y)
\]
For $i\in \Nat$, let us temporarily write $\numeral i$ for $i$ in binary notation\footnote{In fact, any notation will do, but we pick one for concreteness.}, and $\permi i$ for the list $\numeral i, \numeral{i+1}, \dots, \numeral k, \numeral 1, \numeral 2, \dots, \numeral{i-1}$.
Note that, for all $i=1, \dots, k$, $\permi i$ is a permutation (in fact a rotation) of $\numeral 1, \dots, \numeral k$.

Now, let $f(\vec x; \vec y, \vec z)$ over oracles $\vec b$ be given as follows:
\begin{equation}
    \label{eq:sim-rec-reduction}
    f(\vec x; \vec y, \vec z) \dfn \begin{cases}
h_1 (f_1^{\vec x;\vec y}, \dots, f_k^{\vec x; \vec y})(\vec x;\vec y) & \vec z = \permi 1 \\
\vdots \\
h_k (f_1^{\vec x;\vec y}, \dots, f_k^{\vec x; \vec y})(\vec x;\vec y) & \vec z = \permi k \\
0 & \text{otherwise}
\end{cases}
\end{equation}
Note that this really is a finite case distinction since each of the boundedly many $\permi i$ has bounded size, both bounds depending only on $k$, and so is computable in $\bcnorec$ over $\vec h$. 

By definition, then, we have that $f(\vec x;\vec y, \permi i) = f_i (\vec x; \vec y)$.
Moreover note that, for each $j=1,\dots,k$, we have,
\[
\arraycolsep=2pt
\begin{array}{rcl}
   f_j^{\vec x;\vec y} (\vec u'; \vec v')  &=& (\lambda \vec u \permpref \vec x, \lambda \vec v \permprefeq \vec y. f_j(\vec u;\vec v)) (\vec u';\vec v') \\
    & = & (\lambda \vec u \permpref \vec x, \lambda \vec v \permprefeq \vec y. f(\vec u;\vec v,\permi j ))(\vec u';\vec v') \\
    & = & (\lambda \vec u \permpref \vec x, \lambda \vec v \permprefeq \vec y, \lambda \vec w \permprefeq \vec z . f(\vec u;\vec v, \vec w)) (\vec u';\vec v',\permi j)
\end{array}
\]
% \gianluca{A very small comment: from the viewpoint of the lambda calculus, shouldn't the last function in the above chain of  equations have shape $\lambda \vec u \permpref \vec x, \lambda \vec v \permprefeq \vec y, (\lambda \vec w \permprefeq \vec z . f(\vec u;\vec v, \vec w))(\permi j) $? In this case: can we infer the equation below?}
% \anupam{Is this better? I am glossing over the exchange issue a little, since all functions are type 1 so I can assume all inputs are fed. Perhaps we should say more, or give a footnote.}
as long as $\vec z$ is some $\vec i$,
so indeed~\eqref{eq:sim-rec-reduction} has the form,
\[
f(\vec x; \vec y, \vec z) = h (\lambda \vec u \permpref \vec x, \lambda \vec v \permprefeq \vec y, \lambda \vec w \permprefeq \vec z . f(\vec u; \vec v, \vec w)) (\vec x;\vec y)
\]
and $f(\vec x;\vec y,\vec z) \in \bcpp(\vec b)$ by $\srecpp$.
Finally, since  $f_i(\vec x;\vec y) = f(\vec x;\vec y, \permi i)$, we indeed have that each $f_i(\vec x;\vec y) \in \bcpp(\vec b)$.
\end{proof}
\section{Characterisations for function algebras}\label{sec:characterization-results-for-function-algebras}
In this section we characterise the complexities  of the function algebras we introduced in the previous section. 
Namely, despite apparently extending $\bc$, $\bcpp$ still contains just the polynomial-time functions, whereas  both  $\nbc$ and $\nbcpp$ are shown to contain just the elementary functions. 
All such results rely on a `bounding lemma' inspired by~\cite{BellantoniCook}.




\subsection{A relativised Bounding Lemma}
\label{sec:bounding-lemma}
Bellantoni and Cook showed in \cite{BellantoniCook} that any function $f(\vec x;\vec y) \in \bc$ satisfies the `poly-max bounding lemma': there is a polynomial $p_f(\vec n)$ such that:\footnote{Recall that, for $\vec x = x_1, \dots, x_n$, we write $|\vec x|$ for $|x_1|, \dots, |x_n|$.}
\begin{equation}
    \label{eq:bc-polymax-eqn}
    |f(\vec x;\vec y)| \leq p_f (|\vec x|) + \max |\vec y| 
\end{equation}
This provided a suitable invariant for eventually proving that all $\bc$-functions were polynomial-time computable.

In this work, inspired by that result, we generalise the bounding lemma to a form suitable to the relativised algebras from the previous section. 
To this end we establish in the next result a sort of `elementary-max' bounding lemma that accounts for the usual poly-max bounding as a special case, by appealing to the notion of (un)nested safe composition.
Both the statement and the proof are quite delicate due to our algebras' formulation using oracles; we must assume an appropriate bound for the oracles themselves, and the various (mutual) dependencies in the statement are subtle.


To state and prove the Bounding Lemma, let us employ the notation $\sumlen{\vec x} \dfn \sum |\vec x|$. 
% \begin{lem}[Bounding Lemma]
% \label{lem:boundinglemma}
% % \marginpar{\footnotesize See proof at \ref{prf:boundinglemma}.}
% Let $f(\vec a)(\vec x;\vec y) \in \nbcpp(\vec a)$, with $\vec a = a_1, \dots, a_k$.
% There is a function $m_f^{\vec c}(\vec x,\vec y)$ with,
% \[
% m_f^{\vec c}(\vec x,\vec y) = 
% e_f(\sumlen{\vec x}) + d_f\sum\vec c + \max |\vec y|
% \]
% for an elementary function $e_f(n)$ and a constant $d_f \geq 1$, such that whenever there are constants $\vec c=c_1, \dots, c_k$ such that,
% \begin{equation}
%     \label{eq:oracle-const-max-bound}
%     |a_i (\vec x_i;\vec y_i)| \leq c_i + d_f \sum_{j\neq i} c_j + \max|\vec y_i|
% \end{equation}
% % for some constants $\vec c =c_1, \ldots, c_k$, then 
% for $1\leq i \leq k$,
% we have:\footnote{To be clear,  we write $|\vec y_i| \leq m_f^{\vec c}(\vec x,\vec y)  $ here as an abbreviation for $\{|y_{ij}| \leq m_f^{\vec c}(\vec x,\vec y)\}_j$.}
% \begin{equation}
%     \label{eq:elem-const-max-bound}
%     |f(\vec a)(\vec x;\vec y)| \leq m_f^{\vec c}(\vec x,\vec y)
% \end{equation}
% \begin{equation}
%     \label{eq:input-bounding-eqn}
%     f(\vec a)(\vec x; \vec y) \ = \ f\left(\lambda \vec x_i.\lambda |\vec y_i| \leq m_f^{\vec c}(\vec x,\vec y) . a_i(\vec x_i;\vec y_i)\right)_i(\vec x;\vec y)
% \end{equation}
% % \anupam{this last equation a little difficult to parse.}
% % where \red{$m_f^{\vec c}(\vec x,\vec y)$ is shorthand notation for:
% % \[
% % e_f(\sumlen{\vec x}) + d_f\sum\vec c + \max |\vec y|
% % \]
% % }
% Moreover, if in fact $f(\vec x;\vec y) \in \bcpp(\vec a) $, then $d_f=1$ and $e_f(n)$ is a polynomial. 
% \end{lem}

\begin{lem}[Bounding Lemma]
\label{lem:boundinglemma}
% \marginpar{\footnotesize See proof at \ref{prf:boundinglemma}.}
Let $f(\vec a)(\vec x;\vec y) \in \nbcpp(\vec a)$, with $\vec a = a_1, \dots, a_k$, and suppose $g_1, \ldots, g_k$ are safe-normal functions. Then there is an elementary function $e_f(n)$ and a constant $d_f \geq 1$ such that whenever there are constants $\vec c=c_1, \dots, c_k$ satisfying, for $1\leq i \leq k$,
\begin{equation}
    \label{eq:oracle-const-max-bound}
    |g_i (\vec u;\vec v)| \leq c_i + \left(d_f \sum_{j\neq i} c_j \right) + \max|\vec v| \qquad \qquad \text{for all inputs } \vec u,\vec v
\end{equation}
% for some constants $\vec c =c_1, \ldots, c_k$, then 
then we have the following, for all inputs $\vec u, \vec v$\footnote{To be clear,  we write $|\vec y| \leq m_f(\vec c,\vec u,\vec v)  $ here as an abbreviation for $\{|y_{j}| \leq m_f(\vec c,\vec u,\vec v)\}_j$.}:
\begin{equation}
    \label{eq:elem-const-max-bound}
    |f(\vec g)(\vec u;\vec v)| \leq m_f(\vec c, \vec u,\vec v)
\end{equation}
\begin{equation}
    \label{eq:input-bounding-eqn}
    f(\vec g)(\vec u; \vec v) \ = \ f\left(\lambda \vec x.\lambda |\vec y| \leq m_f(\vec c, \vec u,\vec v) . g_i(\vec x;\vec y)\right)_{1 \leq i \leq k}(\vec u;\vec v)
\end{equation}
where 
\[
m_f(\vec c,\vec x,\vec y) = 
e_f(\sumlen{\vec x}) + d_f\sum\vec c + \max |\vec y|
\]
% \anupam{this last equation a little difficult to parse.}
% where \red{$m_f^{\vec c}(\vec x,\vec y)$ is shorthand notation for:
% \[
% e_f(\sumlen{\vec x}) + d_f\sum\vec c + \max |\vec y|
% \]
% }
Moreover, if in fact $f(\vec x;\vec y) \in \bcpp(\vec a) $, then $d_f=1$ and $e_f(n)$ is a polynomial. 
\end{lem}

Notice that in the above statement we employ the notation $\lambda \vec x, \lambda|\vec y|\leq n$, where $\vec y=y_1, \ldots, y_k$, for the guarded abstraction defined by:
\[
(\lambda \vec x, \lambda|\vec y|\leq n. h(\vec x; \vec y))(\vec w; \vec z):= \begin{cases}
    h(\vec w; \vec z) &\text{if } |z_j|\leq n \text{ for all }1 \leq j \leq k\\
    0 &\text{otherwise.}
\end{cases}
\]


Unwinding the statement above, note that $e_f$ and $d_f$ depend \emph{only} on the function $f$ itself, not on the constants $\vec c$ given for the (mutual) oracle bounds in \Cref{eq:oracle-const-max-bound}.
This is crucial for the proof, namely in the case when $f$ is defined by recursion, substituting different values for $\vec c$ during an inductive argument.
% This independency of $e_f$ and $d_f$ from $\vec c$ will be made explicit in their definitions in the proof that follows.

While the role of the elementary bounding function $e_f$ is a natural counterpart of $p_f$ in Bellantoni and Cook's bounding lemma, cf.~\Cref{eq:bc-polymax-eqn}, the role of $d_f$ is perhaps slightly less clear.
Intuitively, $d_f$ represents the amount of `nesting' in the definition of $f$, increasing whenever oracle calls are substituted into arguments for other oracles. 
Hence, if $f$ uses only unnested Safe Composition, then $d_f=1$ as required. 
In fact, it is only important to distinguish whether $d_f=1$ or not, since $d_f$ forms the base of an exponent for defining $e_f$ when $f$ is defined by safe recursion.

Finally let us note that~\Cref{eq:elem-const-max-bound,eq:input-bounding-eqn} are somewhat dual:  while the former bounds the \emph{output} of a function (serving as a modulus of \emph{growth}), the latter bounds the \emph{inputs} (serving as a modulus of \emph{continuity}).


\begin{rem}
Let us also point out that we may relativise the statement of the lemma to any set of oracles including those which $f(\vec x;\vec y)$ is over.
In particular, if $f(\vec x;\vec y)$ is over no oracles, then we may realise \Cref{eq:oracle-const-max-bound} vacuously by choosing $\vec a = \emptyset$ and we would obtain that $|f(\vec x;\vec y)| \leq e_f(\sumlen{\vec x}) + \max |\vec y|$.
More interestingly, in the case when $f(\vec x;\vec y)$ is just, say, $a_i(\vec x;\vec y)$, we may choose to set $\vec a = a_i$ or $\vec a = a_1,\dots, a_n$ in the above lemma, yielding different bounds in each case.
We shall exploit this in inductive hypotheses  in the proof that follows (typically when we write `WLoG').
\end{rem}

\begin{proof}[Proof of~\Cref{lem:boundinglemma}.]
\label{prf:boundinglemma}
We prove~\Cref{eq:elem-const-max-bound} and~\Cref{eq:input-bounding-eqn} by induction on the definition of $f(\vec x;\vec y)$, always assuming that we have $\vec c$ satisfying \Cref{eq:oracle-const-max-bound}.

Throughout the argument we shall actually construct $e_f( n)$ that are \emph{monotone} elementary functions (without loss of generality) and exploit this invariant.
In fact $e_f( n)$ will always be generated by composition from $0, \succ{}, +, \times,  x^y$ and projections.
If $f(\vec x;\vec y) \in \bcpp(\vec a)$ then $e_f(n)$ will be in the same algebra without exponentiation, $ x^y$, i.e.\ it will be a polynomial with only non-negative coefficients.
This property will be made clear by the given explicit definitions of $e_f( n)$ throughout the argument.


Let us start with~\Cref{eq:elem-const-max-bound}. If $f(\vec x;\vec y)$ is an initial function then it suffices to set $e_f(n) \dfn 1 +  n$ and $d_f\dfn 1$.

If $f(\vec x;\vec y) = a_i(\vec x;\vec y)$ then it suffices to set $e_f(n) \dfn 0 $ and $d_f \dfn 1$.

If $f(\vec x;\vec y) = h(\vec x;\vec y, g(\vec x;\vec y))$, let $e_h, e_g, d_h, d_g$ be obtained from the inductive hypothesis.
We have,
\[
\arraycolsep=2pt
\begin{array}{rcl}
    |f(\vec x;\vec y)| & = & |h(\vec x; \vec y, g(\vec x; \vec y))| \\
    & \leq & e_h(\sumlen{\vec x}) + d_h\sum\vec c + \max(|\vec y|, |g(\vec x;\vec y)|) \\
    & \leq & e_h(\sumlen{\vec x}) + d_h\sum\vec c + e_g(\sumlen{\vec x}) + d_g\sum \vec c + \max|\vec y|
    % \\
    % & \leq & (e_h(\sumlen{\vec x}) + e_g(\sumlen{\vec x})) + (d_h+d_g)\sum \vec c + \max |\vec y
\end{array}
\]
so we may set $e_f(n) \dfn e_h(n) + e_g(n)$ and $d_f\dfn d_h+d_g$.

For the `moreover' clause, note that if $f(\vec x;\vec y)$ uses only unnested Safe Composition, then one of $g$ or $h$ does not have oracles, and so we can assume WLoG that either the term $d_h\sum \vec c$ or the term $d_g\sum \vec c$ above does not occur, and we set $d_f \dfn d_g$ or $d_f\dfn d_h$, respectively. In either case we obtain $d_f=1$ by the inductive hypothesis, as required.

If $f(\vec x;\vec y) = h(\vec x, g(\vec x;);\vec y)$, let $e_h,e_g,d_h, d_g$ be obtained from the inductive hypothesis. Note that, by definition of Safe Composition along a normal parameter, we must have that $g$ has no oracles, and so in fact $|g(\vec x;)| \leq e_g(\sumlen{\vec x})$.
We thus have,
\[
\arraycolsep=2pt
\begin{array}{rcl}
    |f(\vec x;\vec y)| & = & |h(\vec x, g(\vec x;);\vec y)| \\
     & \leq & e_h(\sumlen{\vec x}, |g(\vec x;)|) + d_h\sum \vec c + \max|\vec y| \\
     & \leq & e_h(\sumlen{\vec x}, e_g(\sumlen{\vec x})) + d_h\sum \vec c + \max|\vec y| 
\end{array}
\]
so we may set $e_f(n)\dfn e_h(n, e_g(n))$ and $d_f \dfn d_h$.
% 
For the `moreover' clause, note that by the inductive hypothesis $e_h,e_g$ are polynomials and $d_h=1$, so indeed $e_f$ is a polynomial and $d_f=1$.

Finally, if $f(\vec x;\vec y) = h(\lambda \vec u \permpref \vec x, \lambda \vec v. f(\vec u;\vec v) ) (\vec x;\vec y)$, let $e_h,d_h$ be obtained from the inductive hypothesis.
We claim that it suffices to set $d_f \dfn d_h$ and $e_f(n) \dfn nd_h^n e_h(n)$.
Note that, for the `moreover' clause, if $d_h=1$ then also $d_h^n = 1$ and so indeed $e_f(n)$ is a polynomial if $d_h=1$ and $e_h(n)$ is a polynomial.

First, let us calculate the following invariant, for $n>0$:
\begin{equation*}
    \arraycolsep=2pt
    \begin{array}{rcll}
    e_f(n) & = & nd_h^n e_h(n) &  \\
        & = & d_h^n e_h(n) + (n-1)d_h^n e_h(n) & \\
        & = & d_h^n e_h(n) + d_h (n-1) d_h^{n-1}e_h(n) \\
        & \geq & e_h(n) + d_h (n-1) d_h^{n-1}e_h(n) & \text{$d_h\geq 1$} \\
        & \geq & e_h(n) + d_h (n-1) d_h^{n-1}e_h(n-1) \  &  \text{$e_h$~monotone} \\
        & \geq & e_h(n) + d_h e_f(n-1) & \text{def.~of $e_f$}
\end{array}
\end{equation*}
Hence:
\begin{equation}
 \label{eq:bnd-lem-rec-fn-inv}
     e_f(n) \geq  e_h(n) + d_h e_f(n-1) 
\end{equation}




Now, to show that \Cref{eq:elem-const-max-bound} bounds the $f,e_f,d_f$ at hand, we proceed by a sub-induction on $\sumlen{\vec x}$.
For the base case, when $\sumlen{\vec x} = 0 $ (and so, indeed, $\vec x = \vec 0$), note simply that $\lambda \vec u\permpref \vec x, \lambda \vec v . f(\vec u;\vec v)$ is the constant function $0$, and so we may appeal to the main inductive hypothesis for $h(a)$ setting the corresponding constant $c$ for $a$ to be $0$ to obtain,
\[
\arraycolsep=2pt
\begin{array}{rcl}
    |f(\vec 0;\vec y)| & = & |h(0)(\vec 0;\vec y)| \\
        & \leq & e_h(0) + d_h\sum\vec c + \max |\vec y| \\
        & \leq & e_f(0) + d_f\sum\vec c + \max |\vec y|
\end{array}
\]
as required.
For the sub-inductive step,  let $\sumlen{\vec x} > 0$. 
Note that, whenever $\vec u \permpref \vec x$ we have $\sumlen{\vec u} < \sumlen{\vec x} $ and so, by the sub-inductive hypothesis and monotonicity of $e_f$ we have:
\[
|f(\vec u;\vec v)| \leq e_f(\sumlen {\vec x} -1) + d_f \sum \vec c + \max |\vec v|
\]
Now we may again appeal to the main inductive hypothesis for $h(a)$ by setting $c= e_f(\sumlen{\vec x}  -1)$ to be the corresponding constant for $a = \lambda \vec u \permpref \vec x, \lambda \vec v. f(\vec u;\vec v)$. 
We thus obtain:
\[
\arraycolsep=2pt
\begin{array}{rclll}
     |f(\vec x;\vec y)| & = & |h(\lambda \vec u \permpref \vec x, \lambda \vec v . f(\vec u;\vec v)) (\vec x;\vec y)| \\
         & \leq & e_h (\sumlen{\vec x}) + d_h c + d_h\sum\vec c + \max |\vec y| & & \text{main IH} \\
         & \leq & ( e_h (\sumlen{\vec x}) + d_h e_f(\sumlen{\vec x} -1) ) +\\
         & &  + d_h\sum\vec c + \max |\vec y| & & \text{def.~of $c$}\\
         & \leq & e_f(\sumlen{\vec x}) + d_h \sum \vec c + \max|\vec y| && \eqref{eq:bnd-lem-rec-fn-inv} \\ & \leq & e_f(\sumlen{\vec x}) + d_f \sum \vec c + \max|\vec y| && \text{$d_f = d_h$}
\end{array}
\]
Let us now prove~\Cref{eq:input-bounding-eqn}, and let $e_f$ and $d_f$ be constructed by induction on $f$ as above. We proceed again by induction on the definition of $f(\vec a)(\vec x;\vec y)$, always making explicit the oracles of a function.
% Let us temporarily write $m_f(\vec x,\vec y) \dfn e_f(\sumlen{\vec x}) + d_f\sum\vec c + \max |\vec y|$.
\gianluca{Is it better to define $e_f$ and $d_f$ in advance before proving the two inequations?}

The initial functions and oracle calls are immediate, due to the `$\max |\vec y|$' term in \Cref{eq:input-bounding-eqn}. 

If $f(\vec a)(\vec x; \vec y) = h(\vec a)(\vec x; \vec y, g(\vec a)(\vec x; \vec y))$ then, by the inductive hypothesis for $h(\vec a)$, any oracle call from $h(\vec a)$ only takes safe inputs of lengths:
    \[
    \begin{array}{rll}
        \leq & e_h(\sumlen{\vec x}) + d_h \sum \vec c + \max (|\vec y|, |g(\vec a)(\vec x;\vec y)|) \\
        \leq & e_h(\sumlen{\vec x})+ d_h \sum \vec c + e_g(\sumlen{\vec x}) + d_g\sum \vec c + \max |\vec y| & \text{\eqref{eq:oracle-const-max-bound}} \\
        \leq &  (e_h(\sumlen{\vec x} + e_g(\sumlen{\vec x})) + (d_h +d_g)\sum\vec c + \max |\vec y|\\
         \leq & e_f(\sumlen{\vec x}) + d_f\sum\vec c + \max|\vec y|
    \end{array}
    \]
    Note that any oracle call from $g(\vec a)$ will still only take safe inputs of lengths $\leq e_g(\sumlen{\vec x}) + d_g\sum \vec c + \max |\vec y|$, by the inductive hypothesis, and $e_g$ and $d_g$ are bounded above by $e_f$ and $d_f $ respectively.
    
If $f(\vec a)(\vec x;\vec y) = h(\vec a)(\vec x, g(\emptyset)(\vec x;);\vec y)$ then, by the inductive hypothesis, any oracle call will only take safe inputs of lengths:
\[
\arraycolsep=2pt
\begin{array}{rlll}
    \leq & e_h(\sumlen{\vec x} + |g(\emptyset)(\vec x;)|) + d_h\sum\vec c + \max |\vec y| \\
    \leq & e_h(\sumlen{\vec x} + e_g(\sumlen{\vec x})) + d_h\sum\vec c + \max |\vec y| & & \text{\Cref{lem:boundinglemma}} \\
    \leq & e_f(\sumlen{\vec x}) + d_f\sum \vec c + \max |\vec y|
\end{array}
\]

Last, suppose $f(\vec a)(\vec x;\vec y)= h(\vec a, \lambda \vec u\permpref \vec x, \lambda \vec v . f(\vec a)(\vec u;\vec v))(\vec x;\vec y)$. 
We proceed by a sub-induction on $\sumlen{\vec x}$.
Note that, since $\vec u \permpref \vec x \implies \sumlen{\vec u} <\sumlen{\vec x}$, we immediately inherit from the inductive hypothesis the appropriate bound on safe inputs for oracle calls from $\lambda \vec u\permpref \vec x, \lambda \vec v. f(\vec a)(\vec u;\vec v)$.

Now, recall from the Bounding \Cref{lem:boundinglemma}, whenever $\vec u \permpref \vec x$ (and so $\sumlen{\vec u} < \sumlen{\vec x}$), we have $|f(\vec u;\vec v)| \leq e_f(\sumlen{\vec x}-1) + d_f\sum \vec c + \max |\vec v|$.
So by setting $c = e_f(\sumlen{\vec x}-1)$ in the inductive hypothesis for $h(\vec a,a)$, with $a = \lambda \vec u\permpref \vec x, \lambda \vec v. f(\vec a)(\vec u;\vec v)$, any oracle call from $h(\vec a,a)$ will only take safe inputs of lengths:
\[
\arraycolsep=2pt
\begin{array}{rlll}
    \leq & e_h(\sumlen{\vec x}) + d_h e_f(\sumlen{\vec x}-1) + d_h\sum \vec c + \max |\vec y| \\
    \leq & e_f(\sumlen{\vec x}) + d_f \sum \vec c + \max |\vec y| & & \text{\eqref{eq:bnd-lem-rec-fn-inv}}
\end{array}
\]
% Now, for an oracle call from $\lambda \vec u\permpref \vec x, \lambda \vec v. f(\vec a)(\vec u;\vec v)$ we may simply apply the same reasoning again, essentially proceeding by induction on $\sumlen{\vec x}$.
This completes the proof.  
\end{proof}



% Before moving on to our soundness results, let us first present an important consequence of the bounding lemma that will be important for characterising $\nbcpp$.

% \begin{prop}
% % \label{input-elem-bound}
% Let $f(\vec a)(\vec x; \vec y) \in \nbcpp (\vec a) $ and let $e_f$ and $d_f$ be as constructed in the proof of \Cref{lem:boundinglemma} and write:
% \[
% m_f^{\vec c}(\vec x,\vec y) \ \dfn \ e_f(\sumlen{\vec x}) + d_f\sum\vec c + \max |\vec y|
% \]
% If there are constants $\vec c$ such that $\vec a$ satisfy \Cref{eq:oracle-const-max-bound},
% then:\footnote{To be clear, here we write $|\vec y_i| \leq m_f^{\vec c}(\vec x,\vec y)  $ here as an abbreviation for $\{|y_{ij} \leq m_f^{\vec c}(\vec x,\vec y)\}_j$.}
% \begin{equation}
%     \label{eq:input-bounding-eqn}
%     f(\vec a)(\vec x; \vec y) \ = \ f(\lambda \vec x_i.\lambda |\vec y_i| \leq m_f^{\vec c}(\vec x,\vec y) . a_i(\vec x_i;\vec y_i))_i(\vec x;\vec y)
% \end{equation}
% \end{prop}





% \toappendix{\begin{proof}
% We proceed by induction on the definition of $f(\vec a)(\vec x;\vec y)$, always making explicit the oracles of a function.
% % Let us temporarily write $m_f(\vec x,\vec y) \dfn e_f(\sumlen{\vec x}) + d_f\sum\vec c + \max |\vec y|$.

% The initial functions and oracle calls are immediate, due to the `$\max |\vec y|$' term in \Cref{eq:input-bounding-eqn}. 

% If $f(\vec a)(\vec x; \vec y) = h(\vec a)(\vec x; \vec y, g(\vec a)(\vec x; \vec y))$ then, by the inductive hypothesis for $h(\vec a)$, any oracle call from $h(\vec a)$ only take safe inputs of lengths:
%     \[
%     \begin{array}{rll}
%         \leq & e_h(\sumlen{\vec x}) + d_h \sum \vec c + \max (|\vec y|, |g(\vec a)(\vec x;\vec y)|) \\
%         \leq & e_h(\sumlen{\vec x})+ d_h \sum \vec c + e_g(\sumlen{\vec x}) + d_g\sum \vec c + \max |\vec y| & \text{by Bounding \Cref{lem:boundinglemma}} \\
%         \leq &  (e_h(\sumlen{\vec x} + e_g(\sumlen{\vec x})) + (d_h +d_g)\sum\vec c + \max |\vec y|\\
%          \leq & e_f(\sumlen{\vec x}) + d_f\sum\vec c + \max|\vec y|
%     \end{array}
%     \]
%     Note that any oracle call from $g(\vec a)$ will still only take safe inputs of lengths $\leq e_g(\sumlen{\vec x}) + d_g\sum \vec c + \max |\vec y|$, by the inductive hypothesis, and $e_g$ and $d_g$ are bounded above by $e_f$ and $d_f $ respectively.
    
% If $f(\vec a)(\vec x;\vec y) = h(\vec a)(\vec x, g(\emptyset)(\vec x;);\vec y)$ then, by the inductive hypothesis, any oracle call will only take safe inputs of lengths:
% \[
% \begin{array}{rll}
%     \leq & e_h(\sumlen{\vec x} + |g(\emptyset)(\vec x;)|) + d_h\sum\vec c + \max |\vec y| \\
%     \leq & e_h(\sumlen{\vec x} + e_g(\sumlen{\vec x})) + d_h\sum\vec c + \max |\vec y| & \text{by Bounding \Cref{lem:boundinglemma}} \\
%     \leq & e_f(\sumlen{\vec x}) + d_f\sum \vec c + \max |\vec y|
% \end{array}
% \]

% Finally suppose $f(\vec a)(\vec x;\vec y)= h(\vec a, \lambda \vec u\permpref \vec x, \lambda \vec v . f(\vec a)(\vec u;\vec v))(\vec x;\vec y)$. 
% We proceed by a sub-induction on $\sumlen{\vec x}$.
% Note that, since $\vec u \permpref \vec x \implies \sumlen{\vec u} <\sumlen{\vec x}$, we immediately inherit from the inductive hypothesis the appropriate bound on safe inputs for oracle calls from $\lambda \vec u\permpref \vec x, \lambda \vec v. f(\vec a)(\vec u;\vec v)$.

% Now, recall from the Bounding \Cref{lem:boundinglemma}, whenever $\vec u \permpref \vec x$ (and so $\sumlen{\vec u} < \sumlen{\vec x}$), we have $|f(\vec u;\vec v)| \leq e_f(\sumlen{\vec x}-1) + d_f\sum \vec c + \max |\vec v|$.
% So by setting $c = e_f(\sumlen{\vec x}-1)$ in the inductive hypothesis for $h(\vec a,a)$, with $a = \lambda \vec u\permpref \vec x, \lambda \vec v. f(\vec a)(\vec u;\vec v)$, any oracle call from $h(\vec a,a)$ will only take safe inputs of lengths:
% \[
% \begin{array}{rll}
%     \leq & e_h(\sumlen{\vec x}) + d_h e_f(\sumlen{\vec x}-1) + d_h\sum \vec c + \max |\vec y| \\
%     \leq & e_f(\sumlen{\vec x}) + d_f \sum \vec c + \max |\vec y| & \text{by \Cref{eq:bnd-lem-rec-fn-inv}}
% \end{array}
% \]
% % Now, for an oracle call from $\lambda \vec u\permpref \vec x, \lambda \vec v. f(\vec a)(\vec u;\vec v)$ we may simply apply the same reasoning again, essentially proceeding by induction on $\sumlen{\vec x}$.
% This completes the proof.
% \end{proof}
% }

\subsection{Soundness results}
In this subsection we  show that the function algebras $\bcpp$ and $\nbcpp$ (as well as $\nbc$)  capture precisely the classes  $\fptime$ and $\felementary$, respectively. We start with $\bcpp$:

\begin{thm}
\label{thm:fp-soundness}
Suppose $\vec a$ satisfies \Cref{eq:oracle-const-max-bound} for some constants $\vec c$.
We have the following:
\begin{enumerate}
    \item\label{item:soundness-fp} If $f(\vec x;\vec y) \in \bcpp (\vec a)$ then $f(\vec x,\vec y) \in \fptime(\vec a)$.
    % \item\label{item:soundness-fpspace} if $f(\vec x;\vec y) \in \sbcpp(\vec a)$ then $f (\vec x,\vec y) \in \fpspace(\vec a)$, as long as each $d_i = 1$.
    \item\label{item:soundness-felementary} If $f(\vec x;\vec y) \in \nbcpp (\vec a)$ then $f(\vec x,\vec y) \in \felementary(\vec a)$.
\end{enumerate}
\end{thm}


Note in particular that, for $f(\vec x;\vec y)$ in $\bcpp$ or $\nbcpp$, i.e.\ not using any oracles, we immediately obtain membership in $\fptime$ or $\felementary$, respectively.
However, the reliance on intermediate oracles during a function definition causes some difficulties that we must take into account.
At a high level, the idea is to use the Bounding Lemma   (namely \Cref{eq:input-bounding-eqn}) to replace certain oracle calls with explicit appropriately bounded functions computing their graphs. 
From here we compute $f(\vec x;\vec y)$ by a sort of `course-of-values' recursion on $\permpref$, storing previous values in a lookup table.
In the case of $\bcpp$, it is important that this table has polynomial-size, since there are only $m! \prod |\vec x|$ permutations of prefixes of a list $\vec x= x_1, \dots, x_m$ (which is a polynomial of degree $m$).

\begin{proof}[Proof of~\Cref{thm:fp-soundness}.]
We proceed by induction on the definition of $f(\vec x;\vec y)$.

Each initial function is polynomial-time computable, and each (relativised) complexity class considered is under composition, so it suffices to only consider the respective recursion schemes.
We shall focus first on the case of $\bcpp (\vec a)$, \Cref{item:soundness-fp} above, so that $e_f$ is a polynomial and $d_f=1$.



Suppose we have $h(a)(\vec x;\vec y) \in \bcpp(a, \vec a)$ and let:
\[
f(\vec x;\vec y) = h(\lambda \vec u \subset \vec x, \lambda \vec v \subseteq \vec y. f(\vec u; \vec v))(\vec x;\vec y)
\]
We start by making some observations:
\begin{enumerate}
    \item\label{item:poly-growth-rate} First, note that $|f(\vec x;\vec y)| \leq e_f(|\vec x|) +d_f\sum \vec c + \max |\vec y| $, by the Bounding \Cref{lem:boundinglemma}, and so $|f(\vec x;\vec y)|$ is polynomial in $|\vec x,\vec y|$.
%     \begin{itemize}
% \item $|f(\vec x;\vec y)| \leq e_f(|\vec x|) +d_f\sum \vec c + \max |\vec y| $ by the Bounding \Cref{lem:boundinglemma}.
%     \item $h(a)(\vec x;\vec y) \in \fptime(a,\vec a)$, by the inductive hypothesis, whenever $|a(\vec x;\vec y)| \leq c + d_f \sum \vec c + \max |\vec y|$, cf.~\Cref{eq:oracle-const-max-bound}.
%     In particular, by the above point, this holds when we set $a = \lambda \vec u \permpref \vec x, \lambda \vec v \permprefeq \vec y . f(\vec u; \vec v)$, by setting $c = e_f(\sumlen{\vec x}-1)$.
% \end{itemize}
\item\label{item:poly-many-pp} Second, note that the set $[\vec x;\vec y] \dfn \{(\vec u,\vec v) \ | \ \vec u \permpref \vec x , \vec v \permprefeq \vec y\} $ has size polynomial in $|\vec x,\vec y|$:
\begin{itemize}
    \item write $\vec x = x_1 , \dots, x_m$ and $\vec y = y_1, \dots, y_n$.
    \item Each $x_i$ and $y_j$ have only linearly many prefixes, and so there are at most $|x_1|\cdot \cdots \cdot |x_m||y_1| \cdot \cdots \cdot |y_n| \leq \sumlen{\vec x,\vec y}^{m+n}$ many choices of prefixes for all the arguments $\vec x, \vec y$. 
    (This is a polynomial since $m$ and $n$ are global constants).
    \item Additionally, there are $m!$ permutations of the arguments $\vec x$ and $n!$ permutations of the arguments $\vec y$.
    Again, since $m$ and $n$ are global constants, we indeed have
    $|[\vec x;\vec y]| = O(\sumlen{\vec x, \vec y}^{m+n})$, which is polynomial in $|\vec x,\vec y|$.
\end{itemize}
\end{enumerate}

We describe a polynomial-time algorithm for computing $f(\vec x;\vec y)$ (over oracles $\vec a$) by a sort of `course-of-values' recursion on the order $\permpref \times \permprefeq $ on $[\vec x;\vec y]$.

First, for convenience, temporarily extend $\permpref \times \permprefeq$ to a total well-order on $[\vec x;\vec y]$, and write $S$ for the associated successor function.
    Note that $S$ can be computed in polynomial-time from $[\vec x;\vec y]$.


Define $F(\vec x,\vec y) \dfn \langle f(\vec u;\vec v)\rangle_{\vec u \permpref \vec x, \vec v \permprefeq \vec y}$, i.e.\ it is the graph of $\lambda \vec u \permpref \vec x, \lambda \vec v \permprefeq \vec y . f(\vec u;\vec v)$ that we shall use as a `lookup table'.
    Note that $|F(\vec x,\vec y)| $ is polynomial in $|\vec x,\vec y|$ by \Cref{item:poly-growth-rate} and \Cref{item:poly-many-pp} above.
    Now, we can write:\footnote{Here, as abuse of notation, we are now simply identifying $F(\vec x;\vec y)$ with $\lambda \vec u \permpref \vec x, \lambda \vec v \permprefeq \vec y . f(\vec x;\vec y)$.}
    \[
    \arraycolsep=2pt
    \begin{array}{rcl}
        F(S(\vec x,\vec y)) & =&  \langle f(S(\vec x,\vec y)), F(\vec x,\vec y) \rangle  \\
        & = & \langle h(F(\vec x ,\vec y))(\vec x;\vec y),F(\vec x,\vec y)\rangle
    \end{array}
    \]
    Again by \Cref{item:poly-many-pp} (and since $F$ is polynomially bounded), this recursion terminates in polynomial-time.
    We may now simply calculate $f(\vec x;\vec y)$ as $h(F(\vec x,\vec y))(\vec x;\vec y)$.
    
The argument for $\nbcpp$ is similar, though we need not be as careful about computing the size of the lookup tables ($F$ above) for recursive calls.
The key idea is to use the Bounding Lemma (\Cref{eq:elem-const-max-bound}) to bound the safe inputs of recursive calls so that we can adequately store the lookup table for previous values.
\end{proof}


\subsection{Completeness and characterisations}
We are now ready to give our main function algebraic characterisation results for polynomial-time:

\begin{cor}\label{cor:bcpp-bc-fptime}
The following are equivalent:
\begin{enumerate}
    \item\label{item:complfp-bc} $f(\vec x;) \in \bc$.
    \item\label{item:complfp-bcpp} $f(\vec x;) \in \bcpp$.
    \item\label{item:complfp-fp} $f(\vec x) \in \fptime$.
\end{enumerate}
\end{cor}
\begin{proof}
% [Proof sketch]
$\eqref{item:complfp-bc}\implies \eqref{item:complfp-bcpp}$ is trivial, and $\eqref{item:complfp-bcpp} \implies \eqref{item:complfp-fp}$ is given by \Cref{thm:fp-soundness}.\eqref{item:soundness-fp}.
Finally, $\eqref{item:complfp-fp} \implies \eqref{item:complfp-bc}$ is from \cite{BellantoniCook}, stated in \Cref{thm:bellantoni} earlier.
\end{proof}

The remainder of this subsection is devoted to establishing a similar characterisation for $\nbc$, $ \nbcpp$ and $\felementary$.
% To this end, we naturally require the following result:



% \begin{proof}[Proof sketch] Using a formulation of $\felementary$ in binary notation, and we prove by induction on the definition of $f(\vec{x}) \in \felementary$  that there exists  $f^*(x; \vec{x})\in \nbc$ and a monotone function $t_f\in \felementary$ such that for all integers $\vec{x}$ and all $w \geq t_f(\vec{x})$ we have $f^*(w; \vec{x})=f(\vec{x})$. Then, we use the function $\ex(x; y)\in \nbc$ in~\eqref{eq:safe-exp-by-snrec} to achieve arbitrary elementary growth rate, which implies  $f(\vec{x}) \in \nbc$. The proof technique based on the construction of the pair $(f^*(x; \vec x),t_f )$ is well-known since~\cite{BellantoniCook}, and has been adapted to the case of the elementary functions  in~\cite{Wirz99characterizingthe}.
% \end{proof}



% \toappendix{The proof can be divided into three steps:
% \begin{itemize}
%     \item We define an equivalent formulation of the class $\felementary$ for binary notation, having 
%     an exponential growth rate function $\varepsilon^1(x)$, whence one may define arbitrary elementary growth rate by setting $\varepsilon^{m+1}(x) \dfn  \varepsilon^1(\varepsilon^{m}(x))$.
%     % 
% %     the following exponential growth rate function:
% %   \begin{equation*}
% %     \begin{aligned}
% %              \varepsilon(0,y)&\dfn \succ 0 (y)\\
% %              \varepsilon(\succ ix,y)&\dfn \varepsilon(x,  \varepsilon(x,y))
% %     \end{aligned}
% %     \end{equation*}
% %     and we show that the unary and the binary presentation give raise to the same class of functions; 
%     Now we show  that any function $f(\vec{x})\in \felementary$   in binary notation is bounded above by $ \varepsilon^m(\max(\vec{x}))$ for some $m$.
%     % , where $ \varepsilon^1(x)  \dfn  \varepsilon(x,1) $ and $  \varepsilon^{m+1}(x) \dfn  \varepsilon^1(\varepsilon^{m}(x))  $
%     \item  We show that for any function  $f(\vec{x}) \in \felementary$ in binary notation there are a function $f^*(x; \vec{x})\in \nbc$ and a monotone function\footnote{{Bear in mind that the function $e_f$ here is not necessarily the same as that of \Cref{lem:boundinglemma} in the previous section.}} $e_f\in \felementary$ such that for all integers $\vec{x}$ and all $w \geq e_f(\vec{x})$ we have $f^*(w; \vec{x})=f(\vec{x})$. 
%     \item Now, given the function $\ex(x; y)$ in~\eqref{eq:safe-exp-by-snrec}, we construct the function $\ex^m(x;)$ by induction on $m$:
% \begin{equation*}
%     \begin{aligned}
%      \ex^1(x;)&= \ex(x;1)\\
%       \ex^{m+1}(x;)&=   \ex^{m}(  \ex^{1}(x;);) 
%     \end{aligned}
% \end{equation*}
% so that, by increasing $m$, $\ex^m(x;)$ achieves arbitrary elementary growth rate, just like $\varepsilon^m$. If  $f(\vec x) \in \felementary$, by the previous point we obtain a function $f^*(w; \vec x)\in\nbc$ and a monotone function $e_{f}\in \felementary$. By the first point,  there exists $m \geq 1$ such that $  e_f(\vec{x})\leq  \ex^{m}(\max_{\sharp \vec{x}}(\vec{x};);) $, where $\max(\vec{x};)$ in $\nbc$. Therefore, $ f(\vec{x};)=f^*(\ex^{m}(\max_{\sharp \vec{x}}(\vec{x};););\vec{x})\in \nbc$.
% \end{itemize}
% In particular, the first point is folklore, while the proof technique adopted for proving the last two points is well-known since~\cite{BellantoniCook}, and has been adapted to the case of the elementary functions   by~\cite{Wirz99characterizingthe}.}


% \red{
% \begin{rem}\label{rem:restriction-scheme-oracles}
% It is important to stress that the above construction applies  the  scheme  in~\ref{defn:oracles-functions}.\ref{enum:safe-composition-condition} only once, i.e.~when we compose the  function $\ex^m\in \nbc(\emptyset)$ along the normal parameter of the  function $f^*\in \nbc(\emptyset)$. Thanks to this observation, by combining the above theorem and~\Cref{thm:fp-soundness}.\ref{item:soundness-felementary},  we can conclude that the  safe composition scheme in~\ref{defn:oracles-functions}.\ref{enum:safe-composition-condition} can be restricted w.l.o.g.~to oracle-free functions. Moreover, as we are mainly interested in $\nbc$, so that   functions over oracles are essentially introduced to prove properties about the latter, we shall extend our restriction to the function algebras $\nbc(\vec a)$.
% \end{rem}
% }


% \begin{proof}
% Point~\ref{enum:elementary1} implies point~\ref{enum:elementary2} trivially. Point~\ref{enum:elementary2} implies point~\ref{enum:elementary3} by Lemma~\ref{lem:boundinglemma} and because \todo{Here explain why from bounding lemma we conclude elementary time soundness}. Last, point~\ref{enum:elementary3}  implies point~\ref{enum:elementary1} by Theorem~\ref{thm:elementary-in-nbc}. 
% \end{proof}


 To begin with, we recall the definition of the class $\felementary$: 

\begin{defn} \label{defn:felementary}  $\felementary$ is the smallest set of functions containing:
\begin{itemize}
    \item $0()\dfn 0 \in \Nat$,  
    \item $\pi^n_i(x_1, \ldots, x_{n})\dfn x_j$, whenever $1 \leq j \leq n$;
    \item $\mathsf{s}(x)\dfn x+1$;
    \item the function $E_2$ defined as follows:
    \begin{equation*}
    \begin{aligned}
            E_1(x)&=x^2+2\\
             E_2(0)&=2\\
             E_2(x+1)&=E_1(E_2(x))
    \end{aligned}
    \end{equation*}
    \end{itemize}
    and closed under the following:
    \begin{itemize}
    \item (Composition) If $f(\vec x, x), g(\vec x)\in \felementary$ then so is $f(\vec x, g(\vec x))$; 
    \item (Bounded recursion) If $g(\vec x), h (x,\vec x,y),  j(x, \vec x)$ are functions in $ \felementary$ then so is $f(x, \vec x)$ given by:
    \begin{equation*}
    \arraycolsep=2pt
        \begin{array}{rcl}
              f(0, \vec{x})&\dfn& g(\vec{x})\\
              f(x+1, \vec{x})&\dfn& h(x, \vec{x}, f(x, \vec{x}))
        \end{array}
    \end{equation*}
    provided that $f(x, \vec{x})\leq j(x, \vec{x})$.
\end{itemize}
\end{defn}


\begin{prop}[\cite{rose1984subrecursion}]\label{prop:rose} Let $f\in \felementary$ be a $k$-ary function. Then, there exists an integer $m$ such that:
\begin{equation*}
    f(\vec{x})\leq E^m_{2}(\max_k(\vec{x}))
\end{equation*}
where $E^0_{2}(x)=x$ and $E^{m+1}_{2}(x)=E_2(E^{m}_{2}(x))$.
\end{prop}

For our purposes we shall consider a  formulation of this class in binary notation, that we call  $\felementary_{0,1}$. 

\begin{defn} \label{defn:E3}  $\felementary_{0,1}$ is the smallest set of functions containing:
\begin{itemize}
    \item $0()\dfn 0 \in \Nat$,  
    \item $\pi^n_i(x_1, \ldots, x_{n})\dfn x_j$, whenever $1 \leq j \leq n$;
    \item $\succ i(x)\dfn 2x+i$, for $i \in \{ 0,1\}$
    \item the function $\varepsilon(x,y)$ defined as follows:
    \begin{equation*}
    \arraycolsep=2pt
    \begin{array}{rcl}
             \varepsilon(0,y)&\dfn& \succ 0 (y)\\
             \varepsilon(\succ ix,y)&\dfn& \varepsilon(x,  \varepsilon(x,y))
    \end{array}
    \end{equation*}
    \end{itemize}
    and closed under the following:
    \begin{itemize}
    \item (Composition) If $f(\vec x, x), g(\vec x)\in \felementary_{0,1}$ then so is $f(\vec x, g(\vec x))$; 
    \item (Bounded recursion on notation) If $g(\vec x), h_i (x,\vec x,y)$, $j(x, \vec x)$ are functions in  $\felementary_{0,1}$  then so is $f(x, \vec x)$ given by:
    \begin{equation*}
    \arraycolsep=2pt
        \begin{array}{rcl}
              f(0, \vec{x})&\dfn &g(\vec{x})\\
              f(\succ ix, \vec{x})&\dfn & h_i(x, \vec{x}, f(x, \vec{x}))
        \end{array}
    \end{equation*}
    provided that $f(x, \vec{x})\leq j(x, \vec{x})$.
\end{itemize}
\end{defn}

It is easy to show that the unary and the binary definition of the class of elementary time computable functions coincide. To see this, we first define $\varepsilon^n(x)$ as
\begin{equation}\label{eqn:vaerpsilon}
\arraycolsep=2pt
    \begin{array}{rcl}
   \varepsilon^1(x)  &\dfn&  \varepsilon(x,1)  \\
   \varepsilon^{m+1}(x)& \dfn & \varepsilon^1(\varepsilon^{m}(x))   
\end{array}
\end{equation}
which allows us to prove that $\varepsilon^n(x)$ plays the role of rate growth function as the function $E_2^n(x)$ (see Proposition~\ref{prop:rose}). 

\gianluca{Stress the equivalence between  binary notation and unary notation for elementary time functions is folklore?}

\begin{prop}\label{prop:EsubseteqEbin}{\ }
\begin{enumerate}
    \item \label{enum:EsubseteqEbin1} $\felementary = \felementary_{0,1}$;
    \item \label{enum:EsubseteqEbin2} for any $f \in \felementary_{0,1}$ $k$-ary function there is an integer $m$ such that:
\begin{equation*}
    f(\vec{x})\leq \varepsilon^m(\max_k(\vec{x}))
\end{equation*}
\end{enumerate}
\end{prop}
\begin{proof}
 Let us first prove point~\ref{enum:EsubseteqEbin1}. For the $\supseteq$ direction we show that for any $f  \in \felementary_{0,1}$ there exists 
 $n$ such that:
\begin{equation}\label{eqn:feleementary01-bound}
    \s{f(\vec x)}\leq 2_n(\sum\s{\vec x}) 
\end{equation}
where $2_0(x)=x$ and  $2_{n+1}(x)=2^{2_n(x)}$. 
 Since $\s{x}=\lceil \log_2(x+1)\rceil$, from the above inequation we would have that, for some $m$:
\begin{equation*}
    f(\vec x)\leq 2_{n+m}(\sum \vec x)
\end{equation*}
 which allows us to conclude $f  \in \felementary$, as  the elementary time computable functions are exactly the elementary space  ones. The inequation~\eqref{eqn:feleementary01-bound} can be proved by induction on $f$,  noticing that $\s{\varepsilon(x, y)}= 2^{\s{x}}+ \s{y}$. Concerning the $\subseteq$ direction, we prove  by induction on  $f\in \felementary$ that there exists a function $\hat{f}\in \felementary_{0,1}$ such that, for all $\vec x= x_1, \ldots, x_n$:
\begin{equation*}
   f(\vec x)= \s{\hat{f}(\num{x_1}, \ldots, \num{x_n})}
\end{equation*}
where $\num{m}\dfn \succ 1 ^m(0)=\succ 1(\overset{m}{\ldots}  \succ 1(0))$.  Since the functions $\s{\cdot}$ and $x \mapsto \num{x}$ are both in $\felementary_{0,1}$, we are able to conclude $f \in \felementary_{0,1}$. The case $f=0$ is trivial. As for the cases $f=\mathsf{s}$ and $f= \pi^n_i$, we first notice that $\s{\num{x}}=x$. Then, we have: 
\begin{equation*}
\def\arraystretch{1.5}
\begin{array}{c}
\mathsf{s}(x)=\s{\num{s(x)}}
\\
  \pi^n_i(\vec x)=x_i=\s{\num{x_i}}=  \s{\pi^n_i(\num{x_1}, \ldots, \num{x_n})}
\end{array}
\end{equation*}
Concerning the case of $E_2(x)$, we first notice that the following property holds for any $m$ and some $k$:
 \begin{equation}\label{eqn:boundgrowthrate}
     E^m_2(x)\leq \varepsilon^{m+k}(x)
 \end{equation}
where $\varepsilon^{n}(x)$ is as in~\eqref{eqn:vaerpsilon}. Moreover, the  function $E_2(x)$ can be defined by two applications of bounded recursion proceeding from the successor and the projection functions, where each recursion can be bounded by $E_2(x)$, and hence by $\varepsilon^{k}(x)$ for some $k$. This means that the case of  $E_2(x)$ can be reduced to the case of bounded recursion. Suppose now that $f(\vec x)=h(\vec x, g(\vec x))$.  We define  $\hat{f}(\vec x)= \hat{h}(\vec x, \hat{g}(\vec x))$ so that, by induction hypothesis:
\begin{equation*}
\arraycolsep=2pt
    \begin{array}{rcl}
     f(\vec x)&=& h(\vec x, g(\vec x))\\
     &=&\s{\hat{h}(\num{x_1}, \ldots, \num{x_n}, \num{g(\vec x)})}\\
     &=&\s{\hat{h}(\num{x_1}, \ldots, \num{x_n}, \num{\s{\hat{g}(\num{x_1}, \ldots, \num{x_n})}})}\\
      &=&\s{\hat{h}(\num{x_1}, \ldots, \num{x_n}, \hat{g}(\num{x_1}, \ldots, \num{x_n}))}\\
      &=& \s{\hat{f}(\vec x)}= \s{\hat{h}(\vec x, \hat{g}(\vec x))}
    \end{array}
\end{equation*}
Last, suppose that $f$ has been obtained by  bounded recursion from $h, g, j$, i.e.
    \begin{equation*}
    \arraycolsep=2pt
        \begin{array}{rcl}
              f(0, \vec{x})&=&g(\vec{x})\\
              f(y+1, \vec{x})&=& h(y, \vec{x}, f(y, \vec{x}))
        \end{array}
    \end{equation*}
    provided that $f(y, \vec{x})\leq j(y, \vec{x})$. We define $\hat{f}$ as follows:
    \begin{equation*}
    \arraycolsep=2pt
        \begin{array}{rcl}
              \hat{f}(0, \vec{x})&=&\hat{g}(\vec{x})\\
              \hat{f}(\succ iy, \vec{x})&=& \hat{h}(y, \vec{x}, \hat{f}(y, \vec{x}))
        \end{array}
    \end{equation*}
We show by induction on $y$ that:     
    \begin{equation*}
   f(y, \vec x)= \s{\hat{f}(\num{y}, \num{x_1}, \ldots, \num{x_n})}
\end{equation*}
We have:
\begin{equation*}
\arraycolsep=2pt
    \begin{array}{rcl}
  f(0, \vec x)&=& g(\vec x)\\
  &=&\s{\hat{g}(\num{x_1}, \ldots, \num{x_n})}\\
  &=&
    \s{\hat{f}(\num{0}, \num{x_1}, \ldots, \num{x_n})}\\ \\
    f(y+1, \vec x)&=& h(y,\vec x, f(y, \vec x))\\
    &=& \s{\hat{h}(\num{y}, \num{x_1}, \ldots, \num{x_n}, \num{f(y, \vec x)})}\\
    &=& \s{\hat{h}(\num{y}, \num{x_1}, \ldots, \num{x_n}, \num{\s{\hat{f}(\num{y}, \num{x_1}, \ldots, \num{x_n})}})}\\
     &=& \s{\hat{h}(\num{y}, \num{x_1}, \ldots, \num{x_n}, \hat{f}(\num{y}, \num{x_1}, \ldots, \num{x_n})}\\
    &=& \s{\hat{f}(\num{y+1}, \num{x_1}, \ldots, \num{x_n})}
   \end{array}
\end{equation*}
since    $f(y, \vec{x})\leq j(y, \vec{x})$ and $ j(y, \vec x)= \s{\hat{j}(\num{y}, \num{x_1}, \ldots, \num{x_n})}$  by induction hypothesis, we are done.     
    
    
    
     Point~\ref{enum:EsubseteqEbin2} follows by point~\ref{enum:EsubseteqEbin1},  Proposition~\ref{prop:rose} and~\eqref{eqn:boundgrowthrate}.
\end{proof}



Completeness for $\nbc$ is based on a standard technique (see~\cite{BellantoniCook}),  adapted to the case of $\felementary$  in~\cite{Wirz99characterizingthe}.

\begin{lem}\label{lem:padding}
For any $f(\vec{x}) \in \felementary_{0,1}$  there are a function $f^*(x; \vec{x})\in \nbc$ and a monotone function $t_f\in \felementary_{0,1}$ such that for all integers $\vec{x}$ and all $w \geq t_f(\vec{x})$ we have $f^*(w; \vec{x})=f(\vec{x})$.
\end{lem}
\begin{proof}
The proof is by induction on the definition of $f$.  If $f$ is the zero, successor or projection  function then $f^*\in \nbc$.  In this case we choose $t_f=0$. The function $\varepsilon$ has a definition by one application of bounded recursion on notation, proceeding from the successors and the projection functions, where each recursion is bounded by $\varepsilon$. Since the treatment of bounded recursion  does not make use of the induction hypothesis for the bounding function, we can use this method to get functions $\varepsilon^*\in \nbc$ and $t_{\varepsilon}\in \felementary_{0,1}$ with the required properties.  If $f(\vec x)= h(\vec x, g(\vec{x}))$ then we set  $f^*(w;\vec{x})=h^*(w; \vec x, g^*(w;\vec{x}))$, which is in $\nbc$.  Since the function $g^*$ is clearly bounded by a monotone function $b \in \felementary_{0,1}$, we set $ t_f(\vec{x})= t_h(\vec x, b(\vec{x}))+  t_{g}(\vec{x})$,  which is monotone. By applying the induction hypothesis,  if   $w \geq t_f(\vec{x})$ then:
\begin{equation*}
         f^*(w; \vec{x})=  h^*(w; \vec{x}, g^*(w; \vec{x})) 
         =h^*(w;\vec x, g(\vec{x})) 
         = h(\vec x, g(\vec{x}))
\end{equation*}
Let us finally suppose that $f(x, \vec{y})$ is defined by bounded recursion on notation from $g(\vec y)$, $h_i(x, \vec y, f(x, \vec y))$ and $j(x, \vec y)$. By applying the induction hypothesis we set:
\begin{equation*}
\arraycolsep=2pt
    \begin{array}{rcl}
     \hat{f}(0,w; x, \vec{y})&=& g^*(w; \vec{y})\\
      \hat{f}(\succ i(x),w; x, \vec{y})&=& \cnd(; W(\succ i(v),w;x), g^*(w; \vec{y}),  h^*_i(w; W(v,w;x), \vec{y}, f^*(v,w;x,\vec{y})) )\\
     f^*(w;x, \vec{y})&= & \hat{f}(w,w; x, \vec{y}) 
    \end{array}
\end{equation*}
where $ W(v,w;x)=\dot-( \dot-(v;w) ; x)$ and $\dot-(x;y)$ is the truncated subtraction, which is in $\bc$, and hence in  $\nbc$. We can easily show that $f^*(w;x, \vec y)\in \nbc$. We define $t_f(x, \vec{y})=t_{g}(\vec{y})+\sum_i t_{h_i}(x, \vec{y}, j(x, \vec{y}))$, where $j$ is the bounding function. Assuming $j$ to be monotone, $t_f$ is monotone too.  We now show by induction on $u$ that, whenever $w \geq t_f(x,\vec{y})$ and $w-x\leq u \leq w$:
\begin{equation}
\label{eqn:toshowpadding}
    \hat{f}(u,w;x,\vec{y})=f(x-(w-u), \vec{y})
\end{equation}
If $u=w-x$ then we have two cases:
\begin{itemize}
    \item if $u=0$ then $\hat{f}(0, w; x, \vec{y})=g^*(w; \vec{y})$;
    \item if $u=\succ i(v)$ then, since $W(\succ i(v),w;x)=0$, we have $\hat{f}(\succ i(v), w; x, \vec{y})=g^*(w; \vec{y})$.
\end{itemize}
Hence, in any case:
\begin{equation*}
\hat{f}(u, w; x, \vec{y})=g^*(w; \vec{y})=g(\vec y)=f(0; \vec{y})= f(x-(w-u), \vec{y})    
\end{equation*}
Let us now suppose that $w-x<u\leq w$. This means that $u=\succ i(v)$ and $W(\succ i(v),w;x)>0$. Moreover,  by monotonicity of $t_f$ and definition of $j$:
\begin{equation*}
\arraycolsep=2pt
\begin{array}{rcl}
     w &\geq&     t_f(x, \vec{y})\\ 
     &\geq& t_f(x-(w-\succ i(v)), \vec{y}) \\
     &\geq& t_f(x-(w-v), \vec{y}) \\
     &\geq& t_{h_i}(x-(w-v), \vec{y}, j(x-(w-v), \vec{y}))\\
      &\geq& t_{h_i}(x-(w-v), \vec{y}, f(x-(w-v), \vec{y}))
\end{array}
\end{equation*}
By applying the induction hypothesis:
\begin{equation*}
\arraycolsep=2pt
\begin{array}{rcl}
      \hat{f}(\succ i(v), w; x, \vec{y})&=&
        h^*_i(w; W(v, w;x), \vec{y}, \hat{f}(v,w;x, \vec{y}))\\
         &=&  h^*_i(w; W(v, w;x), \vec{y}, f(x-(w-v), \vec{y}))\\
         &=& h^*_i(w; x-(w-v), \vec{y}, f(x-(w-v), \vec{y}))\\
         &=& h_i(x-(w-v), \vec{y}, f(x-(w-v), \vec{y}))\\
           &=&f(\succ i((x-(w-v))), \vec{y})\\
           &=& f((x-(w-\succ i(v))), \vec{y})
\end{array}
\end{equation*}
Now, by~\eqref{eqn:toshowpadding}, for all $w \geq t_f(x, \vec{y})$ we have:
\begin{equation*}
  f^*(w; x, \vec{y})=  \hat{f}(w,w;x, \vec{y})= f(x, \vec{y})
\end{equation*}
and this concludes the proof.
\end{proof}

\begin{thm}
\label{thm:elementary-in-nbc}
If $f(\vec x) \in \felementary$ then $f(\vec x;) \in \nbc$.
\end{thm}
\begin{proof}
First, given the function $\ex(x; y)$ in~\eqref{eq:safe-exp-by-snrec}, we construct the function $\ex^m(x;)$ by induction on $m$:
\begin{equation*}
\arraycolsep=2pt
    \begin{array}{rcl}
     \ex^1(x;)&=& \ex(x;1)\\
       \ex^{m+1}(x;)&=&   \ex^{m}(  \ex^{1}(x;);)
    \end{array}
\end{equation*}
Hence  $\varepsilon^m(x)=\ex^m(x;)$, for all $m\geq 1$. Now, let $f(\vec x) \in \felementary$. By Proposition~\ref{prop:EsubseteqEbin}.\ref{enum:EsubseteqEbin1}, we have that $f(\vec x) \in \felementary_{0,1}$. By Lemma~\ref{lem:padding},   there exist $f^*(w; \vec x)\in\nbc$ and a monotone function $t_{f}\in \felementary_{0,1}$  such that, for all $w, \vec{x}$ with  $w \geq t_f(\vec{x})$, it holds that $f^*(w; \vec{x})= f(\vec{x})$. By Proposition~\ref{prop:EsubseteqEbin}.\ref{enum:EsubseteqEbin2}  there exists $m \geq 1$ such that:
\begin{equation*}
    t_f(\vec{x})\leq  \ex^{m}(\max_{\sharp \vec{x}}(\vec{x};);) 
\end{equation*}
where $\max_{\sharp \vec{x}}(\vec{x};)$ is the $k$-ary maximum function, which is in $\bc$ by Theorem~\ref{thm:bellantoni}, and hence in $\nbc$. Therefore:
\begin{equation*}
    f(\vec{x};)=f^*(\ex^{m}(\max_{\sharp \vec{x}}(\vec{x};););\vec{x})\in \nbc
\end{equation*}
\end{proof}

Now, by the same argument as for \Cref{cor:bcpp-bc-fptime}, only using \Cref{lem:padding} above instead of appealing to \cite{BellantoniCook}, we can give our main characterisation result for algebras for elementary computation:

\begin{cor}\label{cor:nb-elementary-characterization} The following are equivalent:
\begin{enumerate}
    \item\label{enum:elementary1} $f(\vec x;)\in \nbc$.
    \item \label{enum:elementary2}  $f(\vec x;)\in \nbcpp$. 
    \item \label{enum:elementary3}  $f(\vec x)\in \felementary$.
\end{enumerate}
\end{cor}  \section{Characterisations for circular systems}
\label{sec:completeness}
We now return our attention to the circular systems $\cbc$ and $\ncbc$ that we introduced in \Cref{sec:two-tiered-circular-systems-on-notation}. 
We will address the complexity of their definable functions by `sandwiching' them between function algebras of \Cref{sec:some-variants}, given their characterisations that we have just established.

% In particular, in this section we show that $\cbc$ is `complete' for $\fptime$  and that $\ncbc$ is complete for $\felementary$. 


\subsection{Completeness}
To show that $\cbc$ contains all polynomial-time functions, we may simply simulate Bellantoni and Cook's algebra:



\begin{thm}
\label{thm:bc-in-cbc}
If $f(\vec x;\vec y) \in \bc$ then $f(\vec x; \vec y)\in \cbc$.
\end{thm}
\begin{proof}
By Proposition~\ref{prop:bc-type-system-characterisation} it suffices to show that for any $\bc$-derivation $\der$  there is a $\cbc$-coderivation $\der^*$  such that $\denot\der(\vec x; \vec y)=\denot{\der^*}(\vec x,  \vec y)$. The proof is by induction on $\der$. The only non-trivial case is when   $\der$ is the following derivation:
\[
\small
 \vlderivation{
    \vliiin{\srec}{}{\sn, \Gamma \seqar \n}{
        \vltr{\der_0}{\Gamma \seqar \n}{\vlhy{\ }}{\vlhy{\ }}{\vlhy{\ }}
    }{
        \vltr{\der_1}{\sn, \Gamma, \n \seqar \n}{\vlhy{\ }}{\vlhy{\ }}{\vlhy{\ }}
    }{
        \vltr{\der_2}{\sn, \Gamma, \n \seqar \n}{\vlhy{\ }}{\vlhy{\ }}{\vlhy{\ }}
    }
    }
\]
We define $\der^*$ as follows:
\[
\small
\vlderivation{
\vliin{\cnd_\sq}{\bullet}{\blue{\underline{\sn}},\Gamma \seqar \n}{
\vltr{\der_0^*}{\Gamma \seqar \n}{\vlhy{\ }}{\vlhy{\ }}{\vlhy{\ }}
}{
\vliin{\cut_{\n}}{{\scriptstyle i=1,2}}{\blue{\sn}, \Gamma  \seqar \n}{
   \vlin{\cnd_{\sq}}{\bullet}{
   \blue{\sn}, \Gamma  \seqar \red{\n}
   }{
   \vlhy{\vdots}
   }
}{
   \vltr{\der^*_i}{\blue{\sn}, \Gamma , \red{\n} \seqar \n}{\vlhy{\ }}{\vlhy{\ }}{\vlhy{\ }}
}
}
% {
%  \vliin{\cut_{\n}}{}{\blue{\sn}, \sq \vec \n, \vec \n \seqar \n}{
%   \vlin{\cnd_{\sq}}{\bullet}{
%   \blue{\sn}, \sq \vec \n, \vec \n \seqar \red{\n}
%   }{
%   \vlhy{\vdots}
%   }
% }{
%   \vltr{\hfunction_1}{\blue{\sn}, \sq \vec \n, \vec \n, \red{\n} \seqar \n}{\vlhy{\ }}{\vlhy{\ }}{\vlhy{\ }}
% }
% }
}
\]
where $\Gamma= \sq \vec \n, \vec \n$ and we identify the coderivations corresponding to the second and the third premise of the conditional rule, as they only differ on the sub-coderivation $\der^*_{i}$ ($i=1$ for the former and $i=2$ for the latter). 

The above coderivation is clearly safe and left-leaning, by the inductive hypotheses for $\der_0,\der_1,\der_2$. 
{To see that it is progressing, note that any infinite branch is either eventually entirely in $\der_0^*$, $\der_1^*$ or $\der_2^*$, in which case it is progressing by the inductive hypotheses, or it simply loops on $\bullet$ forever, in which case there is a progressing thread along the blue $\blue{\sn}$.} 

Moreover, the equational program associated with $\der^*$  is equivalent to:
 \[
 \arraycolsep=2pt
    \begin{array}{rcll}
         \denot{\der^*_\epsilon} (0, \vec x;\vec y) & \dfn &  \denot{\der^*_0} (\vec x;\vec y) \\
         \denot{\der^*_\epsilon} (\succ 0 x, \vec x; \vec y) & \dfn & \denot{\der^*_1} (x, \vec x; \vec y, \denot{\der^*_\epsilon}(x, \vec x; \vec y)) & \text{if $x \neq 0$} \\
         \denot{\der^*_\epsilon} (\succ 1 x , \vec x; \vec y) & \dfn & \denot{\der^*_2} (x, \vec x; \vec y,  \denot{\der_\epsilon}(x, \vec x; \vec y)) 
    \end{array}
    \]
so that $\denot\der(\vec x; \vec y)=\denot{\der^*}(\vec x,  \vec y)$.
\end{proof}

% \begin{proof}[Proof sketch] The proof is by induction on the definition of $f$. 
% The only interesting case is when $f(x, \vec x; \vec y)$ is defined by safe recursion on notation from the functions $g(\vec x; \vec y)$ and $h_i(x, \vec x; \vec y, y)$. Given $\der_g$ and $\der_{h_i}$ defining $g$ and $h_i$ by induction hypothesis, we may define $f(x,\vec x;\vec y)$ by the coderivation,
% \[
% \vlderivation{
% \vliin{\cnd_\sq}{\bullet}{\blue{\underline{\sn}},\Gamma \seqar \n}{
% \vltr{\der_g}{\Gamma  \seqar \n}{\vlhy{\ }}{\vlhy{\  \ \ }}{\vlhy{\ }}
% }{
% \vliin{\cut_{\n}}{{\scriptstyle i=0,1}}{\blue{\sn}, \Gamma \seqar \n}{
%    \vlin{\cnd_{\sq}}{\bullet}{
%    \blue{\sn}, \Gamma  \seqar \red{\n}
%    }{
%    \vlhy{\vdots}
%    }
% }{
%    \vltr{\der_{h_i}}{\blue{\sn}, \Gamma , \red{\n} \seqar \n}{\vlhy{\ }}{\vlhy{\ \ \ \  }}{\vlhy{\ }}
% }
% }
% % {
% %  \vliin{\cut_{\n}}{}{\blue{\sn}, \sq \vec \n, \vec \n \seqar \n}{
% %   \vlin{\cnd_{\sq}}{\bullet}{
% %   \blue{\sn}, \sq \vec \n, \vec \n \seqar \red{\n}
% %   }{
% %   \vlhy{\vdots}
% %   }
% % }{
% %   \vltr{\hfunction_1}{\blue{\sn}, \sq \vec \n, \vec \n, \red{\n} \seqar \n}{\vlhy{\ }}{\vlhy{\ }}{\vlhy{\ }}
% % }
% % }
% }
% \]
% where $\Gamma=  \sq \vec \n, \vec \n$.
% The only infinite branch (outside $\der_g$ and $\der_{h_i}$) loops on $\bullet$ and has a progressing thread in blue.
% \end{proof}




% \begin{proof}
% We proceed by induction on the definition of $f(\vec x;\vec y)$. We shall use the context $\Gamma$ to denote modal formulas, and the context $\Delta$ to denote non-modal formulas.   The cases where $f$ is a  Constant or a Successor function are straightforward. If $f$ is the Predecessor function then $\pder f$ is as follows:
% \begin{prooftree}
% \AxiomC{}
% \LeftLabel{$\zero$}
% \UnaryInfC{$\seqar \n$}
% \AxiomC{}
% \LeftLabel{$\id$}
% \UnaryInfC{$ \n \seqar \n$}
% \AxiomC{}
% \LeftLabel{$\id$}
% \UnaryInfC{$ \n \seqar \n$}
% \LeftLabel{$\cnd$}
% \TrinaryInfC{$\underline{\n} \seqar \n$}
% \end{prooftree}
% If $f$ is the Conditional function then $\pder f$ is the following, where $A,B,C$ are the non-modal formula $\n$, then $\pder f$ is as follows:
% \begin{prooftree}
% \AxiomC{}
% \LeftLabel{$\id$}
% \UnaryInfC{$\n \seqar \n$}
% \LeftLabel{$\wk$}
% \UnaryInfC{$\n,\underline{\n} \seqar \n$}
% \AxiomC{}
% \LeftLabel{$\id$}
% \UnaryInfC{$\n \seqar \n$}
% \LeftLabel{$\wk$}
% \UnaryInfC{$ \n,\underline{\n}\seqar \n$}
% \LeftLabel{$\wk$}
% \UnaryInfC{$\underline{\n},\n, \n \seqar \n$}
% \AxiomC{}
% \LeftLabel{$\id$}
% \UnaryInfC{$\n \seqar \n$}
% \LeftLabel{$\wk$}
% \UnaryInfC{$\underline{\n},\n \seqar \n$}
% \LeftLabel{$\wk$}
% \UnaryInfC{$\underline\n,\n, \n \seqar \n$}
% \LeftLabel{$\cnd$}
% \TrinaryInfC{$\underline{\n},\n, \n \seqar \n$}
% \end{prooftree}
% If $f$ is the Projection function, then $\pder f$ is as follows:
% \begin{prooftree}
% \AxiomC{}
% \LeftLabel{$\id$}
% \UnaryInfC{$\n \seqar \n$}
% \LeftLabel{$\wk$}
% \doubleLine
% \UnaryInfC{$\Gamma, \Delta  \seqar \n$}
% \end{prooftree}
% If  $f(\vec x; \vec y)= h(\vec x, g(\vec x;); \vec y)$ then $\pder f$ is as follows:
% \begin{prooftree}
%   \AxiomC{$\pder {g}$}
%   \noLine
%   \UnaryInfC{$\Gamma \seqar\n$}
%   \LeftLabel{$\sqr$}
%   \UnaryInfC{$\Gamma \seqar \sn$}
%   \doubleLine
%   \LeftLabel{$\wk$}
%   \UnaryInfC{$\Gamma, \Delta \seqar \sn$}
%   \AxiomC{$\pder h$}
%   \noLine
%   \UnaryInfC{$\Gamma, \sn, \Delta \seqar \n$}
%   \LeftLabel{$\cut$}
%   \BinaryInfC{$\Gamma, \Delta \seqar \n$}
%   \end{prooftree}
%  If $f(\vec x; \vec y)= h(\vec x; \vec y, g(\vec x;\vec y))$ then $\pder f$ is as follows:
% \begin{prooftree}
%   \AxiomC{$\pder {g}$}
%   \noLine
%   \UnaryInfC{$\Gamma, \Delta \seqar \n$}
%   \AxiomC{$\pder h$}
%   \noLine
%   \UnaryInfC{$\Gamma,  \Delta, \n \seqar \n$}
%   \LeftLabel{$\cut$}
%   \BinaryInfC{$\Gamma, \Delta \seqar \n$}
%   \end{prooftree}
%  Finally, suppose $f(x, \vec x; \vec y)$ is defined  by Safe recursion from  $g(\vec x; \vec y)$ and $h_i(x, \vec x; \vec y, z)$ ($i=0,1$). If $\Sigma= \Gamma, \Delta$  then $\pder f$ is as follows:
%   \begin{prooftree}
%   \def\defaultHypSeparation{\hskip .2cm}
%   \def\ScoreOverhang{1pt} 
%   \AxiomC{$\pder g$}
%   \noLine
%   \UnaryInfC{$\Sigma \seqar \n$}
%  \AxiomC{}
%   \LeftLabel{$\bullet$}
%   \UnaryInfC{$\sn,\Sigma \seqar \n$}
%   \AxiomC{$\pder {h_0}$}
%   \noLine
%   \UnaryInfC{$\sn, \Sigma, \n \seqar \n$}
%   \LeftLabel{$\cut$}
%   \BinaryInfC{$\sn, \Sigma \seqar \n$}
%  \AxiomC{}
%   \LeftLabel{$\bullet$}
%   \UnaryInfC{$\sn,\Sigma \seqar \n$}
%   \AxiomC{$\pder {h_1}$}
%   \noLine
%   \UnaryInfC{$\sn, \Sigma, \n \seqar \n$}
%   \LeftLabel{$\cut$}
%   \BinaryInfC{$\sn, \Sigma \seqar \n$}
%   \LeftLabel{$\cnd$}
%   \RightLabel{$\bullet$}
%   \TrinaryInfC{$\underline{\sn}, \Sigma \seqar \n$}
%   \end{prooftree}
%   It is easy to check that  $\pder f$ is in $\cbc$.
%   \end{proof}

% \begin{rem}\label{rem:auxiliary-functions}
%  As an immediate consequence of the above theorem, the following functions are in $\cbc$:
%  \allowdisplaybreaks
% \begin{align*}
%         \dot-(0; y)&=0 \\
%         \dot-(\succ i(x); y)&=\pred(; \dot-( x;y))\\ 
%         \\
%          +(0; y)&=y\\
%          +(\succ i(x); y)&=\succ i(; +( x;y))
%         \\ \\
%       \max_2(x,y;)&=+(\dot-(x;y);y)\\
%      \max_{k+1}(x_1, \ldots, x_{k+1};)&= \max_2(\max_{k}(x_1, \ldots, x_{k};), x_{k+1};)\qquad \qquad  (k \geq 2)
%       \end{align*}
% \end{rem}


We can also show that $\ncbc$ is complete for elementary functions by  simulating  our nested algebra $\nbc$. 
First, we need to introduce the notion of \emph{oracle} for coderivations.\footnote{Notice that a similar notational convention discussed in~\Cref{rem:variables-vs-functions} applies \emph{mutatis mutandis} to~\Cref{defn:derivation-oracles}.}





\begin{defn}
[Oracles for coderivations]
\label{defn:derivation-oracles}
Let $\vec{a}=a_1,\ldots, a_n$ be a set of safe-normal functions. A $\bcnorec(\vec a)$-coderivation is just a usual $\bcnorec$-coderivation that may use initial sequents of the form
$\vlinf{a_i}{}{\sn^{n_i},\n^{m_i}\seqar \n }{}$, when $a_i $ takes $n_i$ normal and $m_i$ safe inputs. 
We write:
\begin{equation*}
    \toks0={0.5}
    \vlderivation{
    \small
\vltrf{\mathcal{D}(\vec a)}{\Gamma\seqar A}{\vlhy{}}{\scriptsize \vlin{a_i}{i}{\sn^{n_i},\n^{m_i}\seqar \n }{\vlhy{\ }}}{\vlhy{\ }}{\the\toks0}
}
\end{equation*}
for a coderivation $\der$ whose initial sequents are among the initial sequents  $\vlinf{a_i}{}{\sn^{n_i},\n^{m_i}\seqar \n }{}$, with $i=1,\dots, n$.  We write $ \ncbc(\vec{a})$ for the set of $\ncbc$-coderivations  with initial functions $\vec{a}$. 
We may sometimes omit indicating some oracles $\vec a$ if it is clear from context.

The semantics of such coderivations  and the notion of $\ncbc(\vec a)$-definability are as expected, with coderivations representing functions over the oracles $\vec a$, and   $\denot{\der(\vec a)}\in \ncbc(\vec{a})$ denoting the induced interpretation of $\der(\vec a)$. 
\end{defn}




Before giving our main completeness result for $\ncbc$, we need the following lemma allowing us to `pass' parameters to oracle calls.
It is similar to the notion $\der^{\vec \rho}$ from \cite[Lemma 42]{Das2021}, only we must give a more refined argument due to the unavailability of contraction in our system.



\begin{lem}\label{lem:to-prove-ncbc-complete-for-nbc}
Let $\der(a)$ be a regular coderivation over initial sequents $\vec a, a$ of the form:
\[
\toks0={0.5}
\vlderivation{
\small
\vltrf{\der(a)}{\sn, \overset{k}{\ldots},\sn, \Gamma  \seqar \n}{\scriptsize \vlin{a}{}{    \Delta \seqar \n}{\vlhy{\ }}}{\vlhy{}}{\scriptsize \vlin{a_i}{i}{ \Delta_i \seqar \n}{\vlhy{\ }}}{\the\toks0}
}
\]
where $\Gamma$ and $\Delta$ are lists of non-modal formulas, and   the path from the conclusion to each initial sequent $a$ does not contain $\cut_\sq$-steps, $\sql$-steps and the leftmost premise of a $\cnd_\sq$-step. 
Then, there exists an $a^*$ and a regular coderivation   $\der^*( a^*)$  with shape:
\[
\toks0={0.5}
\vlderivation{
\small
\vltrf{\der^*( a^*)}{\sn, \overset{k}{\ldots},\sn, \Gamma \seqar \n}{\scriptsize \vlin{a^*}{}{ \sn, \overset{k}{\ldots},\sn,\Delta  \seqar \n}{\vlhy{\ }}}{\vlhy{}}{\scriptsize \vlin{a_i}{i}{  \Delta_i \seqar \n}{\vlhy{\ }}}{\the\toks0}
}
\]
such that:
\begin{itemize}
\item  $\denot{\der^*( a^*)}( \vec x; \vec y)=\denot 
     {\der( a( \vec x))}(\vec x; \vec y) $;
    \item there exists a $\sn$-thread from the $j$\textsuperscript{th} $\sn$ in the LHS of the end-sequent to the $j$\textsuperscript{th} $\sn$ in the context of any occurrence of the initial sequent $a^*$ in $\der^*( a^*)$, for $1\leq j\leq k$.
    \end{itemize}
    Moreover, if $\der(a)$ is progressing, safe or left-leaning, then $\der^*(a^*)$ is also progressing, safe or left-leaning, respectively.
\end{lem}
\begin{proof}[Proof sketch]
Let us consider a path $B$ from the root of $\der( a)$ to an occurrence of the initial sequent $a$. 
Since the conclusion of   $\der( a)$ contains $k$ modal formulas, and $B$ cannot cross $\cut_\sq$-steps, $B$ contains exactly $k$ $\sn$-threads, and all such threads  start from the root. 
Moreover, since $a$ has only non-modal formulas and $B$ cannot cross $\sql$-steps or the leftmost premise of a $\cnd_\sq$-step, we conclude that each of the $k$  $\sn$-threads must end in the principal formula of a $\wk_\sq$-step. 
For each $j$ we remove the corresponding  $\wk_\sq$-step in $B$  we add an extra $\sn$ to the antecedent of all higher sequents in $B$ (this operation may require us to introduce weakening steps for other branches of the proof).  
By repeatedly applying the above procedure for each possible path  from the root of $\der(a)$ to an initial sequent $a$ we obtain a coderivation with the desired properties.
\end{proof}
% \begin{proof}[Proof sketch]
% We introduce coderivations having special initial sequents for oracles. The proof then relies on  an auxiliary lemma allowing us to `pass' parameters to oracle calls, and is 
%  similar to the notion $\der^{\vec \rho}$ from \cite[Lemma 42]{Das2021}, only we must give a more refined argument due to the unavailability of contraction in our system.
% \end{proof}



Finally we can give our main simulation result for $\ncbc$:


\begin{thm}
\label{thm:nbc-in-ncbc}
If $f(\vec x;\vec y) \in \nbc$ then $f(\vec x; \vec y)\in \ncbc$.
\end{thm}
\begin{proof}
We show by induction on $f(\vec x;\vec y)\in \nbc(\vec a)$ that there is a $ \ncbc(\vec{a})$-coderivation $\der_f$ such that:
\begin{enumerate}
    \item $\denot{\der_f} (\vec x;\vec y) = f(\vec x;\vec y)$;
    \item \label{enum:construction} the path from the conclusion of $\der_{f}$ to each initial sequent $a_i$ does not  contain $\cut_\sq$-steps, $\sql$-steps and the leftmost premise of a $\cnd_\sq$-step.
\end{enumerate}
When $f(\vec x;\vec y)$ is an initial function the definition of $\der_f$ is  straightforward, as $\vec a= \emptyset$ and $\der_f$ can be constructed as in~\Cref{prop:bc-type-system-characterisation} . 
If $f(\vec x;\vec y)=a_i(\vec x;\vec y)$ then $\der_f$ is the  initial sequent  $
 \vlinf{a_i}{}{\sq \vec \n , \vec \n\seqar \n }{}
$. 

Suppose that
 $f(\vec x; \vec y)= h(\vec x, g(\vec x;); \vec y)$ with $g(\vec x;)\in \nbc(\emptyset)$
 and
 $h(\vec x, z; \vec y)\in \nbc(\emptyset)$.  Then $f$ can be  $\ncbc(\emptyset)$-defined by:
\[
\vlderivation{
\small
\vliin{\cut_\sq}{}
{\sq \vec \n, \vec \n \seqar \n}
{\vlin{\sqr}{}{\sq \vec \n\seqar \red{\sn}}{\vltr{\der_{g}}{\sq \vec \n\seqar \n}{\vlhy{\ }}{\scriptsize  \vlhy{\ \ \  \emptyset\ \ \ }  }{\vlhy{\ }}}}
{\vltr{\der_{h}}{\red{\sn}, \sq \vec \n, \vec \n \seqar \n}{\vlhy{\ }}{\scriptsize  \vlhy{\ \ \  \emptyset\ \ \ } }{\vlhy{\ }}}
}
\]
Note that, while we introduce a $\cut_\sq$ here, there crucially remains no $\cut_\sq$ between the conclusion and an oracle sequent, thanks to the condition that $g$ and $h$ are over no oracles.

Suppose $f(\vec x; \vec y)= h(\vec x; \vec y, g(\vec x; \vec y))$. 
Then $f$ is $\ncbc(\vec a)$-defined by:
\[
\vlderivation{
\small
\vliin{\cut_\n}{}
{\sq \vec \n, \vec \n \seqar \n}
{\vltr{\der_g}{\sq \vec \n, \vec \n \seqar \red{\n}}{\vlhy{\ }}{\scriptsize \vlin{a_i}{i}{\sq \vec \n, \vec \n \seqar \n}{\vlhy{\ }}}{\vlhy{\ }}}
{\vltr{\der_h}{ \sq \vec \n, \vec \n, \red{\n} \seqar \n}{\vlhy{\ }}{\scriptsize  \vlin{a_i}{i}{\sq \vec \n,  \vec \n \seqar \n}{\vlhy{\ }}}{\vlhy{\ }}}
}
\]

Finally, suppose that $f(x, \vec x; \vec y)$ is obtained from $g(\vec x; \vec y)$ {over $\emptyset$}, and  $h_j( a)(x,\vec x;\vec y)$ ($j=0,1$) over $a,\vec a$ by $\snrec$. Then,  $f(0, \vec x; \vec y) = g(\vec x;\vec y)$ and $f(\succ j x, \vec x; \vec y) = h_j(\lambda \vec v . f(x,\vec x;\vec v))(x,\vec x;\vec y)$. 
% By Remark~\ref{rem:fast-growing} we can assume that $g(\vec x; \vec y)\in \nbc(\emptyset)$. 
Note that $a$ has same type as $\lambda \vec v . f(x,\vec x;\vec v)$, so it is a function taking safe arguments only. 
Thus by induction hypothesis, $h_j(a)(x,\vec x;\vec y)$ is $\ncbc(\vec a, a)$-defined by:
\[
\vlderivation{
\small
\vltr{\der_{h_j}(a)}{ \sn, \sq \vec \n, \vec \n \seqar \n}{\scriptsize \vlin{a}{}{ \vec \n \seqar \n}{\vlhy{\ }}}{\vlhy{}}{\scriptsize \vlin{a_i}{i}{   \sq \vec \n , \vec \n \seqar \n}{\vlhy{\ }}}
}
\]
where the path from the conclusion of $\der_{h_j}( a)$ to each initial sequent $a_i$ does not  contain $\cut_\sq$-steps, $\sql$-steps and the leftmost premise of a $\cnd_\sq$-step. By Lemma~\ref{lem:to-prove-ncbc-complete-for-nbc} we obtain the following coderivation:
\[
\vlderivation{
\small
\vltr{\der_{h_j}^*( a^*)}{\sn, \sq \vec \n, \vec \n \seqar \n}{\scriptsize \vlin{a^*}{}{\sn, \sq \vec \n , \vec \n \seqar \n}{\vlhy{\ }}}{\vlhy{}}{\scriptsize \vlin{a_i}{i}{\sn, \sq \vec \n, \vec \n \seqar \n}{\vlhy{\ }}}
}
\]
where $\denot{\der_{h_j}^*( a^*)}(x, \vec x; \vec y)=\denot {\der_{h_j}( a(x, \vec x))}(x, \vec x; \vec y) $ and there exists a $\sn$-thread from the $k$-th modal formula in the antecedent of the conclusion to the $k$-th modal formula in the antecedent of any $a^*$ initial sequent in $\der^*(a^*)$.  
We define $\der_f$ as follows: 
\[
\vlderivation{
\small
\vliin{\cnd_{\sq}}{\bullet}{\underline{\blue{\sn}},  \sq \vec \n, \vec \n \seqar \n}
{
\vltr{\der_g}{ \sq \vec \n, \vec \n \seqar \n}{\vlhy{\ }}{\vlhy{\scriptsize  \ \emptyset\ \ \ }}{\vlhy{}}
}
{\vltr{\der_{h_j}^*{(\der_f)}}{\blue{\sn},  \sq \vec \n, \vec \n \seqar \n}{\scriptsize \vlin{\cnd_{\sq}}{\bullet }{\blue{\sn},  \sq \vec \n, \vec \n \seqar \n}{\vlhy{\vdots }}}{\vlhy{}}{\scriptsize \vlin{a_i}{i}{\sn,  \sq \vec \n, \vec \n \seqar \n}{\vlhy{\ }}}}
}
\]
where we identify the sub-coderivations corresponding to the second and the third premises of the conditional rule. 
By construction and induction hypothesis, the above coderivation is regular and safe. 
{To see that it is progressing, note that any infinite branch $B$ either hits $\bullet$ infinitely often, in which case there is a progressing thread along the blue $\blue{\sn}$, by the properties of $\der_{h_j}^*$ inherited from \Cref{lem:to-prove-ncbc-complete-for-nbc}, or $B$ shares a tail with an infinite branch of $\der_g$ or $\der_{h_j}^*(a^*)$, which are progressing by the inductive hypotheses.}

We show that $\der_f$ $\ncbc(\vec a)$-defines $f$  by induction on $x$. For the base case, $ \denot{\der^*_f}(0, \vec x; \vec y)=   \denot{\der_g}(\vec x; \vec y)= g(\vec x; \vec y)= f(0, \vec x; \vec y)$. For the inductive step:
\[
\renewcommand{\arraystretch}{1.5}
\arraycolsep=2pt
\begin{array}{rcl}
        \denot{\der^*_f}(\succ jx, \vec x; \vec y)&=&  \denot{\der^*_{h_j}(\lambda \vec v.  \denot{\der^*_f}(x, \vec x; \vec v) )}(x, \vec x; \vec y)  \\ 
     & =& \denot{\der^*_{h_j}( \lambda \vec v.  f(x, \vec x; \vec v) )}(x, \vec x; \vec y) \\
       & =& \denot{\der_{h_j}( \lambda \vec v.  f(x, \vec x; \vec v) )}(x, \vec x; \vec y) \\
       & =& h_j( \lambda \vec v.  f(x, \vec x; \vec v) )(x, \vec x; \vec y)\\
       &=& f(\succ j x, \vec x; \vec y) .
\end{array}
\]
This completes the proof.
\end{proof}

% \red{====================================================}

% To show that $\ncbc$ is complete for elementary functions, we shall simulate our nested algebra $\nbc$. First, we need to introduce the notion of \emph{oracle} for $\bcnorec$-coderivations.

% \begin{defn}
% [Oracles for coderivations]
% Let $\vec{a}=a_1,\ldots, a_n$ be a set of safe-normal functions. A $\bcnorec(\vec a)$-coderivations is just a usual $\bcnorec$-coderivation that may use initial sequents of the form
% $\vlinf{a_i}{}{\sn^{n_i},\n^{m_i}\seqar \n }{}$, when $a_i $ takes $n_i$ normal and $m_i$ safe inputs. 
% We write:
% \begin{equation*}
%     \vlderivation{
%     \small
% \vltr{\mathcal{D}(\vec a)}{\Gamma\seqar A}{\vlhy{}}{\scriptsize \vlin{a_i}{i}{\sn^{n_i},\n^{m_i}\seqar \n }{\vlhy{\ }}}{\vlhy{\ }}
% }
% \end{equation*}
% for a coderivation $\der$ whose initial sequents are among the initial sequents  $\vlinf{a_i}{}{\sn^{n_i},\n^{m_i}\seqar \n }{}$, with $i=1,\dots, n$.  We write $ \ncbc(\vec{a})$ for the set of $\ncbc$-coderivations  with initial functions $\vec{a}$. 
% We may sometimes omit indicating some oracles $\vec a$ if it is clear from context.

% The semantics of such coderivations  and the notion of $\ncbc(\vec a)$-definability are as expected, with coderivations representing functions over the oracles $\vec a$, and   $\denot{\der(\vec a)}\in \ncbc(\vec{a})$ denoting the induced interpretation of $\der(\vec a)$. 
% \end{defn}




% {Before giving our main completeness result for $\ncbc$, we need the following lemma allowing us to `pass' parameters to oracle calls.
% It is similar to the notion $\der^{\vec \rho}$ from \cite[Lemma 42]{Das2021}, only we must give a more refined argument due to the unavailability of contraction in our system.}



% \begin{lem}\label{lem:to-prove-ncbc-complete-for-nbc}
% Let $\der(a)$ be a regular coderivation over initial sequents $\vec a, a$ with shape:
% \[
% \vlderivation{
% \small
% \vltr{\der(a)}{\sn, \overset{k}{\ldots},\sn, \Gamma  \seqar \n}{\scriptsize \vlin{a}{}{    \Delta \seqar \n}{\vlhy{\ }}}{\vlhy{}}{\scriptsize \vlin{a_i}{i}{ \Delta_i \seqar \n}{\vlhy{\ }}}
% }
% \]
% where $\Gamma$ and $\Delta$ are lists of non-modal formulas, and   the path from the conclusion to each initial sequent $a$ does not contain $\cut_\sq$-steps, $\sql$-steps and the leftmost premise of a $\cnd_\sq$-step. 
% Then, there exists an $a^*$ and a regular coderivation   $\der^*( a^*)$  with shape:
% \[
% \vlderivation{
% \small
% \vltr{\der^*( a^*)}{\sn, \overset{k}{\ldots},\sn, \Gamma \seqar \n}{\scriptsize \vlin{a^*}{}{ \sn, \overset{k}{\ldots},\sn,\Delta  \seqar \n}{\vlhy{\ }}}{\vlhy{}}{\scriptsize \vlin{a_i}{i}{  \Delta_i \seqar \n}{\vlhy{\ }}}
% }
% \]
% such that:
% \begin{itemize}
% \item  $\denot{\der^*( a^*)}( \vec x; \vec y)=\denot 
%      {\der( a( \vec x))}(\vec x; \vec y) $;
%     \item there exists a $\sn$-thread from the $j$\textsuperscript{th} modal formula in the context of the conclusion to the $j$\textsuperscript{th} modal formula in the context of any occurrence of the initial sequent $a^*$ in $\der^*( a^*)$, for $1\leq j\leq k$.
%     \end{itemize}
%     Moreover, if $\der(a)$ is progressing, safe or left-leaning, then $\der^*(a^*)$ is also progressing, safe or left-leaning, respectively.
% \end{lem}













\subsection{The Translation Lemma} \label{sec:translation} 





%  \red{Recalling the notation in~\Cref{defn:saferec-pp} we can state the following result.}\anupam{I don't think you need this sentence.}
 
The goal of this section is to prove the following theorem, which shows  a translation of  $\ncbc$-coderivations into   functions of $\nbcpp$   mapping, in particular,    $\cbc$-derivations into functions of $\bcpp$:

 
\begin{thm}[Translation Lemma]\label{lem:short-translation-lemma}
Let $\der$ be a  $\ncbc$-coderivation. Then, there exists a set of $n$ functions $(f_i)_{1 \leq i \leq n}$ such that $f_{1}= \model{\der}$ and, for all $i$:
\begin{equation}\label{eqn:naive1}
\arraycolsep=2pt
\begin{array}{rcl}
 f_i(\vec{x}; \vec{y})   & = & h_i\left( \lambda \vec{u}\subset \vec{x},\lambda \vec v. f_j(\vec{u};\vec v)\right)_{1 \leq j \leq n}(\vec{x}; \vec{y})
\end{array}
\end{equation}
where  $h_i \in \nbcpp (a_i)_{1 \leq i \leq n}$. Moreover, if    $\der$ is a $\cbc$-coderivation  then, for all $i$:
\begin{equation}\label{eqn:naive2}
\arraycolsep=2pt
\begin{array}{rcl}
 f_i(\vec{x}; \vec{y})   & = & h_i\left( \lambda \vec{u}\subset \vec{x},\lambda \vec v\subseteq \vec y. f_j(\vec{u};\vec v)\right)_{1 \leq j \leq n}(\vec{x}; \vec{y})
\end{array}
\end{equation}
and  $h_i \in \bcpp  (a_i)_{1 \leq i \leq n}$.
\end{thm} 


% \begin{figure*}
%     \centering
% \subfloat[$\der$ as a definition tree.] { 
%  \label{fig:tree1}
%   \quad  $
%  \phantom{
%    \small
%    \vlderivation{
% \vlin{}{\bullet_1}{\denot{\der_{v}}}
%   {
%   \vlin{\cut_\sq}{}{ 
%   \underset{\vdots}{\denot{\der_{u}}}}{
%     \vlin{}{}{\denot{\der_{u'}}}{
%     \vlin{}{\bullet_n}{ \underset{\vdots}{\denot{\der_{v}}}}{\vlhy{}}
%     }
%      }
%     }
%     }
%     }
%   \small
%      \toks0={0.5}
% \vlderivation{
% \vltrf{}{\denot{\der}(=\denot{\der_\epsilon})}
%   {\vlin{}{\bullet_1}{\denot{\der_{v_1}}}{\vlhy{}}}
%   { \vlhy{\ldots}}
%   {\vlin{}{\bullet_n}{\denot{\der_{v_n}}}{\vlhy{}}}{\the\toks0}
%   }
%   \phantom{
%    \small
% \vlderivation{
% \vlin{}{\bullet}{\denot{\der_{v}}}
%   {
%   \vlin{\cut_\sq}{}{ 
%   \underset{\vdots}{\denot{\der_{u}}}}{
%     \vlin{}{}{\denot{\der_{u'}}}{
%     \vlin{}{\bullet}{ \underset{\vdots}{\denot{\der_{v}}}}{\vlhy{}}
%     }
%      }
%     }
%     }
% }
% $ 
% }
%   \subfloat[A $\cut_{\sq}$ violating safety.]{
%   \label{fig:tree2}  $
% \phantom{
%    \small
% \vlderivation{
% \vlin{}{\bullet}{\denot{\der_{v}}}
%   {
%   \vlin{\cut_\sq}{}{ 
%   \underset{\vdots}{\denot{\der_{u}}}}{
%     \vlin{}{}{\denot{\der_{u'}}}{
%     \vlin{}{\bullet}{ \underset{\vdots}{\denot{\der_{v}}}}{\vlhy{}}
%     }
%      }
%     }
%     }
%     }
%   \small
% \vlderivation{
% \vlin{}{\bullet}{\denot{\der_{v}}}
%   {
%   \vlin{\cut_\sq}{}{ 
%   \underset{\vdots}{\denot{\der_{u}}}}{
%     \vlin{}{}{\denot{\der_{u'}}}{
%     \vlin{}{\bullet}{ \underset{\vdots}{\denot{\der_{v}}}}{\vlhy{}}
%     }
%      }
%     }
% }
% \phantom{
%    \small
% \vlderivation{
% \vlin{}{\bullet}{\denot{\der_{v}}}
%   {
%   \vlin{\cut_\sq}{}{ 
%   \underset{\vdots}{\denot{\der_{u}}}}{
%     \vlin{}{}{\denot{\der_{u'}}}{
%     \vlin{}{\bullet}{ \underset{\vdots}{\denot{\der_{v}}}}{\vlhy{}}
%     }
%      }
%     }
%     }
%     }
% $ 
%   }
%   \subfloat[A $\cut_{\n}$ violating left-leaning.]{
%   \label{fig:tree3}
% \quad  $ 
% \phantom{
%    \small
% \vlderivation{
% \vlin{}{\bullet}{\denot{\der_{v}}}
%   {
%   \vlin{\cut_\sq}{}{ 
%   \underset{\vdots}{\denot{\der_{u}}}}{
%     \vlin{}{}{\denot{\der_{u'}}}{
%     \vlin{}{\bullet}{ \underset{\vdots}{\denot{\der_{v}}}}{\vlhy{}}
%     }
%      }
%     }
%     }
%     }
% \small
% \vlderivation{
% \vlin{}{\bullet}{\denot{\der_{v}}}
%   {
%   \vliin{\cut_\n}{}{ 
%   \underset{\vdots}{\denot{\der_{u}}}}{\vltr{}{\denot{\der_{u'}}}{\vlhy{}}{\vlhy{}}{\vlhy{}}}{
%     \vlin{}{}{\denot{\der_{u''}}}{
%     \vlin{}{\bullet}{ \underset{\vdots}{\denot{\der_{v}}}}{\vlhy{}}
%     }
%      }
%     }
% }
% \phantom{
%    \small
% \vlderivation{
% \vlin{}{\bullet}{\denot{\der_{v}}}
%   {
%   \vlin{\cut_\sq}{}{ 
%   \underset{\vdots}{\denot{\der_{u}}}}{
%     \vlin{}{}{\denot{\der_{u'}}}{
%     \vlin{}{\bullet}{ \underset{\vdots}{\denot{\der_{v}}}}{\vlhy{}}
%     }
%      }
%     }
%     }
%     }
% $\quad  
%   }
%     \caption{Structure of the definition tree of   $\der$.}
%     \label{fig:structure-definition-tree}
% \end{figure*}

This would conclude  our characterisation of $\cbc $ and $\ncbc$ in terms of computational complexity. 

A formal proof of the translation lemma requires a more general statement (i.e.,~\Cref{lem:translation}) as well as a few technical notions and intermediate results. 
% To illustrate some key ideas of that proof, we start with discussing a high-level  and simplified argument for~\Cref{lem:short-translation-lemma} that uses the informal notion of `definition tree'.
First, we observe that a regular coderivation can be naturally seen as a finite tree with `backpointers', a representation known as \emph{cycle normal form}, cf.~\cite{Brotherston05,brotherston2011sequent}.  

% \color{gray}{
% \begin{proof}[Proof idea]
% By~\Cref{defn:semantics-coderivations} there exists a system of equations $S_\der$ containing, for each node $v$ of $\der$,  an equation  that defines the function $\denot{\der_v}$ in terms of the functions  $\denot{\der_u}$ with $u$  immediately above $v$. Moreover,  since $\der$ is regular, by~\Cref{rem:regular-finite-equation} we may assume $S_\der$ is a \emph{finite} system of equations defining  $\denot{\der}$ (i.e.~$\denot{\der_\epsilon}$). It is then suggestive to represent  $\der$ as a `definition tree', by replacing each node $v$ with the corresponding function $\denot{\der_v}$ as in~\Cref{fig:tree1},  where each function in the tree is defined by one of the equations of $S_\der$. This representation of $\der$ has the advantage of  highlighting the inter-dependencies of the functions defined in $S_\der$. 














% We now discuss how the proof-theoretical properties of $\der$ impose conditions on the equations of  $S_\der$:

% \begin{enumerate}
%     \item \label{enum:sketch1} By safety condition, the definition tree cannot contain a branch crossed by infinitely many $\cut_\sq$ rules   as in~\Cref{fig:tree2}. Hence,  whenever $S_\der$ has functions $\denot{\der_{u}}$ and  $\denot{\der_{v}}$ that are inter-dependent, no equation of $S_\der$ can define one of these functions by means of a composition along its normal parameters. 
%     By  inspecting~\Cref{defn:semantics-coderivations} this implies that,   
% whenever an equation in $S_\der$  defines a function $\denot{\der_{u}}$ by affecting its normal parameters, $u$ must be the conclusion of an instance of $\cnd_\sq$, i.e.~this equation must be of the form   $\denot{\der_{u}}(\succ i x, \vec z; \vec y)= \denot{\der_{v}}(x, \vec z; \vec y)$, for some $v$.
%     \item \label{enum:sketch2} By progressiveness of $\der$, there always exists an  occurrence of $\cnd_\sq$ in-between two `backpointers'  $\bullet_i$ of~\Cref{fig:tree1}.
% \item \label{enum:sketch3} If  $\der$ is left-leaning, then the tree in~\Cref{fig:tree1} cannot contain the instance of $\cut_\n$  in~\Cref{fig:tree3}. Hence,  whenever the system of equations for $\der$ contains two functions $\denot{\der_{u}}$ and  $\denot{\der_{v}}$ that are inter-dependent, no equation can define one of these functions by means of a composition along its safe parameters. By  inspecting~\Cref{defn:semantics-coderivations} this  implies  that,   
% whenever an equation in $S_\der$  defines a function $\denot{\der_{u}}$ by affecting its safe parameters, $u$ must be  the conclusion of an instance of $\cnd_\n$, i.e.~the equation must be of the form   $\denot{\der_{u}}(\vec x; \vec y, \succ i z)= \denot{\der_{v}}( \vec z; \vec y, z)$, for some $v$.
% \end{enumerate}

% We can now simplify the equations in $S_\der$ to  obtain  a  `minimal' system of equations $S^*_\der$,  where each equation  defines  a function in  $(f_1)_{1 \leq i \leq n}$,  with $f_1=\denot{\der}$ and $f_i=\denot{\der_{v_i}}$ (see~\Cref{fig:tree1}).  More precisely, for all $1 \leq i \leq n$ there exists $h_i \in \nbcpp(  a_i)_{1 \leq i \leq n}$, for some  oracles  $(a_i)_{1 \leq i \leq n}$, such that  the following equation is in  $S^*_\der$:
% \begin{equation*}
% \arraycolsep=2pt
% \begin{array}{rcl}
%  f_i(\vec{x}; \vec{y})   & = & h_i\left( \lambda \vec{u}.\lambda \vec v. f_j(\vec{u};\vec v)\right)_{1 \leq j \leq n}(\vec{x}; \vec{y})
% \end{array}
% \end{equation*}
% % where lambda abstraction is meta-level notation  for expressing possible nesting of  functions.\anupam{skip this phrase?}
% By point~\ref{enum:sketch1}, each such equation can be rewritten as:  
% \begin{equation}\label{eqn:equation}
% \arraycolsep=2pt
% \begin{array}{rcl}
%  f_i(\vec{x}; \vec{y})   & = & h_i\left( \lambda \vec{u}\subseteq \vec{x}.\lambda \vec v. f_j(\vec{u};\vec v)\right)_{1 \leq j \leq n}(\vec{x}; \vec{y})
% \end{array}
% \end{equation}
% Moreover, by point~\ref{enum:sketch2} the relation ${\subseteq}$  is strict (i.e.~$\vec u \subset \vec x$) when $j=i$. In particular, by point~\ref{enum:sketch3}, if $\der$ is left-leaning then $h_i \in \bcpp  (a_i)_{1 \leq i \leq n}$ and \eqref{eqn:equation} above can be rewritten as: 
% \begin{equation*}
% \arraycolsep=2pt
% \begin{array}{rcl}
%  f_i(\vec{x}; \vec{y})   & = & h_i\left( \lambda \vec{u}\subseteq \vec{x}.\lambda \vec v\subseteq \vec y. f_j(\vec{u};\vec v)\right)_{1 \leq j \leq n}(\vec{x}; \vec{y})
% \end{array}
% \end{equation*}

% It remains to show that  the relation ${\subseteq} $ is strict (i.e.~$\vec u \subset \vec x$) when $j\neq i$ in the above equations. This can be established by repeatedly applying the following operation for each  equation in $S^*_\der$ with shape~\eqref{eqn:equation}: for all $j$, if  ${\subseteq}$ is not strict in $ \lambda \vec{u}\subseteq \vec{x}.\lambda \vec v\subseteq \vec y. f_j(\vec{u};\vec v)$ (i.e.~$\vec u \subset \vec x$), replace $f_j$ with its definition  given by the corresponding equation of  $S^*_\der$. Such a procedure keeps `unfolding' equations, and terminates as a consequence of point~\ref{enum:sketch2}.\qedhere
% \end{proof}
% \color{black}
 
%  \red{
%  It is worth stressing that the above sketched proof can be made more formal. It requires to   consider the  `cycle normal form' of  $\der$, which is essentially a finite tree with backpointers. This allows us to inductively associate with  each node $v$ of the cycle normal form a system of equations defining a given set of functions $f_{i_1}, \ldots, f_{i_k}\subseteq \{f_1, \ldots, f_n \}$, where  $f_{i_1}=\denot{\der_v}$, with the help of  the functions in $\{f_1, \ldots, f_n \}\setminus \{ f_{i_1}, \ldots, f_{i_k}\}$ seen as auxiliary oracles.
% }\anupam{I think just omit all this in red.}




% . We first identify a finite set of nodes $T_\der$ of $\der$ which is downward closed (w.r.t.~$\sqsubseteq$).   Then we show by induction on the height of nodes in $\der$ that any $\nu \in T_\der$ can be given a finite system of equations each one with shape~\eqref{eqn:naive1} and  such that: 
% \begin{itemize}
%     \item $f_1= \denot{\der_\nu}$,  
%     \item  $h_i \in \nbcpp(  (a_i)_{1 \leq i \leq n}, (b_j)_{1 \leq j \leq m})$, for some oracles $b_j$. 
% \end{itemize}
% In particular, if $\der$ is a $\cbc$-coderivation than the equations of the system have form~\eqref{eqn:naive2}, and $h_i \in \bcpp(  (a_i)_{1 \leq i \leq n}, (b_j)_{1 \leq j \leq m})$. Here, the need for auxiliary oracles is required, a


% We now introduce some technical notions and results required for proving~\Cref{lem:translation}. 








\begin{defn}[Cycle normal form]
Let $\der$ be a regular $\bcnorec$-coderivation. 
The  \emph{cycle normal form} (or simply \emph{cycle nf}) of $\der$ is a pair $\derrd$, where  $\rd$ is a partial self-mapping on the nodes of $\der$ whose domain of definition is denoted $\bud(\der)$ and:
\begin{enumerate}[(i)]
    \item every infinite branch of $\der$ contains some (unique) $\nu \in \bud(\der)$;
    % \item for any $\nu \in \bud(\der)$, $\nu$ is the least node (w.r.t.~$\sqsubseteq$) of $\der$  satisfying \red{both} $\rd(\nu) \sqsubset \nu$ and $\pder\nu= \pder{\rd(\nu)}$;
    \item if $\nu \in \bud(\der)$ then both $\rd(\nu) \sqsubset \nu$ and  $\der_{\rd(\nu)} = \der_\nu$;
    \item for any two distinct nodes $\mu \sqsubset \nu$ strictly below $\bud(\der)$, $\der_\mu \neq \der_\nu$
\end{enumerate}
% \anupam{Second condition still seems wrong to me: `least node s.t.\ some property of $\rd(\nu)$' is consistent with not taking the first repetition along a branch, since the first repetition is not in $\rd$'s domain.} 
% \anupam{Can we replace second condition with the two conditions: (a) if $\nu \in \bud(\der)$ then $\der_{\rd(\nu)} = \der_\nu$; and (b) for any two distinct nodes $\mu \sqsubset \nu$ below $\bud(\der)$, $\der_\mu \neq \der_\nu$. This then matches your explanation in the following remark.}

We call any $\nu \in \bud(\der)$ a \emph{bud}, and $\rd(\nu)$ its \emph{companion}. A \emph{terminal} node is either a leaf of $\der$ or a bud. The set of nodes of $\der$ bounded above by a terminal node is denoted $T_\der$.  
Given a node $\nu\in T_\der$, we  define $\bud_\nu(\der)$ as the restriction of buds to those above $\nu$.
\end{defn}

\begin{rem}\label{rem:canonicalnodesfinite}
The cycle normal form of a regular coderivation $\der$ always exists, as  by definition any infinite branch contains a node $\nu$ such that $\pder \nu=\pder \mu$ for some node $\mu$ below $\nu$. 
$\bud(\der)$ is designed to consist of just the \emph{least} such nodes, so that by construction the  cycle normal form is unique. 
Note that $\bud(\der)$ must form an $\sqsubset$-antichain: if $\mu, \nu \in \bud(\der)$ with $\mu\sqsubset \nu$, then $\rd(\mu)\sqsubset \mu $ are below $\bud(\der)$ but we have $\der_{\rd(\mu)} = \der_\mu$ by point (ii) above, contradicting point (iii). 

Also, notice that any branch of $\der$ contains a leaf of $T_\der$. Moreover, since $\bud(\der)$ is an antichain, the leaves of $T_\der$ defines a ‘bar’ across  $\der$,
and so $T_\der$ is a finite tree. 
% Finally, as this tree has finite branches, it must be  \emph{finite}  by weak K\"{o}nig's Lemma.
\end{rem}


% We shall  represent cycle nfs as \emph{finite}  prooftrees (thanks to Remark~\ref{rem:canonicalnodesfinite}) with inference rules  labelled  with coloured bullets, in such a way that bullets with the same colour represent  endoints of the same backlink. 


The following proposition allows us to reformulate   progressiveness, safety and left-leaning conditions for cycle normal forms. 
% \anupam{Could we put this earlier and/or refer to the proofs of NL-decidability of certain conditions?}
\begin{prop}\label{prop:structureofcycles}
Let $\der$ be a regular   $\bcnorec$-coderivation with cycle nf $\derrd$. 
For any  $\nu \in \bud(\der)$, the (finite) path $\pi$ from $\rd(\nu)$ to $\nu$ satisfies:
\begin{enumerate}
\item \label{enum:budcompanion0} if  $\der$ is progressing, $\pi$ must contain the conclusion of an instance of $\cnd_{\sn}$;
    \item \label{enum:budcompanion1} if  $\der$ is a $ \ncbc$-coderivation, $\pi$ cannot  contain  the conclusion of  $\cut_{\sq}$, $\sql$,  $\wk_{\sq}$, and  the   leftmost premise of  $\cnd_{\sq}$;
    \item \label{enum:budcompanion3} if $\der$ is a $\cbc$-coderivation,  $\pi$  cannot contain the conclusion of $\wk_{\n}$, the leftmost premise of  $\cnd_{\n}$, and the rightmost premise of $\cut_{\n}$.
\end{enumerate}
\end{prop}
\begin{proof}
By definition of cycle nf, each path from $\rd(\nu)$ to $\nu$ in $\derrd$ is contained in a branch of $\der$ such that each rule instance in the former appears  infinitely many times in the latter. Hence:
\begin{enumerate}
\item[(i)]  if $\der$ is progressing,  the path contains the conclusion of an instance of $\cnd_{\sq}$;
\item[(ii)]  if $\der$ is safe, the path   cannot contain the conclusion of a $\cut_{\sq}$ rule;
\item[(iii)] if $\der$ is left-leaning, the path    cannot contain the rightmost premise of a $\cut_{\n}$ rule.
\end{enumerate}
This shows point~\ref{enum:budcompanion0}. Let us consider  point~\ref{enum:budcompanion1}.  By point~(ii), if $\der$ is safe then,  going from a node $\mu$ of the path to each of its children $\mu'$,  the number of modal formulas in the context of the corresponding sequents cannot increase.  Moreover, the only cases where this number strictly decreases is when $\mu$ is the conclusion of $\sql$,  $\wk_{\sn}$, or  when $\mu'$ is  the leftmost premise of $\cnd_{\sq}$. Since $\rd(\nu)$ and $\nu$ must be labelled with the same sequent, all such cases are impossible. As for  point~\ref{enum:budcompanion3} we notice that,   by point~(iii) and the above reasoning, if $\der$ is safe and left-leaning then,  going from a node $\mu$ of the path to each of its children $\mu'$,  the number of non-modal formulas in the context of the corresponding sequents cannot increase.  Moreover, the only cases where this number strictly decreases is when $\mu$ is the conclusion of $\wk_{\n}$, or  when $\mu'$ is  the leftmost premise of $\cnd_{\n}$. Since $\rd(\nu)$ and $\nu$ must be labelled with the same sequent, all such cases are impossible. 
\end{proof} 

In what follows we shall indicate circularities in cycle nfs explicitly by extending both $\ncbc$ and  $\cbc$ with a new inference rule called $\com$: 
\[
\vlinf{\com}{X}{\Gamma \Rightarrow A}{\Gamma \Rightarrow A}
\]
where $X$ is a finite set of nodes. 
In this presentation, we insist that each companion $\nu$ of the cycle nf $\derrd$ is always the conclusion of an instance of $\com$, where $X$ denotes the set of buds $\nu'$ such that $\rd(\nu')=\nu$.  
{This expedient will allow us to treat the case of companions separately, avoiding repeating the same argument for each rule.}

% formally 
% distinguish cases when a node of $T_\der$ is a companion from those where it is the conclusion of a standard rule of $\bcnorec$.




To facilitate the translation, we shall define two disjoint  sets  $\close{\nu}$ and $\open{\nu}$. Intuitively, given a cycle nf $\derrd$ 
 and  $\nu\in T_\der$,  $\close{\nu}$ is the set of  companions above  $\nu$, while   $\open{\nu}$  is the set of  buds whose companion is strictly below   $\nu$. 



\begin{defn} 
Let $\derrd$ be the  cycle normal form of a $\bcnorec$-coderivation $\der$. We define the following two sets for any $\nu\in T_\der$:
\begin{equation*}
\begin{aligned}
\close{\nu}&\dfn \{ \mu \in \rd(\bud_\nu(\der))  \ \vert \ \nu \sqsubseteq \mu \}\\
\open{\nu}&\dfn \{ \mu \in \bud_\nu(\der) \ \vert \ \rd(\mu) \sqsubset \nu \}
\end{aligned}
\end{equation*}
% Nodes in $\close \nu$ range over $\mu$, nodes in $\open \nu$ range over $\mu$.  
\end{defn}

\newcommand{\amodel}[1]{a_{#1}}

We now state a generalised version of the translation lemma and prove it in  detail. The idea, here, is to associate with each node $\hnu\in T_\der$  an instance of the scheme  $\ssnrecpp$ such that:
\begin{itemize}
    \item the scheme  simultaneously defines the  functions in  $\mathcal{C}_{\hnu}=\{ \model{\nu} \ \vert \ \nu\in \close {\hnu} \cup \{\hnu \}\}$,   with the help of an additional set of oracles $\mathcal{O}_{\hnu}=\{ \amodel{\mu} \ \vert \ \mu\in \open{\hnu}\}$
    \item if  $\mathcal{O}_{\hnu}=\emptyset$ then,    for every $f_\nu\in \mathcal{C}_{\hnu}$   it holds that $\model{\nu}= \model{\pder \nu}$, i.e.,    $f_\nu$ defines  the interpretation of $\der_{\nu}$ via a set of equations (see~\Cref{defn:semantics-coderivations}). 
\end{itemize}
In particular, when  $\hnu$ is the root of $\der$ then $\open{\hnu} = \emptyset$, and so $\model{\hnu}=\model{\pder \hnu}=\model{\der}$. Moreover, since  $\model{\hnu}$ is defined by an instance of   $\ssnrecpp$ then $\model{\der}\in \nbcpp$ by Proposition~\ref{prop:simultaneous-recursion-admissible}. As a special case,  if $\der$ is a $\cbc$-coderivation then $\model{\hnu}$ can be defined by an instance of   $\ssrecpp$, and so  $\model{\der}\in \bcpp$. 


\begin{lem}[Translation Lemma, general version]\label{lem:translation}
If $\derrd$ is the cycle nf of a $\ncbc$-coderivation $\der$ and  $\hnu\in T_\der$, 
% \begin{enumerate}
%     \item \label{enum:1} If  $\open \hnu
%    =\emptyset$ then   $\model{\pder \hnu}\in \nbcpp$. In particular, if  $\der$ is a $\cbc$-coderivation then $\model{\pder \hnu}\in \bcpp$. 
    % \item 
    % \label{enum:2} If  $\open \hnu \neq \emptyset$ 
    then $\forall \nu \in \close  \hnu \cup \{ \hnu\}$:
    \begin{equation*}
     \label{eqn:translationlemmaequation1}
    \begin{aligned}
     \model{ \nu}(\vec{x}; \vec{y})&=h_{\nu}\left(( \lambda \vec{u}\subset \vec{x},\lambda \vec v. \model{ \mu}(\vec{u};\vec v))_{\mu \in \close \hnu },( \lambda \vec{u}\subseteq \vec{x},\lambda \vec v. \amodel{ \mu}(\vec{u};\vec v))_{\mu \in  \open \hnu}\right)(\vec{x}; \vec{y})
    \end{aligned}
    \end{equation*}
    where:
    \begin{enumerate}
    \item \label{enum:a} $h_\nu \in \nbcpp( (\model{\mu})_{\mu \in \close \hnu}, (\amodel \mu)_{\mu \in \open \hnu})$, and so    $\model{ \nu}\in  \nbcpp( (\amodel{ \mu})_{\mu \in \open \hnu})$;
    \item \label{enum:b}  for all $\mu \in \open \hnu$, the order $\vec u \subseteq \vec x$ in $\lambda \vec{u}\subseteq \vec{x},\lambda \vec v. \amodel{ \mu}(\vec{u};\vec v)$ is strict  if the path from $\nu$ to $\mu$ in $T_\der$ contains the conclusion of an instance of  $\cnd_{\sq}$;
    \item \label{enum:d}  if moreover   $\der$ is a $\cbc$-coderivation  then 
    \begin{equation*}
     \begin{aligned}
     \model{ \nu}(\vec{x}; \vec{y})&=h_{\nu}\left(( \lambda \vec{u}\subset \vec{x},\lambda \vec v\subseteq \vec y. \model{ \mu}(\vec{u};\vec v))_{\mu \in \close \hnu },( \lambda \vec{u}\subseteq \vec{x},\lambda \vec v\subseteq \vec y. \amodel{ \mu}(\vec{u};\vec v))_{\mu \in  \open \hnu}\right)(\vec{x}; \vec{y})
    \end{aligned}
    \end{equation*}
    with $h_\nu \in \bcpp( (\model{\mu})_{\mu \in \close \hnu}, (\amodel \mu)_{\mu \in \open \hnu})$, and so    $\model{ \nu}\in  \bcpp( (\amodel{ \mu})_{\mu \in \open \hnu})$.
    \end{enumerate}
    % \end{enumerate}
\end{lem}


\begin{proof}
By induction on the longest distance of $\nu_0$ from a leaf of  $T_\der$. Concerning the cases where   $\nu_0$ is the conclusion of an instance of $\id$ or $\zero$ we have $\close{\hnu}= \open{\hnu}= \emptyset$, and we simply set $h_{\hnu}\dfn \id$ and $h_{\hnu}\dfn 0$ respectively. If $\hnu$ is the conclusion of a bud then $\close{\hnu}= \emptyset$ and $\open {\hnu}= \{\hnu\}$, and we simply set $h_{\hnu}\dfn \amodel{\hnu}$. The cases where   $\nu$ is the conclusion of a rule in $\{\wk_{\n},\wk_{\sq}, \sql , \exch_\n, \exch_\sq, \sqr,  \succ 0, \succ 1\}$ are straightforward. Notice that for the cases $\wk_{\n}, \wk_{\sq},\sql $ we have $\open{\hnu}=\emptyset$ by Proposition~\ref{prop:structureofcycles} points~\ref{enum:budcompanion1} and \ref{enum:budcompanion3}.

Let us now consider the case where $\hnu$ is the conclusion of an instance of   $\cut_{\sq}$ with premises $\nu_1$ and $\nu_2$. By Proposition~\ref{prop:structureofcycles}.\ref{enum:budcompanion1} we have $\open \hnu= \emptyset$. By induction hypothesis we have  $\model{ {\nu_1}}(\vec x; \vec y), \model{{\nu_2}}(\vec x, x; \vec y)\in \nbcpp$. Since the conclusion of $\pder{\nu_1}$  has modal succedent, by Proposition~\ref{prop:cutbox}  there must be a coderivation $\der^*$ such that $\model{\der^*}(\vec x;)= \model{{\nu_1}}(\vec x; \vec y)\in \nbcpp$. Moreover, by Proposition~\ref{prop:cutbox-for-left-leaning},  if $\pder{\nu_1}$ is a $\cbc$-coderivation then so too is $\der^*$.  Hence, we define $h_{ \hnu}(\vec x; \vec y)\dfn \model{{\nu_2}}(\vec x, \model{\der^*}(\vec x;); \vec y)$. It is easy to see that points~\ref{enum:a}-\ref{enum:d} hold.

Suppose now that $\hnu$ is the conclusion of an instance of $\cut_{\n}$ with premises $\nu_1$ and $\nu_2$.  Then, $\open{\hnu}=\open{\nu_1}\cup \open{\nu_2}$ and  $\close{\hnu}=\close{\nu_1}\cup \close{\nu_2}$.
By induction hypothesis on $\nu_1$ and $\nu_2$ we have:
\[
\arraycolsep=2pt
\begin{array}{rcl}
     \model{{\nu_1}}(\vec{x}; \vec{y})
    & =&
     h_{\nu_1}\left(( \lambda \vec{u}\subset \vec{x}, \lambda \vec v. \model{{\mu}}(\vec{u};\vec v))_{\mu \in \close{\nu_1}}, ( \lambda \vec{u}\subseteq \vec{x}, \lambda \vec v. \amodel{{\mu}}(\vec{u};\vec v))_{\mu \in  \open{\nu_1}}\right)(\vec{x}; \vec{y})
\\
       \model{{\nu_2}}(\vec{x};y, \vec{y})  &=& 
       h_{\nu_2}\left(( \lambda \vec{u}\subset \vec{x}, \lambda v,\vec v. \model{{\mu}}(\vec{u};v,\vec v))_{\mu \in \close{\nu_2}}, ( \lambda \vec{u}\subseteq \vec{x}, \lambda v,\vec v. \amodel{{\mu}}(\vec{u};v,\vec v))_{\mu \in \open{\nu_2}}\right)(\vec{x}; y,\vec{y})
\end{array}
\]
So that we set
\begin{equation*}
    h_{{\hnu}}(\vec{x}; \vec{y})\dfn \model{{\nu_2}}(\vec{x}; \model{{\nu_1}}(\vec{x}; \vec{y}), \vec{y}) 
\end{equation*}
% \begin{equation*}
% % \label{eqn:5}
%      \model{{\hnu}}(\vec{x}; \vec{y})=h_{\hnu}(( \lambda \vec{u}\subseteq \vec{x}, \lambda \vec v. \model{{\mu}}(\vec{u};\vec v))_{\mu \in \close{\hnu}\cup \open{\hnu}})(\vec{x}; \vec{y})
% \end{equation*}
Points~\ref{enum:a} and~\ref{enum:b}  hold by applying the induction hypothesis. Concerning point~\ref{enum:d}, notice that   if $\der$ is a $\cbc$-coderivation then   $\open{\nu_2}= \emptyset$ by Proposition~\ref{prop:structureofcycles}.\ref{enum:budcompanion3}.  By applying the induction hypothesis on $\nu_1$ and $\nu_2$, we have $ \model{{\nu_2}}(\vec{x};y, \vec{y})\in \bcpp$ and   $ \model{{\nu_1}}(\vec{x}; \vec{y})=h_{\nu_1}\left(( \lambda \vec{u}\subset \vec{x}, \lambda \vec v\subseteq \vec y. \model{{\mu}}(\vec{u};\vec v))_{\mu \in \close{\nu_1}},( \lambda \vec{u}\subseteq \vec{x}, \lambda \vec v\subseteq \vec y. \amodel{{\mu}}(\vec{u};\vec v))_{\mu \in  \open{\nu_1}}\right)(\vec{x}; \vec{y})$, so that:
\begin{equation*}
     \model{{\hnu}}(\vec{x}; \vec{y})=h_{\hnu}\left(( \lambda \vec{u}\subset \vec{x}, \lambda \vec v\subseteq \vec y. \model{{\mu}}(\vec{u};\vec v))_{\mu \in \close{\hnu}},( \lambda \vec{u}\subseteq \vec{x}, \lambda \vec v\subseteq \vec y. \amodel{{\mu}}(\vec{u};\vec v))_{\mu \in \open{\hnu}}\right)(\vec{x}; \vec{y})
\end{equation*}


% Notice that in the following situations only point~\ref{enum:1} applies:
% \begin{itemize}
%     \item  $\nu$ is the conclusion of an instance of $\id$ or $\zero$;
%     \item  $\nu$ is the conclusion of an instance of $\wk_{\sq}$, $\sql$ and  $\cut_{\sn}$;
% \item $\nu$ is the conclusion of an instance of $\wk_{\n}$ and  $\der$ is a $\cbc$-coderivation.
% \end{itemize}
% In particular, the  last two cases hold by  Proposition~\ref{prop:structureofcycles}.\ref{enum:budcompanion1}-\ref{enum:budcompanion3}, as it must be that $\open{\nu'}=\emptyset$ for any premise $\nu'$ of $\hnu$. 

% Let us discuss the case where  $\hnu$ is the conclusion of a  $\cut_{\sq}$ rule with premises $\nu_1$ and $\nu_2$.  By induction on point~\ref{enum:1} we have $\model{\pder {\nu_1}}(\vec x; \vec y), \model{\pder{\nu_2}}(\vec x, x; \vec y)\in \nbcpp$. Since the conclusion of $\pder{\nu_1}$  has modal succedent, by Proposition~\ref{prop:cutbox}  there must be a coderivation $\der^*$ such that $\model{\der^*}(\vec x;)= \model{\pder{\nu_1}}(\vec x; \vec y)\in \nbcpp$. Moreover, by Proposition~\ref{prop:cutbox-for-left-leaning},  if $\pder{\nu_1}$ is a $\cbc$-coderivation then $\der^*$ is.  Hence, we define $\model{\pder \hnu}(\vec x; \vec y)= \model{\pder{\nu_2}}(\vec x, \model{\der^*}(\vec x;); \vec y)\in \nbcpp$. If moreover $\der$ is a $\cbc$-coderivation then, by applying the induction hypothesis, we obtain $\model{\pder{\hnu}}\in \bcpp$. 

% Let us now  consider   point~\ref{enum:2}. 
% If  $\hnu$ is the conclusion of a bud  then  $\open \hnu= \{ \hnu\}$, $\close \hnu=\emptyset$, and all points hold trivially.

% The cases where $\hnu$ is an instance of $\wk_\n$, $\exch_\n$, $\exch_\sq$, $\sqr$,  $\succ 0$ or $\succ 1$ are straightforward.  

Suppose that $\hnu$ is the conclusion of a $\cnd_\sq$ step with premises $\nu'$, $\nu_1$, and $\nu_2$. By Proposition~\ref{prop:structureofcycles}.\ref{enum:budcompanion1} we have $\open{\nu'}=\emptyset$, so that $\model{{\nu'}}\in \nbcpp$ (resp., $\model{{\nu'}}\in \bcpp$) by induction hypothesis.  By definition,  $\open \hnu=\open{\nu_1}\cup \open{\nu_2}$ and  $\close \hnu=\close{\nu_1}\cup \close{\nu_2}$. Then, we set:
\[
\arraycolsep=2pt
    \begin{array}{c}
    h_{ \hnu}(x, \vec{x}; \vec{y})
    = \cnd(;x,  \model{{\nu'}}(\vec{x}; \vec{y}), \model{{\nu_1}}(\pred(x;),\vec{x}; \vec{y}) , \model{{\nu_2}}(\pred(x;), \vec{x}; \vec{y}))
  \end{array}
  \]
  where $\pred(x;)$ can be defined from $\pred(;x)$ and projections.  By induction hypothesis on $\nu_i$:
  \[
  \arraycolsep=2pt
      \begin{array}{rcl}
       \model{{\nu_i} }(\pred(x;), \vec{x}; \vec{y})
       &=&
       h_{\nu_i}\left(( \lambda u,\vec{u}\subset \pred(x;),\vec{x}, \lambda \vec{v}. \model{{\mu}}(u,\vec{u};\vec{v}))_{\mu \in \close{\nu_i}}, ( \lambda u,\vec{u}\subseteq \pred(x;),\vec{x}, \lambda \vec{v}. \amodel{{\mu}}(u,\vec{u};\vec{v}))_{\mu \in\open{\nu_i}} \right)(\pred(x;),\vec{x}; \vec{y})\\
       &=&
       h_{\nu_i}\left(( \lambda u,\vec{u}\subset x,\vec{x}, \lambda \vec{v}. \model{{\mu}}(u,\vec{u};\vec{v}))_{\mu \in \close{\nu_i}}, ( \lambda u,\vec{u}\subset x,\vec{x}, \lambda \vec{v}. \amodel{{\mu}}(u,\vec{u};\vec{v}))_{\mu \in\open{\nu_i}} \right)(\pred(x;),\vec{x}; \vec{y})\\
      \end{array}
\]
  hence 
  \[
  \arraycolsep=2pt
  \begin{array}{c}
        \model{ \hnu}(x, \vec{x}; \vec{y})
        =
        h_{\hnu}\left(( \lambda u,\vec{u}\subset x,\vec{x}, \lambda \vec{v}. \model{{\mu}}(u,\vec{u};\vec{v}))_{\mu \in \close{\hnu}}, ( \lambda u,\vec{u}\subset x,\vec{x}, \lambda \vec{v}. \amodel{{\mu}}(u,\vec{u};\vec{v}))_{\mu \in\open{\hnu}} \right)(x,\vec{x}; \vec{y})\\
  \end{array}
  \]
  which shows in particular point~\ref{enum:b}. Points~\ref{enum:a} and~\ref{enum:d} are straightforward.

 
Let us now consider the case where $\hnu$ is an instance of $\cnd_\n$. By Proposition~\ref{prop:structureofcycles}.\ref{enum:budcompanion3} we have $\open{\nu'}=\emptyset$, so that $\model{{\nu'}}\in \nbcpp$ (resp., $\model{{\nu'}}\in \bcpp$) by induction hypothesis.  By definition,  $\open \hnu=\open{\nu_1}\cup \open{\nu_2}$ and  $\close \hnu=\close{\nu_1}\cup \close{\nu_2}$. Then, we set:
\[
\arraycolsep=2pt
\begin{array}{c}
 h_{ \hnu}(\vec{x};y,  \vec{y})
 = 
 \cnd(;y, \model{{\nu'}}(\vec{x}; \vec{y}), \model{{\nu_1}}(\vec{x};\pred(;y), \vec{y}) , \model{{\nu_2}}( \vec{x};\pred(;y), \vec{y}))    
\end{array}
\]
By induction hypothesis on $\nu_i$:
 \[
\arraycolsep=2pt
\begin{array}{rcl}
\model{{\nu_i} }(\vec{x}; \pred(;y), \vec{y}) 
&=&
h_{\nu_i}\left(( \lambda \vec{u}\subset x,\vec{x}, \lambda v,\vec{v}. \model{{\mu}}(\vec{u};v,\vec{v}))_{\mu \in \close{\nu_i}},( \lambda \vec{u}\subseteq x,\vec{x}, \lambda v,\vec{v}. \amodel{{\mu}}(\vec{u};v,\vec{v}))_{\mu \in  \open{\nu_i}}\right)(\vec{x};\pred(;y), \vec{y})
\end{array}
\]
and   hence 
   \[
\arraycolsep=2pt
\begin{array}{c}
\model{ \hnu}( \vec{x};y, \vec{y})
=
h_{\hnu}\left((\lambda \vec{u}\subset \vec{x}, \lambda v,\vec v. \model{{\mu}}(\vec{u};v,\vec v))_{\mu \in \close{\hnu}},(\lambda \vec{u}\subseteq \vec{x}, \lambda v,\vec v. \amodel{{\mu}}(\vec{u};v,\vec v))_{\mu \in \open{\hnu}}\right)(\vec{x};y, \vec{y})
\end{array}
\]
Points~\ref{enum:a} and~\ref{enum:b} are straightforward. Concerning point~\ref{enum:d} notice that by induction hypothesis on $\nu_i$ we have:
 \[
\arraycolsep=2pt
\begin{array}{rcl}
\model{{\nu_i} }(\vec{x}; \pred(;y), \vec{y}) 
&=&
h_{\nu_i}\left(( \lambda \vec{u}\subset x,\vec{x}, \lambda v,\vec{v}\subseteq \pred(;y), \vec{y}. \model{{\mu}}(\vec{u};v,\vec{v}))_{\mu \in \close{\nu_i}},( \lambda \vec{u}\subseteq x,\vec{x}, \lambda v,\vec{v}\subseteq \pred(;y), \vec{y}. \amodel{{\mu}}(\vec{u};v,\vec{v}))_{\mu \in  \open{\nu_i}}\right)(\vec{x};\pred(;y), \vec{y})
\end{array}
\]
and so we have:
 \[
\arraycolsep=2pt
\begin{array}{c}
\model{ \hnu}( \vec{x};y, \vec{y})
=
h_{\hnu}\left((\lambda \vec{u}\subset \vec{x}, \lambda v,\vec v\subset y, \vec{y}. \model{{\mu}}(\vec{u};v,\vec v))_{\mu \in \close{\hnu}},(\lambda \vec{u}\subseteq \vec{x}, \lambda v,\vec v\subset y, \vec{y}. \amodel{{\mu}}(\vec{u};v,\vec v))_{\mu \in \open{\hnu}}\right)(\vec{x};y, \vec{y})
\end{array}
\]


  
Let us finally consider the case where $\hnu$ is the conclusion of an instance of  $\com$  with premise $\nu'$, where  $X$ is the set of buds labelling the rule. We have  $\open{\hnu}= \open{\nu'} \setminus X$ and $\close{\hnu}=\close{\nu'}\cup\{\hnu \}$. We want to find $(h_\nu)_{\nu \in \close \nu \cup \{ \hnu\}}$ defining the equations for $(\model{ \nu})_{\nu \in \close \nu \cup \{ \hnu\}}$   in such a way that points~\ref{enum:a}-\ref{enum:d} hold.  First, note that, by definition of cycle nf, $\model{\pder{\hnu}}(\vec{x}; \vec{y}) =\model{\pder{\nu'}}(\vec{x}; \vec{y})=\model{\pder \mu}(\vec{x}; \vec{y})$ for all $\mu \in X$. Since $ \open{\nu'}=\open \hnu \cup X$ and  the path from $\nu'$ to any  $\mu\in X$ must cross an instance of $\cnd_{\sq}$ by  Proposition~\ref{prop:structureofcycles}.\ref{enum:budcompanion0}, by  induction hypothesis on $\nu'$ there exists a family  $(h'_{\nu})_{\nu \in \close{\nu'}\cup \{ \nu'\}}$ such that:
\begin{equation}\label{eqn:1}
\arraycolsep=2pt
\begin{array}{rcll}
     \model{{\nu'}}(\vec{x}; \vec{y})
     &=& 
      h'_{\nu'}\left(( \lambda \vec{u}\subset \vec{x}, \lambda \vec v. \model{{\mu}}(\vec{u};\vec{v}))_{\mu \in \close{\nu'}},( \lambda \vec{u}\subseteq \vec{x}, \lambda \vec v. \amodel{{\mu}}(\vec{u};\vec{v}))_{\mu \in  \open{\hnu}},( \lambda \vec{u}\subset \vec{x}, \lambda \vec v. \amodel{{\mu}}(\vec{u};\vec{v}))_{\mu \in  X}\right)(\vec{x}; \vec{y})
\end{array}
\end{equation}
and for all $  \nu \in \close{\nu'}$:
\begin{equation}\label{eqn:2}
\arraycolsep=2pt
\begin{array}{rcll}
    \model{{\nu}}(\vec{x}; \vec{y})
     &=& 
      h'_{\nu}\left(( \lambda \vec{u}\subset \vec{x}, \lambda \vec v. \model{{\mu}}(\vec{u};\vec{v}))_{\mu \in \close{\nu'}},( \lambda \vec{u}\subseteq \vec{x}, \lambda \vec v. \amodel{{\mu}}(\vec{u};\vec{v}))_{\mu \in  \open{\hnu}},( \lambda \vec{u}\subseteq \vec{x}, \lambda \vec v. \amodel{{\mu}}(\vec{u};\vec{v}))_{\mu \in  X}\right)(\vec{x}; \vec{y})
   \end{array}
\end{equation}
% \red{Since $ \open{\nu'}=\open \hnu \cup X$ and  the path from $\nu'$ to any  $\mu\in X$ must cross an instance of $\cnd_{\sq}$ by  Proposition~\ref{prop:structureofcycles}.\ref{enum:budcompanion0}, the induction hypothesis on $\nu'$ (point~\ref{enum:b}) allows us to rewrite ~\eqref{eqn:1}  as follows:
% \begin{equation}\label{eqn:translation1}
% \arraycolsep=2pt
%     \begin{array}{c}
%     \model{\pder{\nu'}}(\vec{x}; \vec{y})
%     = 
%     \begin{aligned}
%      g_{\nu'}(&( \lambda \vec{u}\subseteq \vec{x}, \lambda \vec v. \model{\pder \nu}(\vec{u};\vec v))_{\nu \in \close{\nu'}}, 
%       ( \lambda \vec{u}\subseteq \vec{x}, \lambda \vec v. \model{\pder \mu}(\vec{u};\vec v))_{\mu \in  \open \hnu},
%      ( \lambda \vec{u}\subset \vec{x}, \lambda \vec v. \model{\pder{\mu}}(\vec{u};\vec v))_{\mu \in  X})(\vec{x}; \vec{y}) 
%      \end{aligned}
%     \end{array}
% \end{equation}
% On the other hand, for all $\nu \in \mathcal{C}_{\nu'}$,  for all  $\vec u \subseteq \vec x$ and $\vec{v}$, the equation in~\eqref{eqn:2} can be rewritten as:
% \begin{equation}\label{eqn:translationintermediate}
% \arraycolsep=2pt
%     \begin{array}{c}
% \model{\pder \nu}(\vec{u};\vec v)
%      = 
%     \begin{aligned}
%       g_{\nu}( &(\lambda \vec{w}\subset \vec{u}, \lambda \vec {v'}. \model{\pder{\mu}}(\vec{u};\vec {v'}))_{\mu \in \close{\nu'}},    ( \lambda \vec{w}\subseteq \vec{u}, \lambda \vec {v'}. \model{\pder{\mu}}(\vec{w};\vec {v'}))_{\mu \in  \open{\hnu}},  ( \lambda \vec{w}\subseteq \vec{u}, \lambda \vec {v'}. \model{\pder{\mu}}(\vec{w};\vec {v'}))_{\mu \in  X}) (\vec{u}; \vec v)
%          \end{aligned}
%          \end{array}
% \end{equation}
% }
{we fix the oracles as $\amodel{\mu}\dfn \model{\hnu}$ for all $\mu \in X$} and set $h_{\hnu}\dfn h'_{\nu'}$ so that:
\begin{equation}\label{eqn:3}
\arraycolsep=2pt
\begin{array}{rcll}
     \model{{\hnu}}(\vec{x}; \vec{y})
     &=& 
      h_{\hnu}\left(( \lambda \vec{u}\subset \vec{x}, \lambda \vec v. \model{{\mu}}(\vec{u};\vec{v}))_{\mu \in \close{\nu'}},( \lambda \vec{u}\subseteq \vec{x}, \lambda \vec v. \amodel{{\mu}}(\vec{u};\vec{v}))_{\mu \in  \open{\hnu}}, \lambda \vec{u}\subset \vec{x}, \lambda \vec v. \model{{\hnu}}(\vec{u};\vec{v})\right)(\vec{x}; \vec{y})\\
      &=& 
      h_{\hnu}\left(( \lambda \vec{u}\subset \vec{x}, \lambda \vec v. \model{{\mu}}(\vec{u};\vec{v}))_{\mu \in \close{\nu'}\cup \{\hnu\}},( \lambda \vec{u}\subseteq \vec{x}, \lambda \vec v. \amodel{{\mu}}(\vec{u};\vec{v}))_{\mu \in  \open{\hnu}}\right)(\vec{x}; \vec{y})\\
       &=& 
      h_{\hnu}\left(( \lambda \vec{u}\subset \vec{x}, \lambda \vec v. \model{{\mu}}(\vec{u};\vec{v}))_{\mu \in \close{\hnu}},( \lambda \vec{u}\subseteq \vec{x}, \lambda \vec v. \amodel{{\mu}}(\vec{u};\vec{v}))_{\mu \in  \open{\hnu}}\right)(\vec{x}; \vec{y})
\end{array}
\end{equation}
and, for all $  \nu \in \close{\nu'}$:
\begin{equation}\label{eqn:4}
\arraycolsep=2pt
\begin{array}{rcll}
    \model{{\nu}}(\vec{x}; \vec{y})
     &=& 
      h'_{\nu}\left(( \lambda \vec{u}\subset \vec{x}, \lambda \vec v. \model{{\mu}}(\vec{u};\vec{v}))_{\mu \in \close{\nu'}},( \lambda \vec{u}\subseteq \vec{x}, \lambda \vec v. \amodel{{\mu}}(\vec{u};\vec{v}))_{\mu \in  \open{\hnu}}, \lambda \vec{u}\subseteq \vec{x}, \lambda \vec v. \model{{\hnu}}(\vec{u};\vec{v})\right)(\vec{x}; \vec{y})
   \end{array}
\end{equation}
Notice that $\vec{u}\subseteq \vec{x}$ might not be strict in $ \lambda \vec{u}\subseteq \vec{x}, \lambda \vec v. \model{{\hnu}}(\vec{u};\vec{v})$. However, by~\Cref{eqn:3} $f_{\hnu}$ can be rewritten to $  h_{\hnu}\left(( \lambda \vec{u}\subset \vec{x}, \lambda \vec v. \model{{\mu}}(\vec{u};\vec{v}))_{\mu \in \close{\hnu}},( \lambda \vec{u}\subseteq \vec{x}, \lambda \vec v. \amodel{{\mu}}(\vec{u};\vec{v}))_{\mu \in  \open{\hnu}}\right)$, and so $h'_\nu$ can be rewritten to a $h^*_\nu$ such that:
\begin{equation}\label{eqn:5}
\arraycolsep=2pt
\begin{array}{rcll}
    \model{{\nu}}(\vec{x}; \vec{y})
     &=& 
      h^*_{\nu}\left(( \lambda \vec{u}\subset \vec{x}, \lambda \vec v. \model{{\mu}}(\vec{u};\vec{v}))_{\mu \in \close{\hnu}},( \lambda \vec{u}\subseteq \vec{x}, \lambda \vec v. \amodel{{\mu}}(\vec{u};\vec{v}))_{\mu \in  \open{\hnu}}\right)(\vec{x}; \vec{y})
   \end{array}
\end{equation}
we set $h_\nu\dfn h^*_\nu$. Points~\ref{enum:a} and~\ref{enum:b} are straightforward.  Point~\ref{enum:d} follows  by noticing that the construction does not affect the safe arguments.
% \red{====================}
% and so, for all $\nu \in \close{\nu'}$:
% \begin{equation}\label{eqn:3}
% \arraycolsep=2pt
%     \begin{array}{c}
% % \model{\pder \nu}(\vec{u};\vec v)
%  \lambda \vec u \subseteq \vec x, \lambda v. \model{\pder \nu}(\vec{u};\vec v)
%     \\ = \\
%     \begin{aligned}
%      \lambda \vec u \subseteq \vec x, \lambda v.  g_{\nu}( &(\lambda \vec{w}\subset \vec{x}, \lambda \vec {v'}. \model{\pder{\mu}}(\vec{u};\vec {v'}))_{\mu \in \close{\nu'}},  ( \lambda \vec{w}\subseteq \vec{x}, \lambda \vec {v'}. \model{\pder{\mu}}(\vec{w};\vec {v'}))_{\mu \in  \open{\hnu}},  ( \lambda \vec{w}\subseteq \vec{x}, \lambda \vec {v'}. \model{\pder{\mu}}(\vec{w};\vec {v'}))_{\mu \in  X}) (\vec{u}; \vec v)
%          \end{aligned}
%          \end{array}
% \end{equation}
% Now, since  the paths from $\nu'$ to any  $\mu \in X$ in $\der$ must contain an instance of $\cnd_{\sq}$, for all  $\nu \in \close{\nu'}$ and all  $\mu \in X$, we have that either the path from $\nu'$ to $\nu$  contains an instance of $\cnd_{\sq}$ or the path from $\nu$ to $\mu$ does. By applying the induction hypothesis on $\nu'$  (point~\ref{enum:b}), given  $\nu \in \close{\nu'}$  and   $\mu \in X$, either  $\lambda \vec{u}\subseteq \vec{x}, \lambda \vec v. \model{\pder{\nu}}(\vec{u};\vec v)$ in~\eqref{eqn:1} is such that $\vec{u}\subset \vec{x}$,   or $\lambda \vec{w}\subseteq \vec{u}, \lambda \vec {v'}. \model{\pder{\mu}}(\vec{w};\vec {v'})$ in~\eqref{eqn:2} is such that $\vec{w}\subset \vec{u}$. This means that, for any $\mu \in X$,   $\lambda \vec{w}\subseteq \vec{x}, \lambda \vec {v'}. \model{\pder{\mu}}(\vec{w};\vec {v'})$ in~\eqref{eqn:3} is such that $\vec{w}\subset \vec{x}$. For each  $\nu \in \close{\nu'}$, by rewriting $ \lambda \vec{u}\subseteq  \vec{x}, \lambda \vec v. \model{\pder{\nu}}(\vec{u};\vec v)$  in~\eqref{eqn:1} according to the equation in~\eqref{eqn:3} we obtain:
% \[
% \arraycolsep=2pt
% \begin{array}{c}
%    \model{\pder{\hnu}}(\vec{x}; \vec{y})
%    =
%    \begin{aligned}
%    t_{\hnu}( &( \lambda \vec{u}\subset \vec{x}, \lambda \vec v. \model{\pder{\mu}}(\vec{u};\vec v))_{\mu \in  \close{\nu'}\cup X}, ( \lambda \vec{u}\subseteq \vec{x}, \lambda \vec v. \model{\pder{\mu}}(\vec{u};\vec v))_{\mu \in  \open{\hnu}})(\vec{x}; \vec{y})
%    \end{aligned}
% \end{array}
% \]
% for some $t_{\hnu}$. Since  $\model{\pder{\mu}}= \model{\pder{\hnu}}$ for all $\mu \in X$,   and since  $\close{\hnu}= \close{\nu'}\cup \{ \hnu\}$, by setting $h_{\hnu}\dfn t_{\hnu}$ the above equation gives us the following:
% \begin{equation}\label{eqn:translationintermediate2}
% \arraycolsep=2pt
% \begin{array}{c}
%     \model{\pder{\hnu}}(\vec{x}; \vec{y})
%     =
%     h_{\hnu}(( \lambda \vec{u}\subseteq \vec{x}, \lambda \vec v. \model{\pder{\mu}}(\vec{u};\vec v))_{\mu \in \close{\hnu}\cup \open{\hnu}})(\vec{x}; \vec{y})
%     \end{array}
% \end{equation}
% which satisfies point~\ref{enum:b}.   From~\eqref{eqn:translationintermediate2} we are able to find the functions $(h_\nu)_{\nu \in \close \nu}$ defining the equations for $(\model{\pder \nu})_{\nu \in \close \nu}$. Indeed, the induction hypothesis on $\nu'$ gives us~\eqref{eqn:2} for any $\nu \in \close{\nu'}$. We rewrite in each such equation any   $\lambda \vec{u}\subseteq \vec{x}, \lambda \vec{v}.
% \model{\pder{\mu}}(\vec{u};\vec v)$  such that  $\mu \in X$ according to equation~\eqref{eqn:translationintermediate2}, as $\model{\pder{\mu}}=\model{\pder{\hnu}}$ for any $\mu\in X$. We obtain the following equation for any $\nu \in \close{\nu'}$:
% \[
% \arraycolsep=2pt
% \begin{array}{c}
% \model{\pder{\nu}}(\vec{x}; \vec{y})
% =
%    \begin{aligned}
%      t_\nu( &( \lambda \vec{u}\subset \vec{x}, \lambda \vec v. \model{\pder{\mu}}(\vec{u};\vec v))_{\mu \in  \close{\nu'}\cup \{\hnu\}}, ( \lambda \vec{u}\subseteq \vec{x}, \lambda \vec v. \model{\pder{\mu}}(\vec{u};\vec v))_{\mu \in  \open{\hnu}})(\vec{x}; \vec{y})
%    \end{aligned}
% \end{array}
% \]
% for some $t_\nu$. Since $\close{\hnu}=\close{\nu'}\cup \{\hnu\}$ and the above  equation satisfies point~\ref{enum:b}, we set $h_{\nu}\dfn t_{\nu}$ and we obtain, for all $\nu \in \close{\hnu}$:
% \begin{equation}\label{eqn:4}
%    \model{\pder{\nu}}(\vec{x}; \vec{y})= h_\nu( ( \lambda \vec{u}\subseteq \vec{x}, \lambda \vec v. \model{\pder{\mu}}(\vec{u};\vec v))_{\mu \in  \close{\hnu}\cup \open{\hnu}})(\vec{x}; \vec{y})
% \end{equation}
% It remains to show that~\eqref{eqn:translationintermediate2} and~\eqref{eqn:4} satisfy points~\ref{enum:a} and~\ref{enum:d}. Concerning point~\ref{enum:a},   on the one hand for all $\nu \in \close{\hnu}$ we  have  $h_{\nu}\in \nbcpp((\model{\pder{\mu}})_{\mu \in \open{\hnu}}, (a_\mu)_{\mu \in \close{\hnu}} )$, with   $(a_\mu)_{\mu \in \close{\hnu}} $  oracle functions.  On the other hand, by applying the induction hypothesis, we have  $\model{\pder{\hnu}}=\model{\pder{\nu'}}\in \nbcpp((\model{\pder\mu})_{\mu \in \open{\nu'}})$ and  $\model{\pder{\nu}}\in \nbcpp((\model{\pder\mu})_{\mu \in \open{\nu'}})$, for all $\nu \in \close{\hnu}$.  Since $\open{\nu'}= \open{\hnu}\cup X$ and  $\model{\pder{\hnu}}= \model{\pder{\mu}}$ for all $\mu \in X$, we have both  $\model{\pder{\hnu}}\in \nbcpp((\model{\pder\mu})_{\mu \in \open{\hnu}})$ and   $\model{\pder{\nu}}\in \nbcpp((\model{\pder\mu})_{\mu \in \open{\hnu}},  \model{\pder\hnu})$ for all $\nu \in \close{\hnu}$, and hence $\model{\pder{\nu}}\in \nbcpp((\model{\pder\mu})_{\mu \in \open{\hnu}})$, for all $\nu \in \close{\hnu}$. Point~\ref{enum:d} follows  by applying the induction hypothesis, as the construction does not affect the safe arguments.



\end{proof}




Finally, we can establish the main result of this paper:
\begin{cor}\label{cor:main-theorem} 
We have the following:
\begin{itemize}
    \item $f(\vec x; )\in \cbc$ if and only if $f(\vec x)\in \fptime$;
    \item $f(\vec x; )\in \ncbc$ if and only if $f(\vec x)\in \felementary$.
\end{itemize}
\end{cor}
\begin{proof}
[Proof sketch]
Soundness ($\Rightarrow$) follows from \Cref{lem:short-translation-lemma}, by showing that $\bcpp$ and $\nbcpp$ are closed under \emph{simultaneous} versions of their recursion schemes.
Completeness ($\Leftarrow$) follows  from~\Cref{thm:bc-in-cbc} and \Cref{thm:nbc-in-ncbc}. 
\end{proof}
\section{Conclusions and further remarks}\label{sec:conclusion}


In this work we presented two-tiered circular type systems $\cbc$ and $\ncbc$ and showed that they capture polynomial-time and elementary computation, respectively.
This is the first time that methods of circular proof theory have been applied in implicit computational complexity (ICC).
Along the way we gave novel relativised algebras for these classes based on safe (nested) recursion on well-founded relations.

Since the conference version \cite{CurziDas:lics22} of this work was published, other works in Cyclic Implicit Complexity have appeared, building on the present work, in particular in the setting of \emph{non-uniform computation} \cite{CurziDas:csl23,AcclavioCurziGuerrieri:csl24}.

\subsection{Unary notation and linear space}
It is well-known that  $\flinspace$, i.e.~the class of functions computable in linear space, can be captured  by reformulating $\bc$ in  \emph{unary} notation (see~\cite{bellantoniTESIS}). A similar result can be obtained for  $\cbc$ by just defining a unary version of the conditional in $\bcpp$ (similarly to the ones in \cite{Das2021,Kuperberg-Pous21}) and by adapting the proofs  of~\Cref{lem:boundinglemma},~\Cref{thm:bc-in-cbc}  and~\Cref{lem:short-translation-lemma}.  On the other hand, $\ncbc$ is (unsurprisingly) not sensitive to such choice of notation. 


 \subsection{On unnested recursion with compositions}
  We believe that \Cref{lem:boundinglemma} can be adapted to establish a polynomial bound on the growth rate  for the function algebra extending $\bcnorec$ by `Safe Composition During Safe Recursion', cf.~\Cref{rem:errata}. 
  We conjecture that this  function algebra and its extension to recursion on permutation of prefixes (cf.~\Cref{defn:saferec-pp}) capture precisely the class $\fpspace$. 
  Indeed, several function algebras for $\fpspace$ have been proposed in the literature, many of which involve variants of~\eqref{eq:safe-rec-w-comp-dur-rec-no-oracles} (see~\cite{Leivant-pspace,Oitavem-pspace}).
%   As already mentioned in Remark~\ref{rem:characterization-of-fpspace}, the function algebras capturing this class (see e.g.~\cite{Leivant-pspace,Oitavem-pspace}) typically require safe recursion schemes combining multiple recursive calls and composition during recursion. 
% 
  These recursion schemes reflect the parallel nature of polynomial space functions, which can be defined in terms of alternating polynomial time computation. 
  We suspect that a circular proof theoretic characterisation of this class can thus be achieved by extending $\cbc$ with a `parallel' version of the  cut rule and by adapting the left-leaning criterion appropriately.
Parallel cuts might also play a fundamental role for potential circular proof theoretic characterisations of circuit complexity classes, like $\alogtime$ or $\nc$. 


% \red{
% Note that Safe Nested Recursion also admits definitions that are not morally `nested' but rather use a form of `composition during recursion':
% \begin{itemize}
%     \item from $g(\vec x; \vec y), \left(\vec g_j(x, \vec x; \vec y)\right)_{1 \leq j \leq k}$ and  $h_i\left(x, \vec x; \vec y, (z_j)_{1 \leq j \leq k}\right)$  define:
%     \begin{equation}
%         \label{eq:safe-rec-w-comp-dur-rec-no-oracles}
%          \arraycolsep=2pt
%     \begin{array}{rcll}
%          f(0, \vec x; \vec y) &=&g(\vec x;\vec y)\\
%          f(\succ i x , \vec x; \vec y) &=& 
%     h_i \left(x, \vec x; \vec y, \left(f(x, \vec x; \vec g_j(x, \vec x; \vec y))\right)_{1 \leq j \leq k}\right)
%         \end{array}
%     \end{equation}
% \end{itemize}
% In this case, note that we have allowed the safe inputs of $f$ to take arbitrary values given by previously defined functions, but at the same time $f$ never calls itself in a safe position, as in~\eqref{eq:safe-exp-by-snrec}.
% This safe recursion scheme does not require the use of oracles in its statement, but is equivalent to the following one using a notion of `nesting' 
% for Safe Composition:
% \begin{defn}
% [Safe Recursion with Composition During Recursion]
% \label{defn:safe-rec-with-comp-dur-rec}
% An instance of Safe Composition along Safe Parameters  (\Cref{enum:safe-composition-condition-1} from \Cref{defn:oracles-functions}) is called \emph{unnested} if (at least) one of $g(\vec x;\vec y)$ and $h(\vec x; \vec y, y)$ is over $\emptyset$ (i.e.\ is oracle-free). 
% % 
% We define $\sbc(\vec a)$ to be the restriction of $\nbc(\vec a)$ using only unnested Safe Composition along Safe Parameters, 
% % \Cref{enum:safe-composition-condition-1} from \Cref{defn:oracles-functions}, 
% and write simply $\sbc$ for $\sbc(\emptyset)$.
% \end{defn}

% It is not hard to see that, indeed,   $\sbc$ defines the same class of functions of  the extension of $\bcnorec$ by the scheme given in \Cref{eq:safe-rec-w-comp-dur-rec-no-oracles}.
% The `$\mathsf S$' in $\sbc$ thus indicates that the recursion scheme is `shallow', as opposed to nested.
% % \anupam{remainder of paragraph can be omitted} 
% %  On the one hand, given $h_i$ and $\vec{g}_j$, the scheme in~\eqref{eq:safe-rec-w-comp-dur-rec-no-oracles} can be defined in $\nbc(\vec a)$ using $\snrec$  and  only unnested instances of safe composition. On the other hand, given $h(a)(x,\vec x;\vec y)$ over oracles $a,\vec a$, the scheme $\snrec$ allows us to define  $f(\succ i x, \vec x; \vec y)$ by replacing the oracle $a$ in $h$ with the recursive calls $\lambda \vec v . f(x,\vec x;\vec v)$. However, since $h(a)$  is defined using unnested safe composition, no recursive call of $f$ can be found along the safe parameters of another one.
% % 
% It should be said that we will not actually use $\sbc$ in this work, but rather an extension of it defined in the next subsection, but we have included it for completeness of the exposition.
% }

\subsection{Towards higher-order cyclic implicit complexity}
It would be pertinent to pursue higher-order versions of both $\ncbc$ and $\cbc$, in light of precursory works \cite{Das2021,Kuperberg-Pous21} in circular proof theory as well as ICC {\cite{Hofmann97,Leivant99,BellantoniNS00}}.
 In the case of polynomial-time, for instance,
 a soundness result for some higher-order version of $\cbc$ might follow by translation to (a sequent-style formulation of) Hofmann's $\slr$~\cite{Hofmann97}. 
 Analogous translations might be defined for a higher-order version of $\ncbc$ once the linearity restrictions on the recursion operator of $\slr$ are dropped. 
 Finally, as $\slr$ is essentially a subsystem of G\"{o}del's system $\tgodel$, such translations could refine the results on the  abstraction complexity (i.e.~type level) of the circular version of system $\tgodel$  in~\cite{Das2021,Das2021-preprint}. 
 
\paragraph{Acknowledgements}

We would like to thank {the anonymous reviewers for their valuable comments and suggestions. We are also indebted to} Dominik Wehr for pointing out a minor error in the conference version of this paper~\cite{CurziDas:lics22} (see~\Cref{rem:errata}).
This work was supported by a UKRI Future Leaders Fellowship, ‘Structure vs Invariants in Proofs’ (project reference MR/S035540/1),   by the Wallenberg Academy Fellowship Prolongation 
project `Taming Jörmungandr: The Logical Foundations of Circularity' (project reference 251080003), and by the
VR starting grant “Proofs with Cycles in Computation” (project reference 251088801).


% The widespread approach to ICC is based on the introduction of inductively defined languages or calculi endowed with recursion mechanisms whose strength is carefully calibrated in order to increase in complexity  while not overstepping a given bound on computation. The circular proof systems $\cbc$ and $\ncbc$ pave the way to a radically different, top-down approach, where coinductive reasoning  plays a central role. 
% % Regular proofs are able to subsume various forms of recursion, and in some cases  even  non-terminating computation. Therefore, the characterisation results are achieved by imposing global proof-theoretical conditions weakening the computational content of the system.
% % \anupam{I commented two sentences here since they seemed out of place. Can be re-added if developed or put somewhere else.}

% We conclude this paper with some further remarks, some comments about expressivity, and some avenues for future work.



% \anupam{shorten all this to three paragraphs or so}

% \toappendix{\subsection{Further remarks}

% It is well-known that  $\flinspace$, i.e.~the class of functions computable in linear space, can be captured  by reformulating $\bc$ in  \emph{unary} notation (see~\cite{bellantoniTESIS}). A similar result can be obtained for  $\cbc$ by just defining a unary version of the conditional in $\bcpp$ (similarly to the ones in \cite{Das2021,KPP21}) and by adapting the proofs  of~\Cref{lem:boundinglemma},~\Cref{thm:bc-in-cbc}  and~\Cref{lem:translation}.  On the other hand, $\ncbc$ is (unsurprisingly) not sensitive to such choice of notation. 


% Given the deep connection between the notion of regularity for coderivations and the notion of uniformity in computation discussed in Subsection~\ref{subsec:definability in typing derivations.}, we may wonder if dropping the regularity condition from e.g.~$\cbc$ would result in a characterisation of the non-uniform version of $\ptime$, e.g.~$\mathbf{P}/\mathbf{poly}$. In other words, can we capture a non-uniform complexity class using the same proof-theoretic principles  adopted to capture its uniform  formulation? Unfortunately, the answer of this question is negative. To see this, notice that~\Cref{exmp:extensional-completeness-type-1} shows how to define \emph{any}  number-theoretic function by a non-regular progressing $\bcnorec$-coderivation, which happens to be both safe and left-leaning in a trivial way. 
% It would nonetheless be interesting if the regular/non-regular distinction in circular proof theory could somehow be leveraged to model the uniform/non-uniform distinction in computational complexity.


% Let us finally observe that our proof-theoretic approach to `circular complexity' was motivated by the inherently syntactical nature of threads, which define the  progressing criterion, cf.~\Cref{defn:progressing}.  
% Surprisingly enough, however, it seems that a purely  term-oriented presentation of~\Cref{cor:main-theorem} could be possible thanks to~\Cref{prop:regular-safe-progressing-iff-infinite-branch-infinite-cond}, which shows that $\cbc$ and $\ncbc$ can be defined \emph{independently} on the notion of thread. Consequently, our circular  proof systems might be  directly    presented  as  `function coalgebras' satisfying certain conditions or, alternatively, {as `two-sorted'  Kleene-Herbrand-G\"{o}del equational programs (as suggested by~\Cref{defn:semantics-coderivations})}. 
% Such a presentation could be fruitful for future endeavours in this area.


% \subsection{On intensional expressivity of circular systems}
% Notice that~\Cref{cor:main-theorem}  has an \emph{extensional} flavour, as it refers only to usual function complexity classes. 
% What about the \emph{intensional} power of our circular proof systems? 
% Can we gain some algorithmic expressivity by moving from a function algebra to its circular formulation? 
% At least one way of comparing the intensional power of two-tiered systems in ICC is to compare their functions with both normal and safe inputs, not only normal inputs as is usually the case, cf.~\Cref{cor:main-theorem}. 
% % of safe-normal functions is to compare their function , say  $\mathsf{F}$ and $\mathsf{G}$, is the following: we say that  $\mathsf{F}$ is `smaller' than  $\mathsf{G}$ whenever any function $f$ that is definable in $\mathsf{F}$ is also definable in $\mathsf{G}$ with the \emph{same sorting},  and we write  $\mathsf{F} \subseteq \mathsf{G}$ with little abuse of notation.  

% To this end, let us point out that the relations in Figure~\ref{fig:picture-main-results} between two-tiered systems indeed respect the safe-normal distinction of inputs. I.e.\ the  inclusions $\nbc \subseteq \ncbc \subseteq \nbcpp$ and $\bc \subseteq \cbc \subseteq \bcpp$ indeed hold when construing each system as a set of safe-normal functions.
% In this way, note that showing   $\nbcpp\subseteq \nbc $ or $\bcpp \subseteq \bc$ would be enough to infer that all systems have equivalent intensional expressivity, at the level of their definable safe-normal functions.
% We suspect that such a collapse does not hold, but we should point out that existing techniques for showing that a safe-normal function is not definable in a two-tiered system usually exploit a version of the Bounding Lemma from \cite{BellantoniCook} (cf.~\Cref{eq:bc-polymax-eqn}), which all the systems considered here satisfy (cf.~\Cref{lem:boundinglemma}). 

% \subsection{On the `power' of contraction}
% Revisiting~\Cref{rem:explicit-contraction}, it not clear  whether~\Cref{cor:main-theorem} can be restated in presence of the contraction rule  $\cntr_\sq$. Were this possible, $\mathsf{NL}$-checkability of both  $\cbc$ and $\ncbc$ would still hold as a consequence of~\Cref{rem:NL-checkability-with-ecplicit-contraction}. 
% % We conjecture that contraction does not affect the complexity of our circular proof systems. 
% % \anupam{do we?! This seems like a strong thing to say!}
% Nonetheless,  establishing soundness by means of a direct translation would be a rather difficult task, as the rule $\cntr_\sq$ introduces more involved thread structures which invalidate some invariants of~\Cref{lem:translation}. 
% Indeed in a setting where cut is context-splitting rather than context-sharing, it is known that contraction strictly increases extensional expressivity of regular progressing coderivations \cite{KPP21}.



% \subsection{Future work}
%  It would be pertinent to pursue higher-order versions of both $\ncbc$ and $\cbc$, in light of precursory works in circular proof theory \cite{Das2021,KPP21} as well as ICC {\cite{Hofmann97,Leivant99,BellantoniNS00}}.
%  In the case of polynomial-time, for instance,
%  a soundness result for some higher-order version of $\cbc$ might follow by translation to (a sequent-style formulation of) Hofmann's $\slr$~\cite{Hofmann97}. 
%  Analogous translations might be defined for a higher-order version of $\ncbc$ once the linearity restrictions on the recursion operator of $\slr$ are dropped. 
%  Finally, as $\slr$ is essentially a subsystem of G\"{odel}'s system $\tgodel$, such translations could refine the results on the  abstraction complexity (i.e.~type level) of the circular version of system $\tgodel$ presented  in~\cite{Das2021}. 
 
%   We are also investigating how to extend the results of this paper to other relevant complexity classes, like  $\fpspace$. 
%   On the function algebraic side, notice that~\Cref{lem:boundinglemma} implies a polyspace bound for functions in $\sbcpp$, and hence for functions in $\sbc$. We conjecture that both function algebras capture precisely the class $\fpspace$. Indeed, as already observed, the unnnested version of the recursion scheme $\snrec$ can be replaced by the scheme in~\eqref{eq:safe-rec-w-comp-dur-rec-no-oracles}, which allows both  multiple recursive calls and composition during recursion. Several function algebras for $\fpspace$ have been proposed in the literature, and all of them involve variants of~\eqref{eq:safe-rec-w-comp-dur-rec-no-oracles} (see~\cite{Leivant-pspace,Oitavem-pspace}).
% %   As already mentioned in Remark~\ref{rem:characterization-of-fpspace}, the function algebras capturing this class (see e.g.~\cite{Leivant-pspace,Oitavem-pspace}) typically require safe recursion schemes combining multiple recursive calls and composition during recursion. 

%   These recursion schemes reflect the parallel nature of polynomial space functions, which in fact can be defined in terms of alternating polynomial time computation. 
%   We suspect that a circular proof theoretic characterisation of this class can be achieved by extending $\cbc$ with the following parallel version of the  cut rule:
% \[
%     \vliiiinf{\pcut}{}{\Gamma \seqar \n}{\Gamma \seqar \n}{\cdots}{\Gamma \seqar \n}{\Gamma, \n, \dots, \n \seqar \n}
% \]
% and by adapting the left-leaning criterion in such a way as to allow some forms of composition during recursion. 
% Indeed, we believe that $\sbc$ itself characterises $\fpspace$, but this result is beyond the scope of this work.

% Parallel cut might also play a fundamental role for potential circular proof theoretic characterisations of circuit complexity classes, like $\alogtime$ or $\nc$. 
% In such a setting, 
% % the  type for natural numbers would be replaced by the ground type  for binary strings, 
% one would expect a conditional rule that implements a divide and conquer searching mechanism {(see, e.g.,~\cite{Leivant98,LeivantM00,BonfanteKMO16})}.
% }
   








% \section{Introduction}
% ACM's consolidated article template, introduced in 2017, provides a
% consistent \LaTeX\ style for use across ACM publications, and
% incorporates accessibility and metadata-extraction functionality
% necessary for future Digital Library endeavors. Numerous ACM and
% SIG-specific \LaTeX\ templates have been examined, and their unique
% features incorporated into this single new template.

% If you are new to publishing with ACM, this document is a valuable
% guide to the process of preparing your work for publication. If you
% have published with ACM before, this document provides insight and
% instruction into more recent changes to the article template.

% The ``\verb|acmart|'' document class can be used to prepare articles
% for any ACM publication --- conference or journal, and for any stage
% of publication, from review to final ``camera-ready'' copy, to the
% author's own version, with {\itshape very} few changes to the source.

% \section{Template Overview}
% As noted in the introduction, the ``\verb|acmart|'' document class can
% be used to prepare many different kinds of documentation --- a
% dual-anonymous initial submission of a full-length technical paper, a
% two-page SIGGRAPH Emerging Technologies abstract, a ``camera-ready''
% journal article, a SIGCHI Extended Abstract, and more --- all by
% selecting the appropriate {\itshape template style} and {\itshape
%   template parameters}.

% This document will explain the major features of the document
% class. For further information, the {\itshape \LaTeX\ User's Guide} is
% available from
% \url{https://www.acm.org/publications/proceedings-template}.

% \subsection{Template Styles}

% The primary parameter given to the ``\verb|acmart|'' document class is
% the {\itshape template style} which corresponds to the kind of publication
% or SIG publishing the work. This parameter is enclosed in square
% brackets and is a part of the {\verb|documentclass|} command:
% \begin{verbatim}
%   \documentclass[STYLE]{acmart}
% \end{verbatim}

% Journals use one of three template styles. All but three ACM journals
% use the {\verb|acmsmall|} template style:
% \begin{itemize}
% \item {\verb|acmsmall|}: The default journal template style.
% \item {\verb|acmlarge|}: Used by JOCCH and TAP.
% \item {\verb|acmtog|}: Used by TOG.
% \end{itemize}

% The majority of conference proceedings documentation will use the {\verb|acmconf|} template style.
% \begin{itemize}
% \item {\verb|acmconf|}: The default proceedings template style.
% \item{\verb|sigchi|}: Used for SIGCHI conference articles.
% \item{\verb|sigchi-a|}: Used for SIGCHI ``Extended Abstract'' articles.
% \item{\verb|sigplan|}: Used for SIGPLAN conference articles.
% \end{itemize}

% \subsection{Template Parameters}

% In addition to specifying the {\itshape template style} to be used in
% formatting your work, there are a number of {\itshape template parameters}
% which modify some part of the applied template style. A complete list
% of these parameters can be found in the {\itshape \LaTeX\ User's Guide.}

% Frequently-used parameters, or combinations of parameters, include:
% \begin{itemize}
% \item {\verb|anonymous,review|}: Suitable for a ``dual-anonymous''
%   conference submission. Anonymizes the work and includes line
%   numbers. Use with the \verb|\acmSubmissionID| command to print the
%   submission's unique ID on each page of the work.
% \item{\verb|authorversion|}: Produces a version of the work suitable
%   for posting by the author.
% \item{\verb|screen|}: Produces colored hyperlinks.
% \end{itemize}

% This document uses the following string as the first command in the
% source file:
% \begin{verbatim}
% \documentclass[sigconf,authordraft]{acmart}
% \end{verbatim}

% \section{Modifications}

% Modifying the template --- including but not limited to: adjusting
% margins, typeface sizes, line spacing, paragraph and list definitions,
% and the use of the \verb|\vspace| command to manually adjust the
% vertical spacing between elements of your work --- is not allowed.

% {\bfseries Your document will be returned to you for revision if
%   modifications are discovered.}

% \section{Typefaces}

% The ``\verb|acmart|'' document class requires the use of the
% ``Libertine'' typeface family. Your \TeX\ installation should include
% this set of packages. Please do not substitute other typefaces. The
% ``\verb|lmodern|'' and ``\verb|ltimes|'' packages should not be used,
% as they will override the built-in typeface families.

% \section{Title Information}

% The title of your work should use capital letters appropriately -
% \url{https://capitalizemytitle.com/} has useful rules for
% capitalization. Use the {\verb|title|} command to define the title of
% your work. If your work has a subtitle, define it with the
% {\verb|subtitle|} command.  Do not insert line breaks in your title.

% If your title is lengthy, you must define a short version to be used
% in the page headers, to prevent overlapping text. The \verb|title|
% command has a ``short title'' parameter:
% \begin{verbatim}
%   \title[short title]{full title}
% \end{verbatim}

% \section{Authors and Affiliations}

% Each author must be defined separately for accurate metadata
% identification. Multiple authors may share one affiliation. Authors'
% names should not be abbreviated; use full first names wherever
% possible. Include authors' e-mail addresses whenever possible.

% Grouping authors' names or e-mail addresses, or providing an ``e-mail
% alias,'' as shown below, is not acceptable:
% \begin{verbatim}
%   \author{Brooke Aster, David Mehldau}
%   \email{dave,judy,steve@university.edu}
%   \email{firstname.lastname@phillips.org}
% \end{verbatim}

% The \verb|authornote| and \verb|authornotemark| commands allow a note
% to apply to multiple authors --- for example, if the first two authors
% of an article contributed equally to the work.

% If your author list is lengthy, you must define a shortened version of
% the list of authors to be used in the page headers, to prevent
% overlapping text. The following command should be placed just after
% the last \verb|\author{}| definition:
% \begin{verbatim}
%   \renewcommand{\shortauthors}{McCartney, et al.}
% \end{verbatim}
% Omitting this command will force the use of a concatenated list of all
% of the authors' names, which may result in overlapping text in the
% page headers.

% The article template's documentation, available at
% \url{https://www.acm.org/publications/proceedings-template}, has a
% complete explanation of these commands and tips for their effective
% use.

% Note that authors' addresses are mandatory for journal articles.

% \section{Rights Information}

% Authors of any work published by ACM will need to complete a rights
% form. Depending on the kind of work, and the rights management choice
% made by the author, this may be copyright transfer, permission,
% license, or an OA (open access) agreement.

% Regardless of the rights management choice, the author will receive a
% copy of the completed rights form once it has been submitted. This
% form contains \LaTeX\ commands that must be copied into the source
% document. When the document source is compiled, these commands and
% their parameters add formatted text to several areas of the final
% document:
% \begin{itemize}
% \item the ``ACM Reference Format'' text on the first page.
% \item the ``rights management'' text on the first page.
% \item the conference information in the page header(s).
% \end{itemize}

% Rights information is unique to the work; if you are preparing several
% works for an event, make sure to use the correct set of commands with
% each of the works.

% The ACM Reference Format text is required for all articles over one
% page in length, and is optional for one-page articles (abstracts).

% \section{CCS Concepts and User-Defined Keywords}

% Two elements of the ``acmart'' document class provide powerful
% taxonomic tools for you to help readers find your work in an online
% search.

% The ACM Computing Classification System ---
% \url{https://www.acm.org/publications/class-2012} --- is a set of
% classifiers and concepts that describe the computing
% discipline. Authors can select entries from this classification
% system, via \url{https://dl.acm.org/ccs/ccs.cfm}, and generate the
% commands to be included in the \LaTeX\ source.

% User-defined keywords are a comma-separated list of words and phrases
% of the authors' choosing, providing a more flexible way of describing
% the research being presented.

% CCS concepts and user-defined keywords are required for for all
% articles over two pages in length, and are optional for one- and
% two-page articles (or abstracts).

% \section{Sectioning Commands}

% Your work should use standard \LaTeX\ sectioning commands:
% \verb|section|, \verb|subsection|, \verb|subsubsection|, and
% \verb|paragraph|. They should be numbered; do not remove the numbering
% from the commands.

% Simulating a sectioning command by setting the first word or words of
% a paragraph in boldface or italicized text is {\bfseries not allowed.}

% \section{Tables}

% The ``\verb|acmart|'' document class includes the ``\verb|booktabs|''
% package --- \url{https://ctan.org/pkg/booktabs} --- for preparing
% high-quality tables.

% Table captions are placed {\itshape above} the table.

% Because tables cannot be split across pages, the best placement for
% them is typically the top of the page nearest their initial cite.  To
% ensure this proper ``floating'' placement of tables, use the
% environment \textbf{table} to enclose the table's contents and the
% table caption.  The contents of the table itself must go in the
% \textbf{tabular} environment, to be aligned properly in rows and
% columns, with the desired horizontal and vertical rules.  Again,
% detailed instructions on \textbf{tabular} material are found in the
% \textit{\LaTeX\ User's Guide}.

% Immediately following this sentence is the point at which
% Table~\ref{tab:freq} is included in the input file; compare the
% placement of the table here with the table in the printed output of
% this document.

% \begin{table}
%   \caption{Frequency of Special Characters}
%   \label{tab:freq}
%   \begin{tabular}{ccl}
%     \toprule
%     Non-English or Math&Frequency&Comments\\
%     \midrule
%     \O & 1 in 1,000& For Swedish names\\
%     $\pi$ & 1 in 5& Common in math\\
%     \$ & 4 in 5 & Used in business\\
%     $\Psi^2_1$ & 1 in 40,000& Unexplained usage\\
%   \bottomrule
% \end{tabular}
% \end{table}

% To set a wider table, which takes up the whole width of the page's
% live area, use the environment \textbf{table*} to enclose the table's
% contents and the table caption.  As with a single-column table, this
% wide table will ``float'' to a location deemed more
% desirable. Immediately following this sentence is the point at which
% Table~\ref{tab:commands} is included in the input file; again, it is
% instructive to compare the placement of the table here with the table
% in the printed output of this document.

% \begin{table*}
%   \caption{Some Typical Commands}
%   \label{tab:commands}
%   \begin{tabular}{ccl}
%     \toprule
%     Command &A Number & Comments\\
%     \midrule
%     \texttt{{\char'134}author} & 100& Author \\
%     \texttt{{\char'134}table}& 300 & For tables\\
%     \texttt{{\char'134}table*}& 400& For wider tables\\
%     \bottomrule
%   \end{tabular}
% \end{table*}

% Always use midrule to separate table header rows from data rows, and
% use it only for this purpose. This enables assistive technologies to
% recognise table headers and support their users in navigating tables
% more easily.

% \section{Math Equations}
% You may want to display math equations in three distinct styles:
% inline, numbered or non-numbered display.  Each of the three are
% discussed in the next sections.

% \subsection{Inline (In-text) Equations}
% A formula that appears in the running text is called an inline or
% in-text formula.  It is produced by the \textbf{math} environment,
% which can be invoked with the usual
% \texttt{{\char'134}begin\,\ldots{\char'134}end} construction or with
% the short form \texttt{\$\,\ldots\$}. You can use any of the symbols
% and structures, from $\alpha$ to $\omega$, available in
% \LaTeX~\cite{Lamport:LaTeX}; this section will simply show a few
% examples of in-text equations in context. Notice how this equation:
% \begin{math}
%   \lim_{n\rightarrow \infty}x=0
% \end{math},
% set here in in-line math style, looks slightly different when
% set in display style.  (See next section).

% \subsection{Display Equations}
% A numbered display equation---one set off by vertical space from the
% text and centered horizontally---is produced by the \textbf{equation}
% environment. An unnumbered display equation is produced by the
% \textbf{displaymath} environment.

% Again, in either environment, you can use any of the symbols and
% structures available in \LaTeX\@; this section will just give a couple
% of examples of display equations in context.  First, consider the
% equation, shown as an inline equation above:
% \begin{equation}
%   \lim_{n\rightarrow \infty}x=0
% \end{equation}
% Notice how it is formatted somewhat differently in
% the \textbf{displaymath}
% environment.  Now, we'll enter an unnumbered equation:
% \begin{displaymath}
%   \sum_{i=0}^{\infty} x + 1
% \end{displaymath}
% and follow it with another numbered equation:
% \begin{equation}
%   \sum_{i=0}^{\infty}x_i=\int_{0}^{\pi+2} f
% \end{equation}
% just to demonstrate \LaTeX's able handling of numbering.

% \section{Figures}

% The ``\verb|figure|'' environment should be used for figures. One or
% more images can be placed within a figure. If your figure contains
% third-party material, you must clearly identify it as such, as shown
% in the example below.
% \begin{figure}[h]
%   \centering
%   \includegraphics[width=\linewidth]{sample-franklin}
%   \caption{1907 Franklin Model D roadster. Photograph by Harris \&
%     Ewing, Inc. [Public domain], via Wikimedia
%     Commons. (\url{https://goo.gl/VLCRBB}).}
%   \Description{A woman and a girl in white dresses sit in an open car.}
% \end{figure}

% Your figures should contain a caption which describes the figure to
% the reader.

% Figure captions are placed {\itshape below} the figure.

% Every figure should also have a figure description unless it is purely
% decorative. These descriptions convey what’s in the image to someone
% who cannot see it. They are also used by search engine crawlers for
% indexing images, and when images cannot be loaded.

% A figure description must be unformatted plain text less than 2000
% characters long (including spaces).  {\bfseries Figure descriptions
%   should not repeat the figure caption – their purpose is to capture
%   important information that is not already provided in the caption or
%   the main text of the paper.} For figures that convey important and
% complex new information, a short text description may not be
% adequate. More complex alternative descriptions can be placed in an
% appendix and referenced in a short figure description. For example,
% provide a data table capturing the information in a bar chart, or a
% structured list representing a graph.  For additional information
% regarding how best to write figure descriptions and why doing this is
% so important, please see
% \url{https://www.acm.org/publications/taps/describing-figures/}.

% \subsection{The ``Teaser Figure''}

% A ``teaser figure'' is an image, or set of images in one figure, that
% are placed after all author and affiliation information, and before
% the body of the article, spanning the page. If you wish to have such a
% figure in your article, place the command immediately before the
% \verb|\maketitle| command:
% \begin{verbatim}
%   \begin{teaserfigure}
%     \includegraphics[width=\textwidth]{sampleteaser}
%     \caption{figure caption}
%     \Description{figure description}
%   \end{teaserfigure}
% \end{verbatim}

% \section{Citations and Bibliographies}

% The use of \BibTeX\ for the preparation and formatting of one's
% references is strongly recommended. Authors' names should be complete
% --- use full first names (``Donald E. Knuth'') not initials
% (``D. E. Knuth'') --- and the salient identifying features of a
% reference should be included: title, year, volume, number, pages,
% article DOI, etc.

% The bibliography is included in your source document with these two
% commands, placed just before the \verb|\end{document}| command:
% \begin{verbatim}
%   \bibliographystyle{ACM-Reference-Format}
%   \bibliography{bibfile}
% \end{verbatim}
% where ``\verb|bibfile|'' is the name, without the ``\verb|.bib|''
% suffix, of the \BibTeX\ file.

% Citations and references are numbered by default. A small number of
% ACM publications have citations and references formatted in the
% ``author year'' style; for these exceptions, please include this
% command in the {\bfseries preamble} (before the command
% ``\verb|\begin{document}|'') of your \LaTeX\ source:
% \begin{verbatim}
%   \citestyle{acmauthoryear}
% \end{verbatim}

%   Some examples.  A paginated journal article \cite{Abril07}, an
%   enumerated journal article \cite{Cohen07}, a reference to an entire
%   issue \cite{JCohen96}, a monograph (whole book) \cite{Kosiur01}, a
%   monograph/whole book in a series (see 2a in spec. document)
%   \cite{Harel79}, a divisible-book such as an anthology or compilation
%   \cite{Editor00} followed by the same example, however we only output
%   the series if the volume number is given \cite{Editor00a} (so
%   Editor00a's series should NOT be present since it has no vol. no.),
%   a chapter in a divisible book \cite{Spector90}, a chapter in a
%   divisible book in a series \cite{Douglass98}, a multi-volume work as
%   book \cite{Knuth97}, a couple of articles in a proceedings (of a
%   conference, symposium, workshop for example) (paginated proceedings
%   article) \cite{Andler79, Hagerup1993}, a proceedings article with
%   all possible elements \cite{Smith10}, an example of an enumerated
%   proceedings article \cite{VanGundy07}, an informally published work
%   \cite{Harel78}, a couple of preprints \cite{Bornmann2019,
%     AnzarootPBM14}, a doctoral dissertation \cite{Clarkson85}, a
%   master's thesis: \cite{anisi03}, an online document / world wide web
%   resource \cite{Thornburg01, Ablamowicz07, Poker06}, a video game
%   (Case 1) \cite{Obama08} and (Case 2) \cite{Novak03} and \cite{Lee05}
%   and (Case 3) a patent \cite{JoeScientist001}, work accepted for
%   publication \cite{rous08}, 'YYYYb'-test for prolific author
%   \cite{SaeediMEJ10} and \cite{SaeediJETC10}. Other cites might
%   contain 'duplicate' DOI and URLs (some SIAM articles)
%   \cite{Kirschmer:2010:AEI:1958016.1958018}. Boris / Barbara Beeton:
%   multi-volume works as books \cite{MR781536} and \cite{MR781537}. A
%   couple of citations with DOIs:
%   \cite{2004:ITE:1009386.1010128,Kirschmer:2010:AEI:1958016.1958018}. Online
%   citations: \cite{TUGInstmem, Thornburg01, CTANacmart}. Artifacts:
%   \cite{R} and \cite{UMassCitations}.

% \section{Acknowledgments}

% Identification of funding sources and other support, and thanks to
% individuals and groups that assisted in the research and the
% preparation of the work should be included in an acknowledgment
% section, which is placed just before the reference section in your
% document.

% This section has a special environment:
% \begin{verbatim}
%   \begin{acks}
%   ...
%   \end{acks}
% \end{verbatim}
% so that the information contained therein can be more easily collected
% during the article metadata extraction phase, and to ensure
% consistency in the spelling of the section heading.

% Authors should not prepare this section as a numbered or unnumbered {\verb|\section|}; please use the ``{\verb|acks|}'' environment.

% \section{Appendices}

% If your work needs an appendix, add it before the
% ``\verb|\end{document}|'' command at the conclusion of your source
% document.

% Start the appendix with the ``\verb|appendix|'' command:
% \begin{verbatim}
%   \appendix
% \end{verbatim}
% and note that in the appendix, sections are lettered, not
% numbered. This document has two appendices, demonstrating the section
% and subsection identification method.

% \section{Multi-language papers}

% Papers may be written in languages other than English or include
% titles, subtitles, keywords and abstracts in different languages (as a
% rule, a paper in a language other than English should include an
% English title and an English abstract).  Use \verb|language=...| for
% every language used in the paper.  The last language indicated is the
% main language of the paper.  For example, a French paper with
% additional titles and abstracts in English and German may start with
% the following command
% \begin{verbatim}
% \documentclass[sigconf, language=english, language=german,
%                language=french]{acmart}
% \end{verbatim}

% The title, subtitle, keywords and abstract will be typeset in the main
% language of the paper.  The commands \verb|\translatedXXX|, \verb|XXX|
% begin title, subtitle and keywords, can be used to set these elements
% in the other languages.  The environment \verb|translatedabstract| is
% used to set the translation of the abstract.  These commands and
% environment have a mandatory first argument: the language of the
% second argument.  See \verb|sample-sigconf-i13n.tex| file for examples
% of their usage.

% \section{SIGCHI Extended Abstracts}

% The ``\verb|sigchi-a|'' template style (available only in \LaTeX\ and
% not in Word) produces a landscape-orientation formatted article, with
% a wide left margin. Three environments are available for use with the
% ``\verb|sigchi-a|'' template style, and produce formatted output in
% the margin:
% \begin{itemize}
% \item {\verb|sidebar|}:  Place formatted text in the margin.
% \item {\verb|marginfigure|}: Place a figure in the margin.
% \item {\verb|margintable|}: Place a table in the margin.
% \end{itemize}

% %%
% %% The acknowledgments section is defined using the "acks" environment
% %% (and NOT an unnumbered section). This ensures the proper
% %% identification of the section in the article metadata, and the
% %% consistent spelling of the heading.
% \begin{acks}
% To Robert, for the bagels and explaining CMYK and color spaces.
% \end{acks}

%%
%% The next two lines define the bibliography style to be used, and
%% the bibliography file.
\bibliographystyle{ACM-Reference-Format}
\bibliography{main}

%%
%% If your work has an appendix, this is the place to put it.

\clearpage
\appendix

% \section{Proofs for Sections~\ref{sec:preliminaries}-\ref{sec:some-variants}}

% \subsection{Proof-theoretic results}

% \begin{proof}[Proof of~\Cref{prop:modal-succedent-derivations}.]\label{prf:modal-succedent-derivations}
% The proof is by induction on the size of the derivation. The case where the last rule of $\der$ is an instance of $\id, \zero, \sqr, \cnd_{\n}, \cnd_\sq$, or $\srec$ are trivial. If  the last rule of $\der$ is an instance of $\exch_{\n},\exch_{\sq},  \sql,\succ i$, and $\wk_\sq$ then we apply the induction hypothesis. Let us now suppose that $\der$ has been obtained from a derivation $\der_0$ by applying an instance of  $\wk_\n$. By induction hypothesis, there exists a derivation  $\pder0^*:\lists{n}{\sn}{\sn} \seqar \sn$ such that $\model{\pder0}(\vec x; \vec y)= \model{\pder0^*}(\vec x; )$. Since $\model{\der}(\vec x; \vec y, y)= \model{\pder0}(\vec x; \vec y)=\model{\pder0^*}(\vec x; )$ we just set $\der^*= \pder{0}^*$.  Suppose now that $\der$ is obtained from two derivations $\der_0$ and $\der_1$ by applying an instance of  $\cut_\n$. By induction hypothesis, there exists $\pder{1}^*$ such that $\model{\pder{1}}(\vec x;\vec y,y)=\model{\pder{1}^*}(\vec x;)$. Since 
%  $\model{\pder{1}}(\vec x; \vec y,  \model{\pder{0}}(\vec x; \vec y))= \model{\pder{1}^*}(\vec x;)$, we set $\der^*=\pder{1}^*$. As for the case where the last rule is $\cut_\sq$, by induction hypothesis, there exist derivations $\pder{0}^*$ and $\pder{1}^*$ such that $\model{\pder{0}}(\vec x; \vec y)= \model{\pder{0}^*}(\vec x;)$ and $\model{\pder{1}}(\vec x,x; \vec y)= \model{\pder{1}^*}(\vec x, x;)$, so that we define $\der^*$ as the derivation obtained from $\pder{0}^*$ and $\pder{1}^*$ by applying the rule $\cut_{\sq}$.
% \end{proof}










\section{Computational expressivity of (progressing) coderivations}
\label{sec:computational-strength}
\subsection{Examples of coderivations for~\Cref{prop:computational-props-coderivations}}

% First, let us exemplify the capacity of regular progressing $\bcnorec$ coderivations to express arbitrary primitive recursion (on notation):
% \begin{exmp}
% [Regular + progressing $\not\Rightarrow$ safety]
% \label{exmp:non-safe-recursion}
% Let $g(\vec x; \vec y)$ and $h_i( x, \vec x, z; \vec y)$ with $i=0,1$ be functions definable by some (regular and progressing) $\bcnorec$-coderivations $\gfunction$ and $\hfunction_i$, respectively. Then the following \emph{non-safe} recursion scheme,
% % \[
% % \begin{array}{rcll}
% %  f(0, \vec x; \vec y) & = & g(\vec x; \vec y)\\
% %   f(\succ ix,\vec x; \vec y) & = &h_i( x, \vec x, f(x, \vec x; \vec y); \vec y)  
% % \end{array}
% % \]
% % can be defined by the following (regular and progressing) $\bcnorec$-coderivation:
% % \[
% % \vlderivation{
% % \vliin{\cnd_\sq}{\bullet}{\underline{\blue{\sn}}, \sq \vec \n, \vec \n \seqar \sn}
% % {
% % \vltr{\gfunction}{\sq \vec \n, \vec \n \seqar \sn}{\vlhy{\ }}{\vlhy{\ }}{\vlhy{\ }}
% % }
% % {
% %   \vliin{\cut_\sq}{}{\blue{\sn}, \sq \vec \n, \vec \n \seqar \sn}
% %   {
% %   \vlin{\cnd_{\sq}}{\bullet}{\blue{\sn}, \sq \vec \n, \vec \n \seqar \red{\sn}}{\vlhy{\vdots}}
% %   }{
% %   \vltr{\hfunction_i}{\blue{\sn}, \sq \vec \n, \red{\sn}, \vec \n \seqar \sn}{\vlhy{\ }}{\vlhy{\ }}{\vlhy{\ }}
% %   }
% % }
% % {
% % }
% % }
% % \]
% % where we  identified the coderivations corresponding to the second and the third premise of the conditional rule, as they only differ on the sub-coderivation $\hfunction_i$ ($i=0$ for the former and $i=1$ for the latter).
% \begin{equation}
%     \label{eq:norm-prim-rec-scheme}
%       \arraycolsep=2pt
%     \begin{array}{rcll}
%  f(0, \vec x;) & = & g(\vec x; )\\
%   f(\succ ix,\vec x; ) & = &h_i( x, \vec x, f(x, \vec x; ); )  
% \end{array}
% \end{equation}
% can be defined by the following (regular and progressing) $\bcnorec$-coderivation:
% \[
% \small
% \vlderivation{
% \vliin{\cnd_\sq}{\bullet}{\underline{\blue{\sn}}, \sq \vec \n \seqar \n}
% {
% \vltr{\gfunction}{\sq \vec \n \seqar \n}{\vlhy{\ }}{\vlhy{\ }}{\vlhy{\ }}
% }
% {
%   \vliin{\cut_\sq}{{\scriptstyle{i=0,1}}}{\blue{\sn}, \sq \vec \n \seqar \n}
%   {
%   \vlin{\sqr}{}{\blue{\sn}, \sq \vec \n  \seqar \red{\sn}}{
%     \vlin{\cnd_{\sq}}{\bullet}{\blue{\sn}, \sq \vec \n  \seqar {\n}}{\vlhy{\vdots}}
%   }
%   }{
%   \vltr{\hfunction_i}{\blue{\sn}, \sq \vec \n, \red{\sn}  \seqar \n}{\vlhy{\ }}{\vlhy{\ }}{\vlhy{\ }}
%   }
% }
% {
% }
% }
% \]
% % where we  identified the coderivations corresponding to the second and the third premise of the conditional rule, as they only differ on the sub-coderivation $\hfunction_i$ ($i=0$ for the former and $i=1$ for the latter).
% Note in particular that we can use the scheme in \Cref{eq:norm-prim-rec-scheme} to define all primitive recursive functions.
% \end{exmp}


In light of~\Cref{prop:modal-nonmodal-distinction}, we shall simply omit modalities in (regular) (progressing) coderivations in what follows, i.e.~we shall regard any formula in the antecedent of a sequent as modal and we shall omit applications of $\sqr$. Consequently,  we shall write e.g.~$\cut$ instead of $\cut_{\sq}$ (and so on) and avoid writing semicolons in the semantics of a coderivation. 



From here, \Cref{item:type1-complete,item:turing-complete} from \Cref{prop:computational-props-coderivations} are proved by way of the following examples. 

\begin{exmp}[Extensional completeness at type 1]\label{exmp:extensional-completeness-type-1}   For any function  $f: \Nat^k \to \Nat$ there is a progressing coderivation $\completeness$ such that $\denot\completeness=f$. Proceeding by induction on $k$, if $k=0$ then we use the rules $\zero, \succ 0, \succ 1 , \cut$ to construct $\completeness$ defining the natural number $f$. Otherwise, suppose $f:\Nat \times \Nat^k \to \Nat$ and define $f_n$ as $f_n(\vec x)=f(n, \vec x)$. We construct the coderivation defining $f$ as follows:
\[
\tiny
\vlderivation{
\vliin{\cut}{}{\n, \vec \n \seqar \n}
{
   \vltr{\binarytounary}{\n,  \vec \n \seqar \red{\n}}{\vlhy{\ }}{\vlhy{\ }}{\vlhy{\ }}
}
{
 \vliin{\cnd}{}{\red{\n}, \vec \n, \seqar \n}
  {
        \vltr{f_0}{\vec \n \seqar \n}{\vlhy{\ }}{\vlhy{\ }}{\vlhy{\ }}
  }
   {
        \vliin{\cnd}{}{\red{\n}, \vec \n \seqar \n}
  {
        \vltr{f_1}{\vec \n \seqar \n}{\vlhy{\ }}{\vlhy{\ }}{\vlhy{\ }}
  }
  {
        \vliin{\cnd}{}{\red{\n}, \vec \n \seqar \n}
  {
        \vltr{f_2}{\vec \n \seqar \n}{\vlhy{\ }}{\vlhy{\ }}{\vlhy{\ }}
  }
  {
        \vlin{\cnd}{}{\red{\n}, \vec \n \seqar \n}{\vlhy{\vdots}}
  }
    }
  }
}
}
\]
where $\binarytounary$ is a coderivation converting $n$ into $\succ 1 \overset{n}{\ldots} \succ 1\zero$, and we omit the second premise of $\cnd$ as it is never selected. Notice that $\completeness$  is progressing, as red formulas $\red{\n}$ (which are modal) form a progressing thread. 

% \[
% \vlderivation{
% \vliiin{\cnd_{\sq}}{}{\blue{\sn} \seqar \n}
% {
% \vltr{X(0)}{\seqar \n}{\vlhy{\ }}{\vlhy{\ \ \     }}{\vlhy{\ }}
% }
% {
% % \vlin{\wk}{}{\blue{\sn}\seqar \n}{ \vltr{X(0)}{\seqar \n}{\vlhy{\ }}{\vlhy{\ \ \     }}{\vlhy{\ }}}
%   \vliiin{\cnd_{\sq}}{}{\blue{\sn} \seqar \n}
%     {
%     \vltr{X(0)}{\seqar \n}{\vlhy{\ }}{\vlhy{\ \ \     }}{\vlhy{\ }}
%   }
%   {
%         \vliin{\cnd_{\sq}}{}{\blue{\sn} \seqar \n}
%       {
%         \vltr{X(0)}{\seqar \n}{\vlhy{\ }}{\vlhy{\ \ \            
%         }}{\vlhy{\ }}
%          }
%          {
%          \vlhy{\vdots}
%          }
%   }
%   {
%           \vliin{\cnd_{\sq}}{}{\blue{\sn} \seqar \n}
%       {
%         \vltr{X(2)}{\seqar \n}{\vlhy{\ }}{\vlhy{\ \ \        
%         }}{\vlhy{\ }}
%          }
%          {
%           \vlhy{\vdots}
%          }
%   }
% }
% {
%   \vliiin{\cnd_{\sq}}{}{\blue{\sn} \seqar \n}
%     {
%     \vltr{X(1)}{\seqar \n}{\vlhy{\ }}{\vlhy{\ \ \     }}{\vlhy{\ }}
%   }
%   {
%          \vliin{\cnd_{\sq}}{}{\blue{\sn} \seqar \n}
%       {
%         \vltr{X(1)}{\seqar \n}{\vlhy{\ }}{\vlhy{\ \ \             
%         }}{\vlhy{\ }}
%          }
%          {
%           \vlhy{\vdots}
%          }
%   }
%   {
%           \vliin{\cnd_{\sq}}{}{\blue{\sn} \seqar \n}
%       {
%         \vltr{X(3)}{\seqar \n}{\vlhy{\ }}{\vlhy{\ \ \            
%         }}{\vlhy{\ }}
%          }
%          {
%          \vlhy{\vdots}
%          }
%   }
% }
% }
% \]
% where the coderivation labelled $X(n)$   is $\vlinf{\zero}{}{\seqar \n}{\vlhy{\ }}$ if $X(n)=0$ and  $\vlderivation{\vliin{\cut_\n}{}{\seqar \n}
% {\vlin{\zero}{}{\seqar \n}{\vlhy{\ }}}
% {\vlin{\succ 1}{}{\n \seqar \n}{   \vlin{\id}{}{\n \seqar \n}{\vlhy{\ }}}}
% }$ otherwise. Notice that the above coderivation, which is essentially the graph of the function $X(n)$, is not regular in general by a cardinality argument:  there are only countable many regular coderivations, while there are  continuum many subsets of $\Nat$. Moreover, the above coderivation is progressing, as blue formulas $\blue{\sn}$ form progressing threads, each one contained in an infinite branch.
\end{exmp}

The above example illustrates the role of regularity as a \emph{uniformity} condition:  regular coderivations admit a finite description  (e.g.~a finite tree with backpointers), so that they define  computable functions. In fact, it turns out that any (partial) recursive function  is $\bcnorec$-definable by a regular coderivation. This can be easily inferred from  the next two examples using Proposition~\ref{prop:modal-nonmodal-distinction}.

\begin{exmp}[Primitive recursion] \label{exmp:primitive-recursion} Let  $f(x,\vec x)$ be defined by primitive recursion from  $g(\vec x)$ and $h(x, \vec x,y)$. Given  coderivations $\gfunction$ and $\hfunction$ defining $g$ and $h$ respectively, we construct the following coderivation $\primrec$:
\[
\tiny
\vlderivation{
\vliin{\cut}{}{\red{\n}, \vec \n \seqar \n}
{\vltr{\binarytounary}{\red{\n} \seqar \blue{\n}}{\vlhy{\ }}{\vlhy{\ }}{\vlhy{\ }}}{ 
     \vliin{\cnd}{\bullet}{\underline{\blue{\n}},\red{\n},   \vec \n  \seqar \n}
     {
    \vltr{\gfunction}{ \vec \n \seqar \n}{\vlhy{\ }}{\vlhy{\ }}{\vlhy{\ }}
     }
     {
        \vliin{\cut}{}{\blue{\n},\red{\n},  \vec \n \seqar \n}{
                  \vltr{\predecessor}{\red{\n} \seqar  \orange{\n}}{\vlhy{\ }}{\vlhy{\ }}{\vlhy{\ }}
                 }
                 {
            \vliin{\cut}{}{\blue{\n},\orange{\n},  \vec \n \seqar \n}{
                  \vlin{\cnd}{\bullet}{\blue{\n},\orange{\n},  \vec \n \seqar \green{\n}}{\vlhy{\vdots}}
                 }{
             \vltr{\hfunction}{\orange{\n},   \vec \n, \green{\n} \seqar \n}{\vlhy{\ }}{\vlhy{\ }}{\vlhy{\ }}
                 }
                 }
    }
}
}
\]
where $\binarytounary$ is the coderivation converting $n$ into $\succ 1\overset{n}{\ldots} \succ 1 \zero$, $\predecessor$ is a coderivation defining unary predecessor, and we avoid writing the second premise of $\cnd$ as it is never selected.  Notice that $\primrec$  is progressing, as blue formulas $\blue{\n}$ (which are modal) form a progressing thread contained in the infinite branch that loops on $\bullet$. From the associated equational program we obtain:
\[
\arraycolsep=2pt
\begin{array}{rcll}
\denot{\primrec_{\epsilon}}  (x, \vec x) & = &  \denot{\primrec{1}}(\denot{\binarytounary}(x),x, \vec x)\\
    \denot{\primrec_{1}}  (0,x, \vec x) & = & g(\vec x)\\
     \denot{\primrec_{1}}  (\succ 1z, x , \vec x ) & = & h(\denot{\predecessor}(x), \vec x,  \denot{\primrec_{1}}(z,\denot{\predecessor}(x), \vec x))
\end{array}
\]
so that $\denot\primrec=\denot{\primrec{\epsilon}}=f$.
% \[
% \vlderivation{
% \vliiin{\cnd_{\sq}}{\bullet}{\underline{\blue{\sn}},\sn,  \sq \vec \n  \seqar \n}
% {
% \vlin{\wk}{}{\underline{\sn}, \sq \vec \n \seqar \n}{\vltr{g}{\sq \vec \n \seqar \n}{\vlhy{\ }}{\vlhy{\ }}{\vlhy{\ }}}
% }
% {
%   \vliq{\wk}{}{\blue{\sn},\sn, \sq \vec \n \seqar \n}{\vlin{\zero}{}{\seqar \n }{\vlhy{\ }}}
% }
% {
%   \vliin{\cut_{\sq}}{}{\blue{\sn},\sn, \sq \vec \n \seqar \n}
%   {
%   \vlin{\sqr}{}{\blue{\sn},\sn, \sq \vec \n \seqar \red{\sn} }{ \vlin{\cnd_{\sq}}{\bullet}{\blue{\sn},\sn, \sq \vec \n \seqar \n}{\vlhy{\vdots}}}
%   }
%   {
%   \vlin{\wk}{}{\underline{\blue{\sn}},\sn , \sq \vec \n, \red{\sn} \seqar \n}{
%       \vliin{\cut_{\sq}}{}{\sn,  \sq \vec \n, \sn \seqar \n}{\vlin{\sqr}{}{\sn, \sq \vec \n, \sn \seqar \orange{\sn}}{\vliq{\wk}{}{\sn,\sn, \sq \vec \n \seqar \n}{\vltr{\mathcal{E}}{\sn \seqar  \n}{\vlhy{\ }}{\vlhy{\ }}{\vlhy{\ }}}}}{\vlin{\wk}{}{\orange{\sn}, \underline{\sn}, \sq \vec \n, \sn \seqar \n}{\vltr{h}{\sn,  \sq \vec \n, \sn \seqar \n}{\vlhy{\ }}{\vlhy{\ }}{\vlhy{\ }}}}}
%   } 
% }
% }
% \]
% where $\mathcal{E}$ is the coderivation defining unary predecessor in Example~\ref{exmp:predecessor} (See Appendix~\ref{sec:additional-examples}). 
% From the associated equational program we obtain:
% \[
% \begin{array}{rcll}
%     \denot{\der'_{\epsilon}}  (0,x, \vec x;) & = & g(\vec x;)\\
%      \denot{\der'_{\epsilon}}  (\succ 0(;z),x, \vec x;) & = & \zero(;) & x \neq 0 \\ 
%      \denot{\der'_{\epsilon}}  (\succ 1(z;), x , \vec x ;) & = & h(\denot{\mathcal{E}}(x;), \vec x,  \denot{\der'_{\epsilon}}(z,x,)  ;)
% \end{array}
% \]
% Then we set $\der$ as 
% \[
% \vlderivation{
% \vliin{\cut_{\sq}}{}{\sn, \sq \vec \n \seqar \n} 
% {
% \vlin{\sqr}{}{\sn,\sq \vec \n \seqar \red{\sn}}{\vliq{\wk}{}{\sn, \underline{\sq \vec \n} \seqar \n}{\vltr{\mathcal{G}}{\sn \seqar \n}{\vlhy{\ }}{\vlhy{\ }}{\vlhy{\ }}}
% }}
% {
% \vltr{\der'}{\red{\sn},\sn, \sq \vec \n  \seqar \n}{\vlhy{\ }}{\vlhy{\ }}{\vlhy{\ }}
% }
% }
% \]
% where $\mathcal{G}$ is the coderivation in Example~\ref{exmp:binary-to-unary} (See Appendix~\ref{sec:additional-examples}), so that:
% \[
% \begin{array}{rcll}
%     \denot{\der_{\epsilon}}  (x, \vec x;) & = &  \denot{\der1}(\denot{\mathcal{G}}(x;),x, \vec x;)
% \end{array}
% \]
% and hence $\denot\der=f$.
%  \todo{In the next section show that safe recursion is representable by a left-leaning coderivation, so that  $\bc$ can be seen as a subsystem of $\bc^-$.}
\end{exmp}


\begin{exmp}
[Unbounded search] Let $g(x,\vec{x})$ be a function, and let $f(\vec x)\dfn \mu x. (g(x, \vec x)=0)$  be the unbounded search function obtained by applying the minimisation operation on $g(\vec x)$. Given  a coderivation $\gfunction$  defining $g$, we construct the following coderivation $\unbounded$:
\[
\tiny
\vlderivation{
\vliin{\cut}{}{ \vec \n \seqar \n}{
    \vlin{\zero}{}{\seqar\n }{\vlhy{\ }}
}{
\vliin{\cut}{\bullet}{\blue{\n},  \vec \n \seqar \n}
{
\vltr{\gfunction}{\blue{\n},  \vec \n \seqar \red{\n}}{\vlhy{\ }}{\vlhy{\ }}{\vlhy{\ }}
}
{
     \vliin{\cnd}{}{\red{\underline{\n}},\blue{\n},  \vec \n \seqar \n}{
    \vlin{\id}{}{\blue{\n} \seqar \n}{\vlhy{\ }}
     }{ 
         \vliin{\cut}{}{\blue{\n},  \vec \n \seqar \n}{
               \vltr{\successor}{
                \blue{\n} \seqar \orange{\n}
                 }{
                 \vlhy{\ }
                }{
                \vlhy{\ }
                }
                {\vlhy{\ }
                }
      }{
      \vlin{\cut}{\bullet}{\orange{\n},   \vec \n \seqar \n}{\vlhy{\vdots}}
      }
     }
}
{
}
}
}
\]
where the coderivation $\binarytounary$  computes the unary successor, and we identify the sub-coderivations corresponding to the second and the third premises of the conditional rule.   It is easy to check that the above coderivation is regular but not progressing, as  threads containing  principal formulas for $\cnd$ are finite.  From the associated equational program we obtain:
% \[
% \allowdisplaybreaks
% \begin{array}{rcll}
%     \denot{\pder\epsilon}  (x, \vec x;) & = & \denot{\pder1}(\denot{\pder0}(x, \vec x;),x, \vec x;) \\
%   \denot{\pder0}(x,\vec y;) & = & f(x,\vec x;)\\
%   \denot{\pder{1}}(0,x, \vec x;) & = & \denot{\pder{10}}(x, \vec x;)\\
%   \denot{\pder{10}}(x, \vec x;)& = &\proj{n;}{1;}(x, \vec x;)\\
%   \denot{\pder{1}}(\succ 0(;z),x, \vec x;) & = & \denot{\pder{11}}(z, x, \vec x;)\\
%     \denot{\pder{11}}(z, x, \vec x;) & = & \denot{\pder{110}}(x, \vec x;)\\
%   \denot{\pder{110}}(x, \vec x;) & = &   \denot{\pder{1101}}( \denot{\pder{1100}}( x, \vec x;), x, \vec x;)   \\
%   \denot{\pder{1100}}( x, \vec x;)& = &  \denot{\pder{11000}}( x;)\\
%   \denot{\pder{11000}}( x;)& = & \mathsf{s}(x;)\\
%   \denot{\pder{1101}}(z,x, \vec x;)& = &  \denot{\pder{11010}}(z, \vec x;) \\
%   \denot{\pder{11010}}(z, \vec x;)& = &  \denot{\pder{\epsilon}}(z, \vec x;)\\  \denot{\pder{2}}(\succ 1(;z),x, \vec x;) & = & \denot{\pder{21}}(z, x, \vec x;)\\
%     \denot{\pder{21}}(z, x, \vec x;) & = & \denot{\pder{210}}(x, \vec x;)\\
%   \denot{\pder{210}}(x, \vec x;) & = &   \denot{\pder{2101}}( \denot{\pder{2100}}( x, \vec x;), x, \vec x;)   \\
%   \denot{\pder{2100}}( x, \vec x;)& = &  \denot{\pder{21000}}( x;)\\
%   \denot{\pder{21000}}( x;)& = & \mathsf{s}(x;)\\
%   \denot{\pder{2101}}(z,x, \vec x;)& = &  \denot{\pder{21010}}(z, \vec x;) \\
%   \denot{\pder{21010}}(z, \vec x;)& = &  \denot{\pder{\epsilon}}(z, \vec x;)
% \end{array}
% \]
% By purely equational reasoning, we can simplify this program to obtain the following one,
\[
\arraycolsep=2pt
\begin{array}{rcll}
 \denot{\unbounded_{\epsilon}}  (\vec x) & = & \denot{\unbounded_{1}}(0, \vec x) \\
  \denot{\unbounded_{1}}  (x,\vec x) & = & \denot{\unbounded_{11}}(\denot{\gfunction}(x,\vec{x}) ,x, \vec x) \\
   \denot{\unbounded_{11}}(0,x, \vec x) & = & x \\
   \denot{\unbounded_{11}}(\succ 0z,x, \vec x) & = &\denot{\unbounded_{1}}( \denot\successor(x), \vec x) & z \neq 0 \\
   \denot{\unbounded_{11}}(\succ 1z,x, \vec x) & = & \denot{\unbounded_{1}}( \denot\successor(x), \vec x)
\end{array}
\]
Which searches for the least $x \geq 0$ such that $g(x,\vec x)=0$. Hence, $\denot{\unbounded_\epsilon}(\vec x)=\denot{\unbounded}(\vec x)=f(\vec x)$.
\end{exmp}

\subsection{Completeness for type 1 primitive recursion}

\cite{Kuperberg-Pous21} shows that, in the absence of contraction rules, \emph{only} the primitive recursive functions are so definable (even when using arbitrary finite types).
It is tempting, therefore, to conjecture that the regular and progressing $\bcnorec$-coderivations define \emph{just} the primitive recursive functions.

However there is a crucial difference between our formulation of cut and that in \cite{Kuperberg-Pous21}, namely that ours is context-sharing and theirs is context-splitting.
Thus the former admits a quite controlled form of contraction that, perhaps surprisingly, is enough to simulate the type 0 fragment $\ct_0$ from \cite{Das2021} (which has explicit contraction).


\begin{proof}[Proof of \Cref{item:type1-prim-rec} from \Cref{prop:computational-props-coderivations}]
First, by~\Cref{prop:modal-nonmodal-distinction} we can neglect modalities in $\bcnorec$-coderivations (and semicolons in the corresponding semantics). The left-right implication thus follows from the natural inclusion of our system into $\ct_0$ and \cite[Corollary 80]{Das2021}. 

Concerning the right-left implication, 
% it is not hard to prove that, if $f(\vec x; )$ is definable by a regular progressing $\bcnorec+ \{\cntr_\n, \cntr_\sn \}$-coderivation  then it is definable in $\ct_0$,  the type level 0 fragment of the circular version of G\"{o}del's system $\tgodel$. From the results in~\cite{Das2021} we have that 
%  the functions definable in $\ct_0$ are exactly those ones definable in the type level $1$ fragment of system $\tgodel$, i.e.~$\ntgodel{1}$. 
we employ a formulation $\ntgodel 1 (\vec a)$ of the type 1 functions of $\ntgodel 1$ over oracles $\vec a$ as follows. 
$\ntgodel 1(\vec a)$ is defined just like the primitive recursive functions, including oracles $\vec a$ as initial functions, and by adding the following version of type $1$ recursion:
  \begin{itemize}
      \item if $g(\vec x) \in \ntgodel 1 (\vec a)$ and $h(a)(x,\vec x) \in \ntgodel 1 (a, \vec a)$, then the $f(x, \vec x)$ given by,
    \[
    \arraycolsep=2pt
    \begin{array}{rcl}
        f(0, \vec x) & = & g(\vec x) \\
        f( \succ i x , \vec x) & = & h_i (\lambda \vec u . f(x, \vec u)) (x, \vec x)
    \end{array}
    \]
    is in $\ntgodel 1 (\vec a)$.
  \end{itemize}
  It is not hard to see that the type 1 functions of $\ntgodel 1$ are precisely those of $\ntgodel 1 (\emptyset)$ (e.g.\ because of \cite[Appendix A]{Das2021-preprint}).  
  We now proceed by showing how to define the above scheme by regular and progressing $\bcnorec$-coderivations.
% \anupam{I think this last part needs to be rephrased slightly as follows.

% We work with algebras $\ntgodel{1}(\vec a)$ of $\ntgodel{1}$ that extends the usual formulation of primitive recursive functions by relativising to oracles $\vec a$ and adding the following version of type 1 recursion:

% if $g(\vec x) \in \ntgodel 1 (\vec a)$ and $h(a)(x,\vec x) \in \ntgodel 1 (a, \vec a)$, then the function $f(x, \vec x)$ given by,
%     \[
%     \begin{array}{rcl}
%         f(0, \vec x) & = & g(\vec x) \\
%         f( \succ i x , \vec x) & = & h_i (\lambda \vec u . f(x, \vec u)) (x, \vec x)
%     \end{array}
%     \]
%     is in $\ntgodel 1 (\vec a)$.

% In this way, the type 1 functions of $\ntgodel 1$ are precisely those of $\ntgodel 1 (\emptyset)$. 
% }


Given a function $f(\vec x) \in \ntgodel 1 (\vec a)$, we construct a regular progressing coderivation of $\bcnorec$,
\[
\vltreeder{\der_f (\vec a)}{ \vec \n \seqar \n}{}{ \left\{
\vlinf{a_i}{}{ \vec \n \seqar \n}{}
\right\}_i }{}
\]
computing $f(\vec x)$ over $\vec a$, by induction on the definition of $f(\vec x)$.

If $f(\vec x)$ is an initial function,  an oracle, or $f(\vec x) = h(g(\vec x), \vec x)$ then the construction is easy (see, e.g., the proof of~\Cref{thm:nbc-in-ncbc}).

% If $f(\vec x) = h(g(\vec x), \vec x)$ where $\der_g (\vec a)$ and $\der_h (\vec a)$ are already obtained by the inductive hypothesis, then we define $\der_f (\vec a)$ as follows:
% \[
% \vlderivation{
%     \vliin{\cut_\sq}{}{\sq \vec \n \seqar \n}{
%         \vlin{\sqr}{}{\sq \vec \n \seqar \sn}{
%         \vltr{\der_g(\vec a)}{\sq \vec \n \seqar \n }{\vlhy{\quad }}{\vlhy{\  }}{\vlhy{\quad }}
%         }
%     }{
%         \vltr{\der_h (\vec a)}{\sn, \sq \vec \n \seqar \n}{\vlhy{\quad }}{\vlhy{\  }}{\vlhy{\quad }}
%     }
% }
% \]

Let us  consider the case of recursion (which subsumes usual primitive recursion at type 0).
Suppose $f(x, \vec x) \in \ntgodel 1 (\vec a)$ where,
\[
\arraycolsep=2pt
\begin{array}{rcl}
     f(0, \vec x) & = & g(\vec x) \\
     f(\succ 0 x , \vec x) & = & h_0(\lambda \vec u . f(x,\vec u) ) (x, \vec x) \\
     f(\succ 1 x , \vec x) & = & h_1(\lambda \vec u . f(x,\vec u) ) (x, \vec x)
\end{array}
\]
where $\der_g(\vec a)$ and $\der_h(a, \vec a)$ are already obtained by the inductive hypothesis.
We define $\der_f(\vec a)$ as follows:
\[
\vlderivation{
    \vliiin{\cnd}{\bullet}{\n,  \vec \n \seqar \n }{
        \vltr{\der_g (\vec a)}{ \vec \n \seqar \n}{\vlhy{\quad }}{\vlhy{\quad }}{\vlhy{\quad }}
    }{
        \vltr{\n, \der_{h_0}(\der_f(\vec a),\vec a)}{\n ,  \vec \n \seqar \n}{\vlhy{}}{
            \vlin{}{\bullet}{\n,  \vec \n \seqar \n}{\vlhy{\vdots}}
        }{\vlhy{}}
    }{
        \vltr{\n, \der_{h_1}(\der_f(\vec a),\vec a)}{\n ,  \vec \n \seqar \n}{\vlhy{}}{
            \vlin{}{\bullet}{\n,  \vec \n \seqar \n}{\vlhy{\vdots}}
        }{\vlhy{}}
    }
}
\]
Note that the existence of the second and third coderivations above the conditional is given by a similar construction to that of \Cref{lem:to-prove-ncbc-complete-for-nbc} (see also \cite[Lemma 42]{Das2021}).
\end{proof}

\subsection{On the `power' of contraction}
Given the computationally equivalent system $\ct_0$ with contraction from \cite{Das2021}, we can view the above result as a sort of `contraction admissibility' for regular progressing $\bcnorec$-coderivations.
Let us take a moment to make this formal.

Call $\bcnorec+ \{\cntr_\n, \cntr_\sn \}$ the extension of $\bcnorec$ with the rules $\cntr_\n$ and $\cntr_\sn$ below:
\begin{equation} \label{eqn:explicit-contraction}
  \vlinf{\cntr_\n}{}{\orange \Gamma, \blue \n\seqar B}{\orange \Gamma, \blue \n , \blue \n \seqar B}
\qquad 
\vlinf{\cntr_\sq}{}{\orange \Gamma, \blue \sn \seqar B}{\orange \Gamma, \blue \sn , \blue \sn \seqar B}
\end{equation}
 where  the semantics for the new system extends the one for $\bcnorec$ in the obvious way, and the notion of (progressing) thread is induced by the given colouring.\footnote{Note that the totality argument of \Cref{prop:prog-total} still applies in the presence of these rules, cf.~also \cite{Das2021}.} 
% The next proposition shows that, in the setting of regular progressing $\bcnorec$-coderivations,  the rules $\cntr_\n$ and $\cntr_\sn$  do not increase the class of definable functions.
We have:
\begin{cor} $f(\vec x; )$ is definable by a regular progressing  $\bcnorec+ \{\cntr_\n, \cntr_\sn \}$-coderivation  iff it is definable by a regular progressing $\bcnorec$-coderivation.
\end{cor}
\begin{proof}
[Proof idea]
The former system is equivalent to $\ct_0$ from \cite{Das2021}, whose type 1 functions are just those of $\ntgodel 1$ by \cite[Corollary 80]{Das2021}, which are all defined by regular progressing $\bcnorec$-coderivations by \Cref{prop:computational-props-coderivations} \Cref{item:type1-prim-rec}.
\end{proof}


\begin{rem}\label{rem:explicit-contraction}
The reader may at this point wonder if a  direct  ‘contraction-admissibility' argument exists for the rules in~\eqref{eqn:explicit-contraction}. First, notice that  $\cntr_\n$ can be derived in $\bcnorec$:
\[
     \vlderivation{
     \vliin{\cut_\n}{}{\Gamma, \blue \n  \seqar B}
{
   \vliq{\wk_\n}{}{\Gamma, \n \seqar \red \n}{
   \vlin{\id}{}{ \n \seqar \n}{\vlhy{}}}
}
{ 
 \vlhy{\Gamma, \blue \n, \red \n   \seqar B  }
}
}
\]
While a similar derivation exists for $\cntr_{\sq}$,
note that this crucially does not preserve the same notion of thread (cf.~colours above) and so does not, a priori, preserve progressiveness.
\end{rem}



\begin{exmp}
[Ackermann-P\'eter]
As suggested by \Cref{prop:computational-props-coderivations} \Cref{item:type1-prim-rec}, regular progressing $\bcnorec$ coderivations are able to define the Ackermann-P\'eter function $A(x,y)$,
\[
\arraycolsep=2pt
\begin{array}{rcl}
    A(0,y) & = & x+1 \\
    A(x+1,0) & = & A(x,1) \\
    A(x+1,y+1) & = & A(x,A(x+1,y))
\end{array}
\]
which is well-known to not be (type 0) primitive recursive, but is nevertheless type 1 primitive recursive.

The usual cyclic representation of $A(x,y)$ using only base types, e.g.\ as in \cite{Das2021,Kuperberg-Pous21}, mimics the equational program above, namely applying case analysis on the first input before applying case analysis on the second input.
It is important in this construction to be able to explicitly contract the first input, corresponding to $x$, due to the third line of the equational program above.

In our context-sharing setting, without explicit contraction, we may nonetheless represent $A(x,y)$ by conducting a case analysis on the second input $y$ first, and including some redundancy, as in the regular progressing coderivation in \Cref{fig:ack-pet}.
To facilitate readability, we again omit modalities and work purely in unary notation, representing a natural number $n$ by $\succ1 \overset{n}\ldots \succ 10$.

\begin{figure*}
\centering
\begin{equation*}
\begin{array}{c}
\hspace{8em}
\vlderivation{
\vliin{\cnd}{\bullet}{\red{\n},\underline{\blue{\n}}\seqar \n}{
    \vliin{\cnd}{}{\underline{\red{\n}}\seqar \n}{
        \vliq{1}{}{\seqar \n}{\vlhy{}}
    }{
        \vlin{1}{}{\red{\n} \seqar \n}{
        \vlin{\cnd}{\bullet}{\red{\n},\underline{\n}\seqar\n}{\vlhy{\vdots (1)}}
        }
    }
}{
    \vliin{\cut}{}{\red{\n},\blue{\n} \seqar \n}{
        \vlin{\cnd}{\bullet}{\red{\n},\underline{\blue{\n}}\seqar \n}{\vlhy{\vdots(2)}}
    }{
        \vliin{\cnd}{}{\underline{\red{\n}},\blue{\n},\n \seqar \n}{
            \vliq{+2}{}{\blue{\n} \seqar \n}{\vlhy{}}
        }{
            \vlin{\wk}{}{\red{\n},\underline{\blue{\n}},\n \seqar \n}{
            \vlin{\cnd}{\bullet}{\red{\n},\underline{\n} \seqar \n}{\vlhy{\vdots (3)}}
            }
        }
    }
}
}
\end{array}
\end{equation*}
\caption{Ackermann-P\'eter function computed by a regular progressing $\bcnorec$-coderivation (without explicit contractions).}
\label{fig:ack-pet}
\end{figure*}

The verification of the semantics of this coderivation is routine. 
To see that it is progressing, we may conduct a case analysis on the set $L$ of infinitely often visited loops, among $(1),(2),(3)$, in an arbitrary branch:
\begin{itemize}
    \item if $L = \emptyset$ then the branch is finite;
    \item if $L=\{(1)\}$ then there is a progressing thread on the red $\red\n$;
    \item if $L=\{(2)\}$ then there is a progressing thread on the blue $\blue \n$ (note that the red $\red \n $ thread does not progress);
    \item if $L=\{(3)\}$ then there is a progressing thread on the red $\red \n$;
    \item if $L=\{(1),(2)\}$ then there is a progressing thread on the red $\red\n$: on each iteration of $(1)$ the thread progresses, and remains intact (but does not progress) on each iteration of $(2)$;
    \item if $L=\{(1),(3)\}$ then there is a progressing thread on the red $\red \n$: the thread progresses on each iteration of either $(1)$ or $(3)$;
    \item if $L=\{(2),(3)\}$ then there is a progressing thread on the red $\red\n$: on each iteration of $(2)$ the thread remains intact (but does not progress), and on each iteration of $(3)$ the thread progresses;
    \item if $L=\{(1),(2),(3)\}$ then there is a progressing thread on the red $\red\n$: the thread progresses on each iteration of $(1)$ or $(3)$, and remains intact (but does not progress) on each iteration of $(3)$.
\end{itemize}
\end{exmp}

% \section{Proofs of Section~\ref{sec:characterization-results-for-function-algebras}}\label{app:proofs-of-section-characterisation-results-for-function-algebras}
% % \subsection{Proof of the Bounding Lemma}






% \subsection{Soundness proofs}



% \section{Proofs of Section~\ref{sec:completeness}}\label{app:proofs-of-section-completeness}


% \subsection{Completeness}




% \subsection{Proofs of Section~\ref{sec:translation}}



 %%%%%%HERE GARBAGE%%%%%%%%%%%%%

% \clearpage
% \section{Garbage}
% HERE BELOW OLD PROOF OF~\Cref{prop:cutbox}
%  \begin{proof}[Proof of~\Cref{prop:cutbox}.]
%  First, we show that  any infinite branch of $\der$ has finitely many sequents with succedent $\sn$. Suppose towards contradiction that some branch contains infinitely many such sequents. By progressiveness, the branch  also  contains infinitely many conclusions of the $\cnd_\sq$ rule, and hence infinitely many sequents with succedent $\n$. This means that we can construct a family of pairs $(s_i, s'_i)_i$ where $s_i$ is a sequent with shape $\Gamma_i \seqar \n$, $s'_i$ is a sequent with shape $\Delta_i \seqar \sn$, and  $s_i\sqsubset s'_i \sqsubset s_{i+1}$.  By inspecting the rules of $\bcnorec$, we conclude that there exists a family of sequents $(t_i)_i$ such that $s_i\sqsubset t_i\sqsubseteq s'_i $ and $t_i$ is the conclusion of an instance of $\cut_\sq$. Therefore, the infinite branch would contain infinitely many $\cut_\sq$ steps, which is impossible by safety. 

% Now, by the above property, for any infinite branch of $\der$ there exists a node $\nu$ such that all sequents above it have non-modal succedent. The set of all such nodes $X$ forms a bar, and the set of nodes below this bar, say $Y$, is finite by weak  K\"{o}nig lemma. We shall  reason by induction on this fragment of $\der$ and prove that, for any node $\mu\in Y$ such that $\pder\mu: \lists{n}{\sn}{\sn}, \lists{m}{\n}{\n} \seqar \sn$ there exists a coderivation
%  $\pder{\mu}^*:\lists{n}{\sn}{\sn} \seqar \sn$ such that $\model{\pder \mu}(\vec x; \vec y)= \model{\pder{\mu}^*}(\vec x; )$.
%   The cases where $\mu$ is the conclusion of $\id$, $\zero$, $\sqr$,  $\cnd_{\n}$ and $\cnd_{\sq}$ are trivial.  If $\mu$ is the conclusion of $\wk_{\sq}, \exch_{\n}$, $\exch_\sq$, $\sql$ and $\succ i$, then we apply the induction hypothesis. If $\mu$ is the conclusion of $\wk_{\n}$ with premise $\mu'$ then $m>0$ and, by induction hypothesis on $\mu'$, there exists a coderivation  $\pder{\mu'}^*:\lists{n}{\sn}{\sn} \seqar \sn$ such that $\model{\pder{\mu}}(\vec x; \vec y)= \model{\pder{\mu'}^*}(\vec x; )$. Since $\model{\pder\mu}(\vec x; \vec y, y)= \model{\pder{\mu'}}(\vec x; \vec y)=\model{\pder{\mu'}^*}(\vec x; )$ we just set $\pder{\mu}^*= \pder{\mu'}^*$.  Suppose now that $\mu$ is the conclusion of $\cut_{\n}$ with premises $\mu_1$ and $\mu_2$.  Then, by induction hypothesis, there exists $\pder{\mu_2}^*$ such that $\model{\pder{\mu_2}}(\vec x;\vec y,y)=\model{\pder{\mu_2}^*}(\vec x;)$. Since 
%   $\model{\pder{\mu_2}}(\vec x; \vec y,  \model{\pder{\mu_1}}(\vec x; \vec y))= \model{\pder{\mu_2}^*}(\vec x;)$, we set $\pder{\mu}^*=\pder{\mu_2}^*$. Last, suppose that that $\mu$ is the conclusion of $\cut_{\sq}$ with premises $\mu_1$ and $\mu_2$. By induction hypothesis on both $\mu_1$ and $\mu_2$, there exist coderivations $\pder{\mu_1}^*$ and $\pder{\mu_2}^*$ such that $\model{\pder{\mu_1}}(\vec x; \vec y)= \model{\pder{\mu_1}^*}(\vec x;)$ and $\model{\pder{\mu_2}}(\vec x,x; \vec y)= \model{\pder{\mu_2}^*}(\vec x, x;)$, so that we define $\pder{\mu}^*$ as the coderivation obtained from $\pder{\mu_1}^*$ and $\pder{\mu_2}^*$ by applying the rule $\cut_{\sq}$.
%  \end{proof}


%  \begin{prop}  Let  $\der: \sn,\overset{n}{\ldots}, \sn, \n,\overset{m}{\ldots}, \n \seqar C$  be a coderivation of $\ncbc$, and let $\derrd$,  be its cycle nf. If $C=\sn$ then there exists  a coderivation of $\ncbc$  $\der':\sn,\overset{n}{\ldots}, \sn \seqar C$  such that, for all $\vec{x}, \vec{y}\in \Nat$:
%  \begin{equation*}
%      \model{\der}(\vec{x}; \vec{y})= \model{\der'}(\vec{x};).
%  \end{equation*}
%  In particular, if $\der\in \cbc$ then $\der'\in \cbc$.
%  \end{prop}
%  \begin{proof}
%  First we show that, for any $\nu \in \bud(\der)$, the path $\pi$ from $\rd(\nu)$ to $\nu$ contains no sequents with modal succedent. This is an immediate consequence of the following observations:
%  \begin{itemize}
%  \item no inference rule of $\ncbc$ has a conclusion with non-modal succedent and a premise with modal succedent;
%      \item by Proposition~\ref{prop:structureofcycles}.\ref{enum:budcompanion0}, $\pi$ must 
%      contain the conclusion of an  instance of $\cnd_{\sn}$, which has non-modal succedent;
%      \item the nodes $\rd(\nu)$ and $\nu$ are labelled with the same sequent by definition of cycle nf. 
%  \end{itemize}
%  This means that no bud is labelled with a sequent having modal succedent. We now prove that, for any canonical node $\hnu$ such that $\pder\hnu: \sn,\overset{n}{\ldots}, \sn, \n,\overset{m}{\ldots}, \n \seqar \sn$ there exists a coderivation
%  $\pder{\hnu}':\sn,\overset{n}{\ldots}, \sn \seqar \sn$ such that $\model{\pder \hnu}(\vec x; \vec y)= \model{\pder{\hnu}'}(\vec x; )$.
%  The proof is  by induction on the length $\ell$ of the longest path connecting the a canonical node to a terminal node in $\derrd$.  The cases where $\hnu$ is the conclusion of $\id$, $\zero$, $\sqr$ and $\cnd$ are trivial. Moreover, $\hnu$ cannot be the conclusion of a bud by the above reasoning. If $\hnu$ is the conclusion of $\exch$, $\sql$ and $\succ i$, then we apply the induction hypothesis. If $\hnu$ is the conclusion of $\wk$ with premise $\nu$ then we have two cases. If the weakened formula is modal then we simply apply the induction hypothesis. Otherwise $m>0$ and, by induction hypothesis on $\nu$, there exists a coderivation  $\pder{\nu}':\sn,\overset{n}{\ldots}, \sn \seqar \sn$ such that $\model{\pder \nu}(\vec x; \vec y)= \model{\pder{\nu}'}(\vec x; )$. Since $\model{\pder\hnu}(\vec x; \vec y, y)= \model{\pder\nu}(\vec x; \vec y)=\model{\pder{\nu}'}(\vec x; )$ we just set $\pder{\hnu}'= \pder{\nu}'$.  Suppose now that $\hnu$ is the conclusion of $\cut$ with premises $\nu_1$ and $\nu_2$. We have two cases. If the cut-formula is non-modal then, by induction hypothesis, there exists $\pder{\nu_2}'$ such that $\model{\pder{\nu_2}}(\vec x;\vec y,y)=\model{\pder{\nu_2}'}(\vec x;)$. Since 
%   $\model{\pder{\nu_2}}(\vec x; \vec y,  \model{\pder{\nu_1}}(\vec x; \vec y))= \model{\pder{\nu_2}'}(\vec x;)$, we set $\pder{\hnu}'=\pder{\nu_2}'$. Otherwise the cut-formula is modal and, by induction hypothesis on both $\nu_1$ and $\nu_2$, there exist coderivations $\pder{\nu_1}'$ and $\pder{\nu_2}'$ such that $\model{\pder{\nu_1}}(\vec x; \vec y)= \model{\pder{\nu_1}'}(\vec x;)$ and $\model{\pder{\nu_2}}(\vec x,x; \vec y)= \model{\pder{\nu_2}'}(\vec x, x;)$, so that we define $\pder{\hnu}'$ as the coderivation obtained from $\pder{\nu_1}'$ and $\pder{\nu_2}'$ by applying the rule $\cut$.
%  \end{proof}

% \begin{figure*}
%     \centering
% \begin{equation*}
%     \small
% \vlderivation{
% \vliin{\cut_{\n}}{}{\red{\sn}, \blue{\sn} \seqar \n}
%   {
%   \vlin{}{}{\sn \seqar \n}{}{\vliq{\id, \sqr}{}{\blue{\sn} \seqar \blue{\n}}{\vlhy{}}}
%   }{
%   \vliin{\cut_{\n}}{}{\red{\sn}, \blue{\sn},\blue{\n} \seqar \n}
%   {\vliq{\id, \sqr}{}{\red{\sn} \seqar \n}{\vlhy{}}}
%   {
%   \vliin{\cnd_\sn}{\bullet}{\red{\underline{\sn}}, \blue{\sn}, \blue{\n} , \red{\n} \seqar \n }
%      {
%      \vlin{\id}{}{\blue{\n} \seqar \n}{\vlhy{}}
%      }
%      {
%      \vliq{\exch, \exch}{}{\red{\sn}, \blue{\sn}, \blue{\n} , \red{\n} \seqar \n}
%       {
%       \vlin{\cnd_{\sn}}{\bullet}{\blue{\sn}, \red{\sn}, \red{\n}, \blue{\n} \seqar \n}{\vlhy{}}
%       }
%      }
%   }
%   }
% }
% \end{equation*}
% \newcommand{\maximum}{\mathcal{M}}
% The equational program associated to the above coderivation $\maximum$ is the following:
% \[
% \begin{array}{rcl}
%  \denot{\maximum}(x, y; )    &  =& \denot{\der_{11}}(x, y; y, x) \\
%   \denot{\der_{11}}(0, y; u, v)   & =&u\\
%      \denot{\der_{11}}(\succ i x, y; u, v)   & =& \denot{\der_{11}}(y, x; v,u)
% \end{array}
% \]
% so that $\denot{\maximum}(x, y; )= \max(x, y)$.  Notice that by replacing the first equation with $\denot{\maximum}(x, y; )     = \denot{\der_{11}}(x, y; x, y)$ then the resulting equational program computes $\min (x, y)$. The latter funcion can be defined by swapping the order of application of  the occurrences of $\cut_\n$ in $\maximum$.
% \end{figure*}

% \subsection{Example}
%  The proof of the translation lemma is quite technical, so we shall only discuss an example given by the following $\nbc$-coderivation $\der$:
% \[
% \small
%     \vlderivation{
%      \vliin{}{\red{\bullet}}{\tikzmarknode{fc}{\red{\denot \der}}}
%       {
%          \vltr{}{\tikzmarknode{f1c}{\orange{\denot {\der_0}}}}{\vlhy{}}{\vlin{}{\red{\bullet}}{\tikzmarknode{g1b}{\red{\denot\der}}}{\vlhy{}}}{\vlhy{}}
%       }{
%         \vlin{}{\blue{\bullet}}{\tikzmarknode{f2c}{\blue{\denot{\der_1}}}}{
%          \vltr{}{\green{\denot {\der_{10}}}}{\vlin{}{\red{\bullet}}{\tikzmarknode{g2b}{\red{\denot\der}}}
%          {\vlhy{}}}
%          {
%           \vlhy{}
%          }{\vlin{}{\blue{\bullet}}{\tikzmarknode{f2b}{\blue{\denot{\der_1}}}}
%          {\vlhy{}}}
%      }
%      }
%     }
% \begin{tikzpicture}[overlay, remember picture]
%     %  \draw [-, thick, green] ([shift=({0,0.3})]fb.north) 
%     %   .. controls +(30:4cm) and   +(-10:4cm) ..
%     %  ([shift=({0.3,0})]fc.east);
%     %  \draw [-, thick, blue] ([shift=({0,0.3})]f2b.east) 
%     %   .. controls +(30:2cm) and   +(0:1.5cm) ..
%     %  ([shift=({0.3,0})]f2c.east);
%     %  \draw [-, thick, red] ([shift=({0,0.3})]g1b.north) 
%     %   .. controls +(120:3cm) and   +(0:-3.5cm) ..
%     %  ([shift=({-0.3,0})]g1c.west);
%     % %   \draw [-, thick] ([shift=({0,0.3})]a.north) 
%     % %   .. controls +(70:1.5cm) and   +(90:2cm) ..
%     % %   (0.5, 0);
%     % %     \draw [->,>=latex,thick] (0.5, 0)
%     % %   .. controls +(-90:2cm) and   +(-70:0.5cm) ..
%     % %  ([shift=({0,-0.2})]b.south);
%     %%%%%%%
% % \node[anchor=base] (a) at (0,1) {\vphantom{h}}; % push the mark to the top of the line (ie including ascenders)
% % \node[anchor=base] (b) at (0,-0.8){\vphantom{g}}; % push the mark to the bottom of the line (ie including descenders)
% % \node[anchor=base] (c) at (0,-0.7) {\vphantom{h}}; % 
% % \node[anchor=base] (d) at (0,-1.2){\vphantom{g}}; % 
% % \draw [decoration={brace,amplitude=0.5em},decorate, thick,black]
% %  ([shift=({0.4,0})]a.east -| {pic cs:right}) -- node[right]{~~\blue{\textit{\small{Proof}}}} ([shift=({0.4,0})]b.north -| {pic cs:right});
% %  \draw [decoration={brace,amplitude=0.5em},decorate, thick,black]
% %  ([shift=({0.4,0})]c.east -| {pic cs:right}) -- node[right]{~~\green{\textit{\small{Theorem}}}} ([shift=({0.4,0})]d.north -| {pic cs:right});
% \end{tikzpicture}
% \]
% were we replace a node $v$ in position $i$ with the corresponding function  $\denot{\der_i}$, $\denot \der$ being $\denot{\der_\epsilon}$. We shall consider on the \emph{finite} representation of $\der$ using bullets (see~\ref{con}) obtained by~\Cref{conv:convention} and~\Cref{rem:regular-finite-equation}), and reason by induction on the number of nodes, associating with a node $v$ in position $i$ a set of functions defining the programs 




% With $\green{\denot{\der_{10}}}$ is associated a single equation with two oracles, $\red{\denot{\der}}$ and $\blue{\denot{\der_1}}$:
% \[
% \arraycolsep=2pt
% \begin{array}{rcl}
%     \green{\denot{\der_{10}}}(\vec x; \vec y) & =& h_{10}(\vec x; \vec y,  \lambda \vec u\subseteq \vec x. \lambda v. \red{\denot{\der}}(\vec u; \vec v),  \lambda \vec u\subseteq \vec x. \lambda v.  \blue{\denot{\der_1}}(\vec u; \vec v))  
% \end{array}
% \]
% The set of equations for $\blue{\denot{\der_{1}}}$ is obtained by extending the previous one with the following:
% \[
% \arraycolsep=2pt
% \begin{array}{rcl}
%     \blue{\denot{\der_{1}}}(\vec x; \vec y) & =& h_{10}(\vec x; \vec y,  \lambda \vec u\subseteq \vec x. \lambda v. \red{\denot{\der}}(\vec u; \vec v),  \lambda \vec u\subset \vec x. \lambda v.  \blue{\denot{\der_1}}(\vec u; \vec v))
%     \end{array}
% \]
% Notice that we have $\subset$ in the lambda abstraction for $\blue{\denot{\der_1}}$ because of the  progressiveness condition. 







% \section{Research Methods}

% \subsection{Part One}

% Lorem ipsum dolor sit amet, consectetur adipiscing elit. Morbi
% malesuada, quam in pulvinar varius, metus nunc fermentum urna, id
% sollicitudin purus odio sit amet enim. Aliquam ullamcorper eu ipsum
% vel mollis. Curabitur quis dictum nisl. Phasellus vel semper risus, et
% lacinia dolor. Integer ultricies commodo sem nec semper.

% \subsection{Part Two}

% Etiam commodo feugiat nisl pulvinar pellentesque. Etiam auctor sodales
% ligula, non varius nibh pulvinar semper. Suspendisse nec lectus non
% ipsum convallis congue hendrerit vitae sapien. Donec at laoreet
% eros. Vivamus non purus placerat, scelerisque diam eu, cursus
% ante. Etiam aliquam tortor auctor efficitur mattis.

% \section{Online Resources}

% Nam id fermentum dui. Suspendisse sagittis tortor a nulla mollis, in
% pulvinar ex pretium. Sed interdum orci quis metus euismod, et sagittis
% enim maximus. Vestibulum gravida massa ut felis suscipit
% congue. Quisque mattis elit a risus ultrices commodo venenatis eget
% dui. Etiam sagittis eleifend elementum.

% Nam interdum magna at lectus dignissim, ac dignissim lorem
% rhoncus. Maecenas eu arcu ac neque placerat aliquam. Nunc pulvinar
% massa et mattis lacinia.

\end{document}
\endinput
%%
%% End of file `sample-authordraft.tex'.
